% concatenated from main.tex
\documentclass{article}
\usepackage{graphicx} %
\usepackage{xparse}

\usepackage{latexsym,amssymb,amsmath,amsfonts,amsthm}
\numberwithin{equation}{section}
\usepackage{enotez,hyperref,float}
\usepackage{forest}
\usepackage{thmtools, thm-restate}
\usepackage{mathtools}

\usetikzlibrary{arrows.meta, positioning}
\usetikzlibrary{decorations.markings}

\setlength{\textwidth}{6.0in}
\setlength{\textheight}{8.5in}
\setlength{\topmargin}{-.25in}
\setlength{\oddsidemargin}{.0in}
\setlength{\evensidemargin}{.0in}

\usepackage{color}
\usepackage{epsfig}
\usepackage{enumerate}
\usepackage{comment}

\usepackage{svg}
\newenvironment{claimproof}[1]{\par\noindent{Proof:}\space#1}{\hfill $\blacksquare$}
\newtheorem{theorem}{Theorem}
\newtheorem{acknowledgement}[theorem]{Acknowledgement}
\newtheorem{algorithm}[theorem]{Algorithm}
\newtheorem{axiom}[theorem]{Axiom}
\newtheorem{case}[theorem]{Case}
\newtheorem{claim}[theorem]{Claim}
\newtheorem{conclusion}[theorem]{Conclusion}
\newtheorem{condition}[theorem]{Condition}
\newtheorem{conjecture}[theorem]{Conjecture}
\newtheorem{corollary}[theorem]{Corollary}
\newtheorem{criterion}[theorem]{Criterion}
\newtheorem{definition}{Definition}
\newtheorem{example}[theorem]{Example}
\newtheorem{exercise}[theorem]{Exercise}
\newtheorem{lemma}{Lemma}

\graphicspath{ {./assets/} }

\newtheorem{notation}[theorem]{Notational Convention}
\newtheorem{problem}[theorem]{Problem}
\newtheorem{proposition}{Proposition}
\newtheorem*{proposition*}{Proposition}
\newtheorem{remark}{Remark}
\newtheorem{solution}[theorem]{Solution}
\newcommand{\ve}{{\varepsilon}}
\newcommand{\R}{{\mathbb{R}}}
\newcommand{\EE}{{\mathbb{E}}}
\newcommand{\Prob}{{\mathbb{P}}}
\newcommand{\bphi}{{\boldsymbol{\phi}}}
\DeclareMathOperator*{\argmax}{arg\,max}
\newcommand{\Var}{\mathrm{Var}}
\newcommand{\tri}{{\triangle}}
\newcommand{\ceil}{\mathrm{ceil}}
\newcommand{\Cov}{\mathrm{Cov}}
\newcommand{\del}{{\partial}}

\newcommand{\gsimas}[2]{\overset{{#1}\to {#2}}{\gtrsim}}
\newcommand{\lsimas}[2]{\overset{{#1}\to {#2}}{\lesssim}}
\newcommand{\simas}[2]{\overset{{#1}\to {#2}}{\sim}}
\newcommand{\ggas}[2]{\overset{{#1}\to {#2}}{\gg}}
\newcommand{\llas}[2]{\overset{{#1}\to {#2}}{\ll}}
\newcommand{\ordas}[2]{\overset{{#1}\to {#2}}{\asymp}}



\def\beq{ \begin{equation} }
\def\eeq{ \end{equation} }
\def\mn{\medskip\noindent}
\def\ms{\medskip}
\def\bs{\bigskip}
\def\bn{\bigskip\noindent}
\def\ep{\epsilon}
\def\nass{\noalign{\smallskip}}

\def\nams{\noalign{\smallskip}}
\def\square{\vcenter{\vbox{\hrule height .4pt
  \hbox{\vrule width .4pt height 5pt \kern 5pt
        \vrule width .4pt} \hrule height .4pt}}}
\def\eopt{\hfill$\square$}
\def\RR{\mathbb{R}}
\def\ZZ{\mathbb{Z}}
\def\eqd{\,{\buildrel d \over =}\,}
\def\var{\hbox{var}\,}
\def\cov{\hbox{cov}\,}
\def\hbr{\hfill\break}
\def\sqz{\kern-0.2em}
\newcommand*\samethanks[1][\value{footnote}]{\footnotemark[#1]}
\def\clearp{\clearpage}

\newcommand{\tabeq}[3]{%
  \begin{aligned}[t]
    \makebox[#1][l]{$#2$} &= #3
  \end{aligned}%
}


\usepackage{graphicx}
\graphicspath{%
    {converted_graphics/}%
    {/}%
}

\newcommand\blfootnote[1]{%
  \begingroup
  \renewcommand\thefootnote{}\footnote{#1}%
  \addtocounter{footnote}{-1}%
  \endgroup
}

\renewcommand\qedsymbol{$\blacksquare$}

\newtheorem{assumptionstar}{Assumption}
\renewcommand{\theassumptionstar}{*}

\usepackage[
backend=biber,
style=numeric,
sorting=none
]{biblatex}
\addbibresource{bib.bib}

\usepackage{authblk}

\newcommand{\arm}[1]{\textcolor{blue}{[Armaan: #1]}}
\title{The Site Frequency Spectrum in an Exponentially-Growing Population with Selection}
\author[1,2,3]{Armaan Ahmed}
\author[4]{Jasmine Foo}
\author[5]{Einar Bjarki Gunnarsson}
\author[6]{Kevin Leder}

\affil[1]{Department of Applied Mathematics \& Statistics, Johns Hopkins University}
\affil[2]{Department of Mathematics, Johns Hopkins University}
\affil[3]{Department of Statistics, University of California, Berkeley}
\affil[4]{Department of Mathematics, University of Minnesota}
\affil[5]{Division of Applied Mathematics, Science Institute, University of Iceland}
\affil[6]{Department of Industrial and Systems Engineering, University of Minnesota}

\setcounter{secnumdepth}{4} %
\setcounter{tocdepth}{4}
\usepackage{nicematrix}
\newcommand{\red}[1]{\textcolor{red}{#1}}

\providecommand{\classification}[2]
{
  \small	
  \textbf{\textit{#1---}} #2
}



\begin{document}

\maketitle
\begin{abstract}
  We consider a supercritical two-type continuous-time linear birth-death process with mutation and selection, in which wild-type individuals give rise to mutant offspring with a larger net growth rate. In this setting, we investigate the site frequency spectrum (SFS) of driver mutations, describing the number of driver mutations present at any given frequency in the population.
  First, we derive exact moments for the SFS and establish asymptotic power laws at large times and frequencies. %
  Then, strong laws of large numbers for the driver SFS are proven by constructing suitable $L^2$-approximations.  
  These results apply both to the case when all driver clones have the same selective advantage and when the selective advantage is random. 
  Finally, we allow the frequency to vary with time to examine the number of ``intermediate'' and ``large'' driver clones, identifying
  a cutoff 
  frequency at which there are order 1 number of mutant clones. 
  Overall, our results provide quantitative insights into how selection shapes the site frequency spectrum both at small and large frequencies, which can in principle be leveraged to construct estimators of relevant evolutionary parameters, including the selective advantage of driver mutations.
\end{abstract}

\classification{Keywords}{Site Frequency Spectrum, Birth-Death Process, Cancer Evolution}

\classification{MSC2020}{Primary: 60J80. Secondary: 92D10, 92D25, 60F15}

\section{Introduction}

Human cells accumulate mutations throughout the lifetime of a person due to various factors such as DNA replication errors, chronic inflammation, and decreased DNA repair efficacy \cite{takeshima2019}. 
So-called `driver' mutations confer a fitness advantage to a cell, enabling clonal expansion and accumulation of additional mutations, leading to premalignant or cancerous cell populations \cite{bielas2006,hanahan2022}. In addition to these driver mutations, neutral and deleterious mutations accumulate in developing tumors, which can occur at enhanced rates compared to normal cells \cite{lawrence2014,mcfarland2013}. Together, these various mutations define the genetic tumor heterogeneity and contribute towards the total phenotypic tumor heterogeneity. Understanding the genetic tumor heterogeneity remains an important goal as it can provide clinically-relevant information regarding the mechanisms of tumor formation, tumor classification, prediction of therapeutic response, etc. \cite{predicting,salvadores2019,youn2013} 

\begin{figure}[H]
    \centering
    \def\svgwidth{0.8\linewidth}
    \input{assets/intro_mutations.pdf_tex}
    \caption{Example diagram of the two-type population with neutral mutations. Each color represents a distinct mutation, and arrows indicate how the mutation is transferred through the population.}\label{fig:example_mutation}
\end{figure}

A simple summary metric of the genetic heterogeneity of a tumor is the site frequency spectrum (SFS), denoted $S_j(t)$, counting the number of distinct mutations found in $j$ tumor cells (hereafter ``at frequency $j$'') at time $t$. 
The SFS can be used to estimate evolutionary parameters of a phylogeny (mutation, birth, death rates) and to provide evidence for either neutral or selective evolution \cite{tung,gunnarsson2021,tanaka2023,gunnarsson2025}.



\begin{table}[H]
    \centering
    \begin{NiceTabular}{|w{c}{3cm}|c|c|}
\hline
{\diagbox{\textbf{Initiated by}}{\textbf{Present in}}}& \raisebox{-1.5ex}{Type-0} & \raisebox{-1.5ex}{Type-1} \\[3ex]
\hline\hline \\[-2ex]
Type-0 & $S_j^{p,0}(t)$ & $S_j^h(t),\red{S_j^d(t)}$ \\[1ex]
Type-1 & $\emptyset$ & $S_j^{p,1}(t)$ \\[1ex] 
\hline 
\end{NiceTabular}
\caption{Mutations found at time $t$ in Type-0 and Type-1 cells are stratified based on whether they are neutral or selective (in \red{red}), the type of cell in which they originated. Details of each of these contributions to the SFS are present in the main text.}\label{table:SFS_components}
\end{table}

In this work, we characterize the behavior of the SFS in the setting of a two-type continuous-time linear birth-death process. The initial population consists of Type-0 (wild-type) cells which can give rise to Type-1 cells through the acquisition of a selectively advantageous driver mutation.  In the setting where each cell can acquire at most one driver mutation, the SFS can be decomposed as indicated by Table \ref{table:SFS_components}. $S_j^d(t)$ counts the number of driver clonal lineages of size $j$ at time $t$, $S_j^h(t)$ counts the number of ``hitchhiker'' mutations (the neutral mutations present in type-1 cells that first arose in some wild-type cell) at frequency $j$ at time $t$, $S_j^{p,0}(t)$ counts the number of type-0 ``passenger'' mutations (the neutral mutations present in type-0 cells) at frequency $j$ at time $t$, and $S_j^{p,1}(t)$ counts the number of type-1 ``passenger'' mutations (the neutral mutations that first arise in type-1 cells) at frequency $j$ at time $t$. An example diagram with all these mutation types is provided in Figure \ref{fig:example_mutation}. Our ultimate goal is to develop a comprehensive understanding of all contributions to the SFS in a population with selection. In this work, we ignore neutral mutation accumulation and instead focus on characterizing the behavior of the driver SFS ($S^d_j(t)$). Accordingly in order to simplify the notation, we will drop the superscript and write $S_j(t)$ instead of $S_j^d(t)$ for the remainder of the manuscript.

The setting of purely neutral evolution, in which the genetic heterogeneity arises solely from the accumulation of passenger mutations, has been well-studied. In this scenario, every cell in the tumor-initiating population already contains all driver mutations.
In prior works, Gunnarsson et al.~examined continuous-time linear birth-death processes with neutral mutation accumulation. They found that, conditional on non-extinction, at large fixed times $t$ and large frequencies $j$, the number of mutations at frequency $j$ ($S_j^p(t)$) follows a $1/j^2$ power law, both in expectation and almost surely \cite{gunnarsson2021,gunnarsson2025}. Cheek and Antal examined the mean neutral site frequency spectrum in the large-time and large-population small-mutation limit, by categorizing the subpopulations of cells either carrying or not carrying a particular mutation into two types, also finding a $1/j^2$ law \cite{cheek2018}. A special feature of \cite{cheek2018} is that they did not make the infinite-sites assumption. Other authors have obtained similar results, including cumulative versions capturing the number of mutations found in a \textit{proportion larger than $f \in (0,1)$} of cells, in various birth-death, semi-deterministic, or fully deterministic exponential growth models \cite{durrett2013,bozic2016,williams2016,OHTSUKI201743}. In constant-size Moran models with neutral mutations, the SFS follows a $1/j$ law in expectation \cite{durrett_probability_2008}. In both exponentially-growing and fixed-size population models, characterizing $S_j^p(t)$ allows one to estimate the neutral mutation rate, obtain the probability of extinction of the tumor (which with the net growth rate enables the decoupling of birth from death rates), and test treatment response \cite{gunnarsson2021,gunnarsson2025,stein2025,durrett_probability_2008}.

In the setting where the tumor cells can acquire further driver mutations, there are many studies examining the bulk dynamics of the total driver-mutant population (arrival times of mutants, asymptotic total mutant population size, etc.) \cite{iwasa2006,durrett2010,nicholson2023}. However, the dynamics of \textit{each distinct} clonal lineage are what matter in  quantifying $S_j(t)$. In an early study, Lea and Coulson examined a semi-deterministic two-type model of wild-type and mutant bacteria growth \cite{lea1949,zheng1999}. The wild-type population grew deterministically over time, seeding as a Poisson point process mutant clones growing according to a linear birth-death process. They calculated the joint distribution of the total number of mutant cells and the number of mutations that occurred by time $t$, which reveals the average size of all clonal lineages produced by time $t$. Under the infinite sites assumption, Cheek and Antal also provide results concerning the site frequency spectrum and the largest clone under selective evolution in the large time-small mutation or large detection size-small mutation regime \cite{cheek2018}. In their scaling regime, only $\mathcal{O}(1)$ clones contribute to the mutant population whereas in our large time or detection size regime infinitely many clones contribute. Consequently, they obtain Poisson limit theorems whereas we obtain law of large numbers limit theorems for the SFS. This is discussed in Section \ref{sec:lln_SFS} in more detail. Tung and Durrett recently examined a cumulative version of the SFS, the \textit{relative} SFS, describing the number of clones that constitute a proportion larger than $f\in (0,1)$ of the total mutant population \cite{tung}. They find that in the large-time limit, the expected number of clones with frequency $>f$ follows a $f^{-\lambda_0/\lambda_1}$ power-law where $\lambda_0$ (or $\lambda_1$) is the net growth rate of the wild-type (or mutant) population. We will also obtain results for the relative SFS and compare them to the results of Tung and Durrett in Section \ref{sec:convergence_relative_SFS}.

There are comparatively fewer studies that examine frequencies of hitchhiker mutants ($S_j^h(t)$). Braverman et al.~using a coalescent model examines the effect of hitchhikers on the genetic heterogeneity in a bacterial population \cite{braverman1995}. More recently, Bonnet and Leman examine $S_j^h$ in a two-type continuous time linear birth-death process with rare mutants under rescue dynamics (in which the wild-type population is subcritical and the mutant population is supercritical) \cite{bonnet2024}. This is relevant when, for example, modeling susceptible-resistance cancer cell dynamics under treatment. In this setting, they find $\mathbb{E}[S_j^h(t_N)]\lesssim N^{1-\alpha}$ as $N\to\infty$ where $N$ denotes the initial sensitive population, $t_N\propto \log N$ is the characteristic timescale of extinction of the sensitive type, and $0<\alpha\leq 1$ is a parameter related to the rarity of mutations.

\subsection{Outline}

In this paper we study the behavior of the driver SFS ($S_j(t)$) in a two-type continuous-time linear birth-death process, under the assumption that cells can accumulate at most one driver mutation. In Sections \ref{sec:exact_moments_SFS} and \ref{sec:mean_j_scaling}, we obtain exact and asymptotic (in time and frequency) expressions for moments of the SFS,  when the selective advantage of each driver clone is deterministic or random. In Section \ref{sec:lln_SFS}, we extend analyses from Gunnarsson et al. and Harris (\cite{gunnarsson2025} and \cite{harris1963}, respectively) to obtain almost sure convergence results for $S_j(t)$ in the large time and detection size limits. We also provide asymptotic limit theorems when letting the frequency $j(t)$ vary with time in Section \ref{sec:SFSVary}.  In Section \ref{sec:tail_SFS}, we investigate the shape of the tail of the SFS at fixed detection sizes. We then go on to study the frequencies of the largest clones in Section \ref{sec:large_frequencies}. We provide limit theorems for the relative SFS in Section \ref{sec:convergence_relative_SFS}. Finally, the size of the largest mutant clone is examined in Section \ref{sec:smallmut_largeclone}. We further compare some of our results with results obtained from Durrett's semideterministic model and show agreement in the small mutation limit.

\section{Models and notation}

\subsection{Two-Type Population Model}
\label{sec:model_notation}

\begin{figure}[H]
    \centering 

    \begin{minipage}{0.45\textwidth}
        \centering
        \renewcommand{\arraystretch}{1.2}
        \begin{tabular}{cl}
            \hline\hline
            \multicolumn{2}{c}{\textbf{Model Parameters}} \\
            \hline
            \textbf{Symbol} & \textbf{Description}      \\
            \hline
            $a$           & Type 0 per-capita birth rate    \\
            $d$           & Type 0 per-capita death rate     \\
            $v$             & Probability of mutation per birth     \\
            $\mu$           & Birth-independent per-capita mutation rate   \\
            $b$             & Type 1 per-capita increase in birth rate   \\
            \hline\hline
        \end{tabular}
    \end{minipage}
    \hfill 
    \begin{minipage}{0.45\textwidth}
        \begin{align*}
            0&\underset{d}{\to}\emptyset\\
            0&\underset{a}{\to}0,\text{Bernoulli}(v)\\
            0&\underset{\mu}{\to}1\\
            1&\underset{d}{\to}\emptyset\\
            1&\underset{a+b}{\to}1,1
        \end{align*}
    \end{minipage}

    \begin{minipage}{\textwidth}
      \centering
      \begin{forest}
  red node/.style={content={#1}, draw=red, fill=red!20},
  for tree={
    grow=east,
    circle,
    draw=blue,
    fill=blue!20,
    thick,
    l sep=5mm,
    s sep=4mm,
    minimum size=0.5cm,
    edge={-{Stealth[]}, shorten >=2pt},
  }
  [
    [
      [, content={\Large$\emptyset$}, draw=none,fill=none]
    ]
    [
      [
        [
          [[, content={\Large$\emptyset$}, draw=none,fill=none]]
          []
        ]
        [
          []
          [[, red node={\small B}, for children={red node={2}}
            [\small B]
            [\small B]]]
        ]
      ]
      [, red node={\small A}, for children={red node={2}}
        [\small A]
        [\small A]]
    ]
  ]
\end{forest}
    \end{minipage}
    \caption{Parameter descriptions (top left) and model scheme (top right). An example graphical representation of the model is shown in the bottom panel. Blue circles represent type-0 cells and red circles represent type-1 cells. $\emptyset$ denote death events and the differently labeled type-1 cells indicate which clone they belong to. The leaves (cells with no out-going arrows) are the cells currently alive.}
    \label{fig:model}
\end{figure}


Consider a continuous-time Markov branching process with two cell types, denoted type-0 and type-1 \cite{athreya1972}. Each type-0 cell independently divides at exponential rate $a$ and dies at rate $d$,  and the total type-0 population count at time $t$ is denoted $Z_0(t)$. Each type-1 cell  divides at rate $a+b$ and dies at rate $d$, where $b>0$ reflects a selective advantage. The total type-1 population count at time $t$ is denoted $Z_1(t)$. Unless otherwise specified, we begin with the initial condition $(Z_0(0),Z_1(0))=(1,0)$.


Type-0 cells produce type-1 cells through two mechanisms of driver mutation acquisition, (i) {\it replication-dependent:} at each type-0 division exactly one of the two offspring is type-1 with probability $v$, otherwise both are type-0 \footnote{Here, we are making a simplification and ignoring the possibility that both offspring are type-1 cells. This scenario occurs with probability $\mathcal{O}_{v\to 0^+}(v^2)$.}, and (ii) {\it replication-independent:} each type-0 cell mutates to type-1 at rate $\mu$, independently of division \cite{abascal2021}. Thus, the effective birth rate of type-0 cells is $a(1-v)$ and the net growth rate is $\lambda_0 \equiv a(1-v)-d-\mu$. The net growth rate of type-1 cells is $\lambda_1 \equiv a+b-d$. Throughout, we assume supercriticality of the type-$0$ population, namely $\lambda_0>0$. See Figure \ref{fig:model} for a schematic of this model. Observe that clone A is founded by a mutant produced in a replication-dependent fashion and clone B is founded by a mutant produced in a replication-independent fashion. 

In some sections we will assume that driver mutations are replication-dependent ($\mu=0$) and take the small mutation limit $v\to 0^+$. In this case, the notation $\lambda_{0,v}$ will be used to explicitly denote the dependence of $\lambda_0$ on $v$. Denote the probabilities of extinction of a type-0 and of a type-1 clone, each starting from one cell of the respective type, as $p_0=\frac{d+\mu}{a(1-v)}$ and $p_1=\frac{d}{a+b}$, respectively.  These follow from the classical extinction probability formula for supercritical birth-death processes.  For $i\in \{0,1\}$, define $q_i=1-p_i$ as the probability that a clone started by a type-$i$ cell is established, i.e.~survives forever. Where necessary, we will condition on non-extinction, and thus define $\Omega_i^\infty:=\{\exists N\text{ s.t. }t\geq N\text{ implies }Z_i(t)>0\}$ for $i\in \{0,1\}$ as the event that the type-$i$ population does not go extinct. Remark that $ \Omega_0^\infty\subseteq \Omega_1^\infty$ (up to null sets).

Let $S_j(t)$ denote the number of type-1 driver clones of size $j$ at time $t$. $S_j(t)$ is the driver SFS, labeled $S_j^d(t)$ in Table \ref{table:SFS_components}; we usually call it `the SFS' in this paper. Here, we make the infinite sites assumption, so for example the mutations associated with clones A and B in Figure \ref{fig:model} are distinct\footnote{Before defining $S_j(t)$ more carefully, we clarify the relation between mutations at a given frequency versus clones at a given size. Firstly, a type-1 clone is the tree of ancestors rooted at a type-1 founder cell (i.e., a type-1 cell whose parent is a type-0 cell). Associated to each clone is a selective mutation that its founder cell acquired. It may be that different clones have acquired the same selective mutation. If clones $A$ and $B$ in Figure \ref{fig:model} both acquired the same mutation, then at the time of detection clones $A$ and $B$ would both have size $2$ whereas the frequency of the mutation would be $2+2=4$. We avoid this by making the infinite sites assumption, so the frequency of the mutation associated with clone $i$ at time $t$ is the same as the number of cells in clone $i$ at time $t$.}. Define $T_i$ as the initiation time of the $i$th driver clone produced, where $T_i=\infty$ if the $i$th clone is never produced, and define $M_t=\#\{i\in\mathbb{N}:T_i\leq t\}$ as the number of mutant lineages ever produced by time $t$. 
We let $Z_1^{(i)}(s)$ denote the size of the $i$th driver clone $s$ time units after it was produced, and note that 
\[
S_j(t)=\sum_{i:T_i\leq t} \mathbf{1}_{\{Z_{1}^{(i)}(t-T_i)=j\}}.
\]
As our branching process is reducible conditional on the trajectory of the type-0 population, $s\to Z_1^{(i)}(s)$ are i.i.d linear birth-death processes with birth rate $a+b$ and death rate $d$. If $s<0$, $Z_1^{(i)}(s)\equiv 0$.
We will also examine the {\it relative driver SFS}, the number of driver clones above a certain frequency: $R_f(t)=\sum_{j:j>Z_1(t)f}S_j(t)$ for $0<f<1$.

We will investigate features of the driver SFS at a fixed time $t$ and at the random time when the population reaches a fixed detection size $n$, $\tau_n=\inf \{t:Z_0(t)+Z_1(t)=n\}$. On $\Omega_1^\infty$, $\tau_n < \infty$ almost surely. Fixed time results are relevant when considering population growth in a controlled setting. Fixed size results are more relevant in the clinical setting, as a tumor may be readily detected upon reaching a particular size. 

Certain limit theorems are known for $Z_0(t),Z_1(t),\{Z_1^{(i)}(t)\}_{i\in\mathbb{N}}$. For supercritical birth-death processes, it is well-known that $e^{-\lambda_0 t}Z_0(t)$ converges almost surely and in $L^2$ as $t \to \infty$.
The limiting random variable $Y \equiv \lim_{t \to \infty} e^{-\lambda_0 t}Z_0(t)$ is 0 with probability $p_0$, and on the event of survival it is exponential with rate $q_0$ (for example see Theorem 1 in \cite{durrettbranching2015}). Similarly, $Y_1^{(i)}\equiv\lim_{t\to\infty}e^{-\lambda_1 t}Z_1^{(i)}(t)$ exists almost surely and in $L^2$. The $Y_1^{(i)}$ are i.i.d. random variables which are 0 with probability $p_1$ and exponential with rate $q_1$ with probability $q_1$. For the total type-1 population, $W\equiv\lim_{t\to\infty}e^{-\lambda_1 t}Z_1(t)=\sum_{i\in\mathbb{N}}e^{-\lambda_1 T_i}Y_1^{(i)}$ exists almost surely on the event of nonextinction (see Appendix Section \ref{app:W_representation}). 



{\bf \em Random Fitness Advance Model.}
We will also consider a variant of the above model in which we allow each type-1 clone to have random fitness advances. In particular, consider a random perturbation to the birth rate of the $i$-th type-1 clone, $B_i$, where $B_j\sim \mathbb{Q}$ are i.i.d.~for all $j\in \mathbb{N}$ for some distribution $\mathbb{Q}$, and the $B_j$'s are independent of the type-0 process $Z_0(t)$.  Conditional on $B_i$, the birth rate of the $i$-th type-1 clone is $a+B_i$ while the death rate remains $d$. %
We refer to this modified model as the random fitness advance model, and we denote its SFS by $S_j^B$.
We note that a semi-deterministic version of this model was also studied by the second and fourth authors in various previous works \cite{durrettfoo2010, durrettgenetics2010, foo2014escape}. The following bounded support assumption on $\mathbb{Q}$ is sometimes used in our results for this random fitness model:

\begin{assumptionstar}
    $\mathbb{Q}$ has density $h$ with support $[0,b_{\max}]$ which is left-continuous at $b_{\max}$, $h(b_{\max})>0$, and $h$ is bounded.\label{assumption:random}
\end{assumptionstar}

A summary of all the important notation is provided in Table \ref{table:notation}.
\begin{table}[H]
\renewcommand{\arraystretch}{1.5}
  \centering  
  \begin{tabular}{|c p{6cm} p{6cm}|}
\hline
\textbf{Symbol} & \textbf{Description} & \textbf{Mathematical Description}\\
\hline
$Z_0(t)/Z_1(t)$ & Count of total type-0/type-1 population at time $t$ & \rule[0.5ex]{1cm}{0.4pt}\\
$\Omega_i^\infty$ & Extinction event for the type-$i$ population, where $i\in \{0,1\}$, population & $\tabeq{0.9cm}{\Omega_i^\infty}{\{\omega:\exists T\geq 0,\\&t\geq T\Longrightarrow Z_i(t)=0\}}$\\
$T_i$ & Timing of $i$th mutation event & \rule[0.5ex]{1cm}{0.4pt}\\
$M_t$ & Counting process of number of mutation events in $[0,t]$ & $\tabeq{0.9cm}{M_t}{\#\{i\in \mathbb{N}:T_i\leq t\}}$\\
$Z_1^{(i)}(s)$ & Count of cells in the $i$th type-1 clone, $s$ time units after it was produced & $\tabeq{0.9cm}{Z_1(t)}{\sum_{T_i\leq t}Z_1^{(i)}(t-T_i)\\
       &= \sum_{i\leq M_t}Z_1^{(i)}(t-T_i)}$\\
$Y/Y_1^{(i)}/W$ & Scaled limit of the type-0 population/ $i$th type-1 clone / type-1 population &  $\tabeq{0.9cm}{Y}{\lim_{t\to\infty}e^{-\lambda_0 t}Z_0(t)}$\newline$\tabeq{0.9cm}{Y_1^{(i)}}{\lim_{t\to\infty}e^{-\lambda_1 t}Z_1^{(i)}(t)}$\newline$\tabeq{0.9cm}{W}{\lim_{t\to\infty} e^{-\lambda_1 t}Z_1(t)\\&=\sum_{i \in \mathbb{N}}Y_1^{(i)}e^{-\lambda_1 T_i}}$\\
$\tau_n$ & First time the total population reaches size $n$ & $\tabeq{0.9cm}{\tau_n}{\inf\{t\geq 0:Z_0(t)+Z_1(t)=n\}}$\\
$S_j(t)/S_{\geq j}(t)$ & Number of type-$1$ clones composed of $j$/$\geq j$ cells at time $t$ & $\tabeq{0.9cm}{S_j(t)}{\sum_{T_i\leq t}\mathbf{1}_{Z_1^{(i)}(t-T_i)=j}}$\newline$\tabeq{0.9cm}{S_{\geq j}(t)}{\sum_{k\geq j}S_k(t)}$\\
$S_j^B(t)$ & Number of type-$1$ clones composed of $j$ cells at time $t$, where each clone has a random fitness advancement with the law of $B$ & \rule[0.5ex]{1cm}{0.4pt}\\
$R_f(t)$ & Number of clones making up a proportion larger than $f\in (0,1)$ of the type-$1$ population at time $t$&$\tabeq{0.9cm}{R_f(t)}{\sum_{j>fZ_1(t)}S_j(t)}$\\
\hline
\end{tabular}
\caption{Notation used throughout the paper}\label{table:notation}
\end{table}

{\bf \em Asymptotic notation.} Letting $\ell\in \mathbb{R}\cup \{\pm \infty\}$ and $g,h:\mathbb{R}\to \mathbb{R}_{\geq 0}$, we use the following asymptotic notation throughout the paper:

\begin{align*}
    g\simas{t}{\ell}h&\Longleftrightarrow \lim_{t\to\ell}\frac{g(t)}{h(t)}=1\\
    g\lsimas{t}{\ell} h&\Longleftrightarrow \limsup_{t\to\ell}\frac{g(t)}{h(t)}<\infty\\
    g\gsimas{t}{\ell} h&\Longleftrightarrow h\lsimas{t}{\ell} g\\
    g\llas{t}{\ell} h&\Longleftrightarrow \lim_{t\to \ell}\frac{g(t)}{h(t)}=0\\
    g\ggas{t}{\ell} h&\Longleftrightarrow h\llas{t}{\ell} g\\
    g\ordas{t}{\ell} h&\Longleftrightarrow g\lsimas{t}{\ell} h\text{ and }g\gsimas{t}{\ell} h
\end{align*}

Additionally, the Landau-O notation is used to specify in-line error terms. That is, letting $g:\mathbb{R}\to \mathbb{R}$ and $h:\mathbb{R}\to \mathbb{R}_{\geq 0}$ be real-functions, we say:
\begin{align*}
    g=\mathcal{O}_{t\to \ell}(h)&\Longleftrightarrow |g|\lsimas{t}{\ell}h\\
    g=o_{t\to \ell}(h)&\Longleftrightarrow |g|\llas{t}{\ell}h
\end{align*}

Usually, we examine large-time asymptotics ($\ell = \infty$), in which case we drop $t\to\infty$ from the notation as long as the limit is clear.




\section{Results}



\subsection{Simulations reveal characteristics of the SFS under selection}
To motivate our theoretical studies, we first present numerical simulations of the SFS behavior when mutation acquisition is replication-dependent and the driver fitness advantage $b$ is deterministic.
Unless otherwise specified, the set of parameter values in Table \ref{table:default_params} is used throughout all simulations. 



\begin{table}[H]
  \centering
  \begin{tabular}{|c|p{2.5cm}|}
\hline
\textbf{Parameter} & \textbf{Default Value} \\
\hline
$a$ & 0.1\\
$d$ & 0.001\\
$v$ & 0.01\\
$\mu$ & 0\\
$b$ & 0.04\\
\hline
\end{tabular}
\caption{Default simulation parameter values}\label{table:default_params}
\end{table}







Empirically, our simulations demonstrate how selection modifies the shape of the SFS 
observed at a fixed time $t$ (Figure \ref{fig:SFS_characteristics}).  In particular, Figure \ref{fig:SFS_characteristics} demonstrates that as the selective advantage $b$ increases, the size of the largest clone, $\max\{j:S_j(t)\neq 0\}$, grows on average, the frequency of clones at the small $j$-end of the spectrum decreases, and the slope of the SFS at intermediate $j$ increases.


\begin{figure}[H]
	\centering

  \begin{minipage}{0.5\textwidth}
    \centering
    \includegraphics[width=\textwidth]{assets/Figure2.pdf}
  \end{minipage}
  \hfill
  \begin{minipage}{0.45\textwidth}
    \centering
      \begin{tabular}{|c|p{2cm}|}
\hline
\textbf{Features of SFS} & \textbf{Effect from Increased Selection $b$} \\
\hline
Tail of the SFS & Extends\\
\hline
Frequency at small $j$ & Decreases\\
\hline
Slope of SFS & Increases\\ 
\hline
\end{tabular}
  \end{minipage}
	\caption{Impact of selection on SFS at fixed time. (left) Simulations using $t=60$ are run, and the mean for each indicated selection value is plotted. (right) Table of impacts of selection on the shape features of the SFS.}\label{fig:SFS_characteristics}
\end{figure}


Intuitively, the small-frequency end of the SFS is occupied by mutants produced close to the detection time $t$, for example during the time interval $[t-\Delta t,t)$. As the selection strength $b$ increases, the number of mutants produced during $[t-\Delta t,t)$ remains the same, but the number of associated clones that remain at a small size decreases. Thus, the frequency of mutants at small $j$ decreases as $b$ increases. Conversely, the mid-to-large end of the SFS is occupied by more mutant lineages as the selection strength increases, thus modifying the slope. To better understand these effects of selection, we will establish limit theorems studying the shape of the SFS in the following sections. 

\subsection{Shape of the driver site frequency spectrum under selection}\label{sec:time_freq_scaling}


\subsubsection{Exact moments of the driver SFS at fixed time}\label{sec:exact_moments_SFS}

We first consider moments of the driver SFS at a fixed time, both under a deterministic selective advantage and in the random fitness advance setting.

\begin{restatable}{proposition}{SFSmean}\label{prop:SFS_mean}
  Consider the two-type process $(Z_0, Z_1)$ starting from initial condition $(1, 0)$. Denote $P_j(s)=\mathbb{P}(Z_1^{(1)}(s)=j|Z_1^{(1)}(0)=1)$. The exact representations of the mean and second moment of $S_j(t)$ are (for fixed time $t$ and fixed $j$):
  \begin{align}
    \mathbb{E}[S_j(t)]&=\int_0^t (av+\mu)\mathbb{E}[Z_0(s)]P_j(t-s)\,ds, \text{ and}\label{eq:mean_SFS}\\
    \mathbb{E}[{S_j(t)\choose 2}]&=\int_0^t \int_s^t ((\mu+av)^2 \mathbb{E}[Z_0(r)Z_0(s)]-\mu(\mu+av)\mathbb{E}[Z_0(r)])P_j(t-s)P_j(t-r)\,dr\,ds.\label{eq:choose_SFS}
  \end{align}
  Note $\mathbb{E}[Z_0(s)]=\exp(\lambda_0 s)$ from our initial condition. Denote $P_j(s|b)$ as the transition kernel of a single type-1 clone with fitness advantage $b$. Under the random fitness advance model:
  \begin{align}
    \mathbb{E}[S_j^B(t)]&=\int_0^t (av+\mu)\mathbb{E}[Z_0(s)]\mathbb{E}[P_j(t-s|B)]\,ds,\text{ and}\label{eq:mean_SFS_random_fitness}\\
    \mathbb{E}[{S_j^B(t)\choose 2}]&=\int_0^t \int_s^t ((\mu+av)^2 \mathbb{E}[Z_0(r)Z_0(s)]-\mu(\mu+(a+b)v)\mathbb{E}[Z_0(r)])\mathbb{E}[P_j(t-s|B)]\mathbb{E}[P_j(t-r|B)]\,dr\,ds.\label{eq:choose_SFS_random_fitness}
  \end{align}
\end{restatable}
\begin{proof}
  See Appendix Section \ref{app:SFS_mean}.
\end{proof}

The expression for the expected value of the SFS is quite intuitive: contributions to the SFS at time $t$ can be discretized into small bins $[s,s+dt]$. A mutant clone is generated at rate $(av+\mu) Z_0(s)$ at time $s$, and if it forms and grows to size $j$ in the remaining $t-s$ time units, then it contributes towards the SFS, $S_j(t)$. The exact mean (Eq. \ref{eq:mean_SFS}) is visualized in the left panel of Figure \ref{fig:driver_muts}.  For the second moment, we consider the mean of a more combinatorially friendly random variable $S_j(t)\choose 2$, and consider contributing \textit{pairs} of mutants that form on disjoint intervals $[s_1,s_1+dt],[s_2,s_2+dt]$ and that both reach size $j$ at time $t$.

{\bf \em Selection decreases the small-frequency end of the mean driver SFS.} We note that $P_j(t|b)$ is explicitly calculable (see Eq. \ref{eq:P_j}). Using this and assuming $a+b>d$, we can show that $\partial_b P_1(t|b)<0$ for all $t>0$. This implies that the small end of the mean SFS is decreasing in $b$ (using Eq. \ref{eq:mean_SFS}). An analog holds in the random fitness model: if $B_1$ stochastically dominates $B_2$ (for all $x\in\mathbb{R}_{\geq 0}$, $\mathbb{P}(B_1\leq x)\leq \mathbb{P}(B_2\leq x)$), then for all $t>0$, $\mathbb{E}[P_1(t|B_2)]\geq \mathbb{E}[P_1(t|B_1)]$ implying $\mathbb{E}[S_1^{B_2}(t)]\geq \mathbb{E}[S_1^{B_1}(t)]$ (using Eq. \ref{eq:mean_SFS_random_fitness}).

This effect is robust to conditioning on the population being non-zero at the time of detection ($Z_0(t)+Z_1(t)>0$). If $S_j(t)>0$, then automatically $Z_0(t)+Z_1(t)>0$, so
\begin{align}
  \mathbb{E}[S_j(t)|Z_0(t)+Z_1(t)>0]=\frac{\mathbb{E}[S_j(t)]}{1-\mathbb{P}(Z_0(t)=0,Z_1(t)=0)}.\label{eq:SFS_conditional}
\end{align}
Since the extinction probability $\mathbb{P}(Z_0(t)=0,Z_1(t)=0)$ by time $t>0$ is decreasing in $b$, $\mathbb{E}[S_1(t)|Z_0(t)+Z_1(t)>0]$ is still decreasing in $b$. This suggests that we can use the small-frequency end of the driver SFS to estimate the fitness increase $b$, which is discussed further in the Discussion.







\subsubsection{Asymptotic behavior of the mean driver SFS}\label{sec:mean_j_scaling}

\begin{figure}[H]
	\centering
	\includegraphics[width=\textwidth]{assets/time_result.pdf}
	\caption{(left) $5\cdot 10^6$ simulations (deterministic mutant fitness model) are performed to obtain empirical means (lines), and compared with exact means (x-markers). (right) $10^3$ simulations are performed to obtain the rescaled mean driver SFS at time $t=150$. The `large $j$ equivalent' line refers to the $j^{-\lambda_0/\lambda_1-1}$ term in Eq. \ref{eq:large_j_equivalent} in Proposition \ref{prop:SFS_mean_asymptotics}.}\label{fig:driver_muts}
\end{figure}

We can use Proposition \ref{prop:SFS_mean} to obtain large-time ($t\to\infty$) and large-frequency ($j \to \infty$) scaling behavior of the site frequency spectrum: %
\begin{restatable}{proposition}{SFSmeanasymptotics}
  \begin{align}
    \lim_{t\to\infty}e^{-\lambda_0 t}\mathbb{E}[S_j(t)]&=(av+\mu)\int_0^\infty e^{-\lambda_0 s}P_j(s)\,ds\nonumber\\
    &=(av+\mu)\frac{q_1}{\lambda_1}\frac{\Gamma(j)\Gamma(\lambda_0/\lambda_1+1)}{\Gamma(j+\lambda_0/\lambda_1+1)}\sum_{n=0}^\infty \frac{(\lambda_0/\lambda_1)_n(j)_n}{(j+1+\lambda_0/\lambda_1)_n n!}p_1^n\nonumber\\
    &=(av+\mu)\frac{\Gamma(\lambda_0/\lambda_1+1)}{\lambda_1 j^{\lambda_0/\lambda_1+1}}(1-p_1)^{1-\lambda_0/\lambda_1}+\mathcal{O}_{j\to\infty}\left(\frac{1}{j^{\lambda_0/\lambda_1+2}}\right).\label{eq:large_j_equivalent}
  \end{align}
  $\Gamma$ denotes the gamma function and $(x)_n:=\Gamma(x+n)/\Gamma(x)$. We also have the following rate of convergence in time for fixed $j$:
  \[\left(\lim_{t\to\infty}e^{-\lambda_0 t}\mathbb{E}[S_j(t)]\right)-e^{-\lambda_0s}\mathbb{E}[S_j(s)]\overset{s\to\infty}{\sim}(av+\mu)\frac{q_1}{\lambda_0/\lambda_1+1}e^{-(\lambda_0+\lambda_1)s},\]
  and under the random fitness advance model with Assumption \ref{assumption:random}:
  \[\lim_{t\to\infty}e^{-\lambda_0t}\mathbb{E}[S_j^B(t)]\overset{j\to\infty}{\sim} \frac{C}{j^{\frac{\lambda_0}{\lambda_1(b_{\max})}+1}\log j}\]
  where $C$ is a function of the parameters of the model which can be explicitly determined.

  \label{prop:SFS_mean_asymptotics}
\end{restatable}
\begin{proof}
  See Appendix Section \ref{app:SFS_mean_asymptotics}.
\end{proof}
Eq.~\ref{eq:large_j_equivalent} indicates that mutations at large frequencies at large times follow a power law $j^{-\lambda_0/\lambda_1-1}$, as illustrated by the right panel in Figure \ref{fig:driver_muts}.
When $\lambda_0=\lambda_1$, we obtain a $j^{-2}$ power law, characteristic of a single-type model with neutral mutations.
When $\lambda_0 < \lambda_1$, there are more mutations at large frequencies than in the neutral model, and the number of mutations increases as the fitness advantage $b$ increases, which is consistent with the increase in slope we observed in Figure \ref{fig:SFS_characteristics}.
Under the random fitness advance model, mutations at large frequencies come from the clones with the largest fitness advance $b_{\max}$, which introduces a log-correction to the power law due to the rarity of clones with near-maximal fitness. 

\subsubsection{Large time and size asymptotics for the driver SFS}\label{sec:lln_SFS}

We now extend the analysis from Gunnarsson et al.~\cite{gunnarsson2025} to obtain a.s.~results for a scaled version of the SFS with selection, both at fixed detection time and size. 
The fixed-time limit (Eq.~\ref{eq:driver_lln_time}) is influenced by the limiting random variable $Y = \lim_{t \to \infty} e^{-\lambda_0t} Z_0(t)$ for the wild-type process that generates the mutant clones.
The fixed-size limit (Eq.~\ref{eq:driver_lln_size}) is further influenced by the limiting random variable $W = \lim_{t \to \infty} e^{-\lambda_1t} Z_1(t)$ for the type-1 process, which controls the time it takes the population to reach size $n$.
Both limits have a deterministic factor that comes from the first-moment limit in Eq.~\ref{eq:large_j_equivalent}.

\begin{restatable}{theorem}{ConvergenceResults} \label{prop:convergence_results}Conditional on $\Omega_0^\infty$ and for all $j\geq 1$:
  \begin{align}
    &e^{-\lambda_0 t}S_j(t)\overset{\substack{L^2\\a.s.}}{\to}(av+\mu)Y\int_0^\infty e^{-\lambda_0 s}P_j(s)\,ds=Y\lim_{s\to\infty}e^{-\lambda_0 s}\mathbb{E}[S_j(s)],\label{eq:driver_lln_time}\\
    &n^{-\lambda_0/\lambda_1}S_j(\tau_n)\overset{\substack{a.s.}}{\to}(av+\mu)YW^{-\lambda_0/\lambda_1}\int_0^\infty e^{-\lambda_0 s}P_j(s)\,ds=YW^{-\lambda_0/\lambda_1}\lim_{s\to \infty}e^{-\lambda_0 s}\mathbb{E}[S_j(s)].\label{eq:driver_lln_size}
  \end{align}
  Under the random fitness advance model with Assumption \ref{assumption:random}:
  \begin{align*}
    &e^{-\lambda_0 t}S_j^B(t)\overset{\substack{L^2\\a.s.}}{\to}(av+\mu)Y\int_0^\infty e^{-\lambda_0 s}\mathbb{E}[P_j(s|B)]\,ds=Y\lim_{s\to\infty}e^{-\lambda_0 s} \mathbb{E}[S_j^B(s)].
  \end{align*}
\end{restatable}
\begin{proof}
  See Appendix Section \ref{app:SFS_convergence}.
\end{proof}

Eq.~\ref{eq:driver_lln_time} is proved by decomposing $S_j(t)$ into the difference of two monotone processes and devising a succession of $L^2$-approximations to each monotone process. Almost-sure convergence is then established using Proposition \ref{prop:convergence_tool} of the Appendix. Eq.~\ref{eq:driver_lln_size} follows from Eq.~\ref{eq:driver_lln_time} by remarking that on $\Omega_0^\infty$, $|\tau_n-\frac{1}{\lambda_1}\log\left(n/W\right)|\ll 1$, so $\exp(-\lambda_0 \tau_n)\sim n^{-\lambda_0/\lambda_1}W^{\lambda_0/\lambda_1}$ (formally proved using Lemma \ref{lem:large_size_time} of the Appendix). A corresponding fixed detection-size result for the random fitness advance model would require a non-trivial a.s.~convergence result for $\gamma(t)Z_1^B(t)$ for some $\gamma(t)$, which is outside the scope of this manuscript\footnote{Evidence from a semideterministic approximation of the model \cite{durrettfoo2010} and Proposition \ref{prop:scaling} suggests $\gamma(t)=t^{\frac{\lambda_1(b_{\max})}{\lambda_0}}e^{-\lambda_1(b_{\max})t}$, where $\lambda_1(b)=a+b-d$ denotes the type-1 growth rate with fitness advancement $b$.}.

In a single-type neutral model, Gunnarsson et al.~\cite{gunnarsson2025} found that the fixed-size limit of the SFS is deterministic, which was attributed to the fact that there is no variability in the population size at time $\tau_n$.
Here, the two-type dynamics of mutation and selection lead to a random limit even in the fixed-size case.

Before moving on, we provide a brief comment concerning how Theorem \ref{prop:convergence_results} compares with Proposition 4.5 of Cheek and Antal \cite{cheek2018}. If we use a superscript $(n)$ momentarily to denote the parameters of our model under their large detection size-small mutation regime, their proposition says:
\[S_j^{(n)}(\tau_n^{(n)})\Longrightarrow S_j^\ast,\]
weakly where $S_j^\ast\sim \text{Poisson}(\alpha_j)$. $\alpha_j$ is the rate of the Poisson in Proposition 4.5 of \cite{cheek2018}. Hence  $S_j^{(n)}(\tau_n^{(n)})\asymp 1$ in their limiting regime whereas $S_j(\tau_n)\asymp n^{\lambda_0/\lambda_1}$ in our limiting regime, conditional on $\Omega_0^\infty$. By their remark 4.4, $\alpha_j\ordas{j}{\infty}j^{-1-\lambda_0/\lambda_1}$, matching the behavior $\lim_{s\to\infty} e^{-\lambda_0 s}\mathbb{E}[S_j(s)]\ordas{j}{\infty} j^{-1-\lambda_0/\lambda_1}$ for the scaled SFS under our large detection-size/time regime.



{\bf \em Normalized driver SFS.} We next prove that the fraction of mutations present at a given frequency $j$ converges to a deterministic value almost surely. 
Let $S_{\geq j}(t)$ denote the number of mutations present in $j$ or more cells at time $t$ for $j\in \mathbb{N}$. $S_{\geq 1}(t)$ thus represents the total number of driver clones present at time $t$. Clearly $S_{\geq j}(t)=\sum_{k=j}^\infty S_k(t)$ at large times follows a $j^{-\lambda_0/\lambda_1}$ law in mean; in the following proposition we characterize its asymptotic behavior (which has an analogous proof to that of Theorem \ref{prop:convergence_results}).

\begin{restatable}{proposition}{RelativeType1AS}
  Conditional on $\Omega_0^\infty$, for all $k\geq 0$:
  \begin{align*}
    &e^{-\lambda_0 t}S_{\geq j}(t)\overset{\substack{L^2\\a.s.}}{\to}(av+\mu)Y\int_0^\infty e^{-\lambda_0 s}\sum_{k=j}^\infty P_k(s)\,ds\\
    &n^{-\lambda_0/\lambda_1}S_{\geq j}(\tau_n)\overset{\substack{a.s.}}{\to}(av+\mu)YW^{-\lambda_0/\lambda_1}\int_0^\infty e^{-\lambda_0 s}\sum_{k=j}^\infty P_k(s)\,ds
  \end{align*}

  Consequently, a.s.~on $\Omega_0^\infty$ the fraction of total mutations present at a frequency $j$ is asymptotically:
  \begin{align}
    \lim_{t\to\infty}\frac{S_j(t)}{S_{\geq 1}(t)}=\lim_{n\to\infty}\frac{S_j(\tau_n)}{S_{\geq 1}(\tau_n)}=\frac{\int_0^\infty e^{-\lambda_0 s}P_j(s)\,ds}{\int_0^\infty e^{-\lambda_0 s}\left(1-P_0(s)\right)\,ds}=\frac{\frac{\Gamma(j)\Gamma(\lambda_0/\lambda_1+1)}{\Gamma(j+\lambda_0/\lambda_1+1)}{}_2{F}_1(\lambda_0/\lambda_1,j;j+\lambda_0/\lambda_1+1;p_1)}{\frac{\lambda_1}{\lambda_0}{}_2F_1(1,\lambda_0/\lambda_1;\lambda_0/\lambda_1+1;p_1)},\label{eq:normalized_SFS_limit}
  \end{align}
where ${}_2F_1$ is the hypergeometric function.
\label{prop:fraction_as}
\end{restatable}
\begin{proof}
  See Appendix Section \ref{app:SFS_convergence}.
\end{proof}

The limit is deterministic since the random factors $Y$ and $Y W^{-\lambda_0/\lambda_1}$ cancel, at fixed time and fixed detection size, respectively. The limit depends on the model parameters only through $\lambda_0/\lambda_1$ and $p_1$, which may enable estimation of the fitness effect $b$ from data (see Discussion).  


\subsubsection{Driver SFS asymptotics as the frequency varies with time}\label{sec:SFSVary}

When we derived the $j^{-\lambda_0/\lambda_1-1}$ power law for mutations at large frequencies in Proposition \ref{prop:SFS_mean_asymptotics}, we first let $t\to\infty$ and then $j\to\infty$. This represents examining the contribution of clones of size order 1 at large times to the SFS, and then examining the large-frequency-end of the spectrum when ``frequency'' denotes the absolute number of cells that make up the clone. However, different notions of frequency are possible. For example the largest clones are of size $Z_1^{(i)}(s)\asymp e^{\lambda_1 s}$, so that to understand their contribution to the SFS, the number of clones of exponential size -- rather than of size order $1$ -- should be investigated. Thus, we now let $j$ depend on time and study the SFS over the joint limit of $(t,j(t))$ as $t\to\infty$. In fact, in Section \ref{sec:large_frequencies}, an important part of our analysis involves studying the limit of a cumulative version of the SFS given by clones that grow at (or faster than) $j(t)\sim xe^{\lambda_1 t}$. When clones are on the order $1\ll j(t)\ll \exp(\lambda_1 t)$, they are of ``intermediate'' size. Intriguingly, we observe segregated behavior of the SFS depending on whether $j(t)$ grows faster or slower than $\exp\left(\lambda_0\lambda_1/(\lambda_0+\lambda_1)t\right)$:

\begin{restatable}{theorem}{SFSSparsity}\label{thm:SFSSparsity}
    Suppose that $j:\mathbb{R}_{\geq 0}\to \mathbb{N}$ is non-decreasing. Then, we have the following trichotomy:
    \begin{enumerate}
        \item If $j(t)\gg\exp(t\lambda_0\lambda_1/(\lambda_0+\lambda_1))$, $S_{j(t)}(t)\to 0$ in $L^1$;
        \item If $j(t)\sim C\exp(t\lambda_0\lambda_1/(\lambda_0+\lambda_1))$, then conditional on $\Omega_0^\infty$, $S_{j(t)}(t)$ converges in distribution to a random variable with characteristic function:
        \begin{align}
            \mathbb{E}[e^{YL(e^{i\theta}-1)}|\Omega_0^\infty]=\frac{q_1}{q_1-L\left(e^{i\theta}-1\right)}\label{eq:char_function_poisson}
        \end{align}
        where $L=\lim_{t\to\infty} \mathbb{E}[S_{j(t)}(t)]$ (and $L$ exists). \ref{eq:char_function_poisson} is the characteristic function of a mixture of a Poisson random variable with rate $YL|\Omega_0^\infty$;
        \item If $j(t)\ll\exp(t\lambda_0\lambda_1/(\lambda_0+\lambda_1))$, then:
        \[e^{-\lambda_0 t}j(t)^{(\lambda_1+\lambda_0)/\lambda_1}S_{j(t)}(t)\overset{L^1}{\to} YL_{j(\cdot)}\]
        where $L_{j(\cdot)}=\lim_{t\to\infty}e^{-\lambda_0 t}j(t)^{(\lambda_1+\lambda_0)/\lambda_1}\mathbb{E}\left[S_{j(t)}(t)\right]$ (and $L_{j(\cdot)}$ exists).
    \end{enumerate}
    Define $j_k^-(t)=\min\{j:S_j(t)\leq k\}$ and $j_k^+(t)=\max\{j:S_j(t)\geq k\}$. Then:
    \begin{align*}
        &\lim_{k\to\infty}\liminf_{t\to\infty}\frac{\log j_k^-(t)}{t}=\lim_{k\to\infty}\limsup_{t\to\infty}\frac{\log j_k^-(t)}{t}\\
        =&\lim_{k\to\infty}\liminf_{t\to\infty}\frac{\log j_k^+(t)}{t}=\lim_{k\to\infty}\limsup_{t\to\infty}\frac{\log j_k^+(t)}{t}=\frac{\lambda_0\lambda_1}{\lambda_0+\lambda_1}
    \end{align*}
    almost surely on $\Omega_0^\infty$.
\end{restatable}
\begin{proof}
  See Appendix Section \ref{app:SFSSparsity}
\end{proof}

Point 1 follows since $S_j(t)\geq 0$ and by showing that the exact expression for the mean of the SFS (Eq.~\ref{eq:mean_SFS}) tends to $0$ for the specified $j(t)$. Point 2 is proven by examining the characteristic function of $S_{j(t)}(t)$ and applying the succession of $L^2$ approximations used to prove Theorem \ref{prop:convergence_results}. Point 3 also uses this same succession of approximations. The quantity $j_k^-(t)$ represents the {\em first} frequency that {\em few} ($\leq k$) clones occupy, and $j_k^+(t)$ represent the {\em last} frequency that {\em many} ($\geq k$) clones occupy. To prove the last part of the theorem, we first note that $j_k^-(t)\leq j_k^+(t)+1$, so it suffices to prove that for every $\varepsilon>0$, $\lambda_0\lambda_1/(\lambda_0+\lambda_1)-\varepsilon\leq \liminf_{t\to\infty}\log j_k^-(t)/t$ and $\limsup_{t\to\infty}\log j_k^+(t)/t\leq \lambda_0\lambda_1/(\lambda_0+\lambda_1)+\varepsilon$ for $k$ sufficiently large, almost surely. We then use a Borel-Cantelli argument to show that for sufficiently large $k$ there are more than $k$ clones occupying each of the sizes $[1,\exp((\lambda_0\lambda_1/(\lambda_0+\lambda_1)-\varepsilon)t)]\cap \mathbb{N}$ and fewer than $k$ clones occupying each of the sizes $[\exp((\lambda_0\lambda_1/(\lambda_0+\lambda_1)+\varepsilon)t),\infty)\cap \mathbb{N}$ for all large $t$.

Points 1--3 show that $e^{\frac{\lambda_0\lambda_1}{\lambda_0+\lambda_1}t}$ sets a threshold for the ``sparsity'' of mutations at a particular frequency $j(t)$. Many mutations occupy larger frequencies and few occupy smaller frequencies, asymptotically in $L^1$. This is also seen via simulations (Figure \ref{fig:SFS_heatmap}). Intuitively, this threshold represents a balance between two competing forces that affect the size of $S_{j(t)}(t)$: production of driver clones and their growth to size $j(t)$ contributes towards $S_{j(t)}(t)$, while a change in the size of a clone with $j(t)$ members detracts from $S_{j(t)}(t)$. When $j(t)\sim C\exp(t\lambda_0\lambda_1/(\lambda_0+\lambda_1))$, these forces counter-balance one another and yield a Poisson limit theorem.
The last result in Theorem \ref{thm:SFSSparsity} says that almost surely on $\Omega_0^\infty$, there is a sequence of numbers $\varepsilon_k\to 0^+$ and a sequence of times $(t_k)_k$ with $t_k \to \infty$ so that $t\geq t_k$ implies:
\[j_k^-(t),j_k^+(t)\in \left[\exp\left(\left(\frac{\lambda_0\lambda_1}{\lambda_0+\lambda_1}-\varepsilon_k\right)t\right),\exp\left(\left(\frac{\lambda_0\lambda_1}{\lambda_0+\lambda_1}+\varepsilon_k\right)t\right)\right].\]
This result is an almost sure version of point 2 in that the transition from having many clones of a particular frequency to few clones of a particular frequency occurs at sizes $j(t)$ growing on order $\exp\left(\frac{\lambda_0\lambda_1}{\lambda_0+\lambda_1}t\right)$.



We may then get results at fixed detection sizes as well as almost sure variants of points 1 and 3 of Theorem \ref{thm:SFSSparsity}.
From the asymptotics of $j_k^-,j_k^+$ above, we obtain the following two results almost surely on $\Omega_0^\infty$.
\begin{restatable}{corollary}{SFSSparsityCor}\label{cor:SFSSparsity}
 \begin{enumerate}
        \item The scaled time limits of $j_k^-,j_k^+$ can be extended to a fixed detection-size:
        \begin{align*}
        &\lim_{k\to\infty}\liminf_{n\to\infty}\frac{\log j_k^-(\tau_n)}{\log n}=\lim_{k\to\infty}\limsup_{n\to\infty}\frac{\log j_k^-(\tau_n)}{\log n}\\
        =&\lim_{k\to\infty}\liminf_{n\to\infty}\frac{\log j_k^+(\tau_n)}{\log n}=\lim_{k\to\infty}\limsup_{n\to\infty}\frac{\log j_k^+(\tau_n)}{\log n}=\frac{\lambda_0}{\lambda_0+\lambda_1}
    \end{align*}
    almost surely on $\Omega_0^\infty$.
    \item For every $\varepsilon>0$ if $j_{\varepsilon}(t)\overset{t\to\infty}{\lesssim} \exp\left(\left(\frac{\lambda_0\lambda_1}{\lambda_0+\lambda_1}-\varepsilon\right)t\right)$ and $J_{\varepsilon}(t)\overset{t\to\infty}{\gtrsim} \exp\left(\left(\frac{\lambda_0\lambda_1}{\lambda_0+\lambda_1}+\varepsilon\right)t\right)$, then:
    \begin{align*}
        \lim_{t\to\infty}S_{j_\varepsilon(t)}(t)&=\infty & \lim_{n\to\infty} S_{j_\varepsilon(\tau_n)}(\tau_n)&=\infty\\
        \limsup_{t\to\infty} S_{J_\varepsilon(t)}(t)&\overset{\varepsilon\to 0^+}{\lesssim}\varepsilon^{-1} & \limsup_{n\to\infty}S_{J_\varepsilon(\tau_n)}(\tau_n)&\overset{\varepsilon\to 0^+}{\lesssim}\varepsilon^{-1}\\
        \liminf_{t\to\infty}S_{J_\varepsilon(t)}(t)&=0.
    \end{align*}
    almost surely on $\Omega_0^\infty$. If in fact $J_\varepsilon(t)\ordas{t}{\infty}\exp\left(\left(\frac{\lambda_0\lambda_1}{\lambda_0+\lambda_1}+\varepsilon\right)t\right)$ and $t\mapsto J_\varepsilon(t)$ is monotone, then:
    \begin{align*}
        \limsup_{t\to\infty} S_{J_\varepsilon(t)}(t)\ordas{\varepsilon}{0^+}\varepsilon^{-1}
    \end{align*}
    on $\Omega_0^\infty$. The asymptotics terms are uniform over the sample points in $\Omega_0^\infty$ (up to null sets).
    \end{enumerate}
\end{restatable}

\begin{proof}
    See Appendix Section \ref{app:SFSSparsity}.
\end{proof}

We remark that on $\Omega_0^\infty$, for $\varepsilon>0$ small enough such that $\frac{\lambda_0\lambda_1}{\lambda_0+\lambda_1}+\varepsilon<\lambda_1$, $\limsup_{t\to\infty}S_{J_\varepsilon(t)}(t)\geq 1$ because the production times of clones $T_i$ tend to $\infty$ as $i\to\infty$ and established clones grow at order $e^{\lambda_1 t}$. So, $S_{J(t)}(t)$ has no almost sure limit. For large $t$, the oscillations of $S_{J_\varepsilon(t)}(t)$ vary within $[0,\varepsilon^{-1}]$, up to constant factors, which increases as $\varepsilon\to 0^+$. By the last part of Corollary \ref{cor:SFSSparsity}, the upper bound of the range for these oscillations is tight (up to order).

\subsubsection{Tail of the driver SFS at fixed detection size}\label{sec:tail_SFS}
We next investigate tail characteristics of the driver SFS at fixed detection size.  As remarked in \cite{gunnarsson2021} (and as seen in Figure 1 of \cite{cheek2018}), the fixed-size SFS for a neutrally evolving population forms a point mass at $j=n$. This remains true in the case with selection (Figure \ref{fig:det_time}A).  Interestingly, the size of the point mass is asymptotically independent of the detection size $n$.
In the following we show that $S_n(\tau_n)$, conditioned on $\tau_n<\infty$, is a Bernoulli-distributed random variable. 

 
\begin{figure}[H]
	\centering
	\includegraphics[width=\textwidth]{assets/Figure_det_time.pdf}
	\caption{Site frequency spectrum at fixed detection size. \textbf{(A)} $\mathbb{E}[S_j(\tau_n)|\tau_n<\infty]$ are plotted for $n\in \{10^3,10^4,10^5\}$. The colored tick markers indicate the stopping size. \textbf{(B)} $\mathbb{E}[S_j(\tau_n)|\tau_n<\infty]$ is plotted for varying selection amounts using the stopping size $n=10^4$.}\label{fig:det_time}
\end{figure}


\begin{restatable}{lemma}{TauNSFS}
    Conditional on $\Omega_1^\infty = \cap_n \{\tau_n<\infty\}$, $S_n(\tau_n)$ converges almost surely to a Bernoulli random variable with success probability $\frac{\mathbb{P}(\text{There is only 1 established mutation}|\Omega_0^{\infty,c})(1-\pi_0)}{\pi_1}$ where $\pi_i=\mathbb{P}(\Omega_i^\infty)$. All probabilities are explicitly calculable.\label{lem:tau_n_SFS}
\end{restatable}
\begin{proof}
  See Appendix Section \ref{app:tau_n_SFS}.
\end{proof}
The intuition for Lemma \ref{lem:tau_n_SFS} is that at time $\tau_n$, $S_n(\tau_n) \in \{0,1\}$ as at most one type-1 clone can have reached size $n$; asymptotically, the event $\{S_n(\tau_n)=1\}$ agrees with the event that the type-0 population goes extinct and there is a single established type-1 clone.  The probability of this event is independent of $n$ since it is a property of the early competition dynamics of the population. Each of the probabilities in the lemma can be computed by using standard first-step analysis. Note in Figure \ref{fig:det_time} that the mass immediately prior to $n$ goes to zero; this is discussed further in the next section.

In Figure \ref{fig:det_time} we observe significant mass near the end of the tail at fixed detection size and large selective advantage $b$ (Figure \ref{fig:det_time}B). This is because for sufficiently large $b$, upon forming the first type-1 mutant clone at time $T_1$, the events that occur at times $t \in [T_1, \tau_n]$ will mostly be composed of births within that clone, so the SFS at time $\tau_n$ provides a snapshot of the type-0 population at time $T_1$. On the other hand, if $b$ is more moderate, the time between first mutation and detection is longer, and the type-0 population will have time to evolve away from its state at $T_1$, which dampens the mass in the tail. We formalize this intuition with the following proposition.
\begin{restatable}{proposition}{Freeze}
    Let $n(b):\mathbb{R}\to \mathbb{N}$ be non-decreasing and $n(b)\geq 2$.  Also define $\tau_n^{(0)}=\inf\{t\geq 0:Z_0(t)=n\}$. We have the following dichotomy:
    \begin{itemize}
        \item If $\log n(b)\llas{b}{\infty}b$, then $\mathbb{P}(\tau_{n(b)}<\infty,Z_0(\tau_{n(b)})=Z_0(T_1)|T_1\leq \tau_{n(b)}^{(0)},T_1<\infty)\to 1$ as $b\to\infty$. 
        \item If $\log n(b)\ggas{b}{\infty}b$, then $\mathbb{P}(\tau_{n(b)}<\infty,Z_0(\tau_{n(b)})=Z_0(T_1)|T_1\leq \tau_{n(b)}^{(0)},T_1<\infty)\to \mathbb{P}(Z_0(T_1)=0|T_1<\infty)$ as $b\to\infty$.
    \end{itemize}
\end{restatable}
\begin{proof}
  See Appendix Section \ref{app:freeze}.
\end{proof}

There is a slight asymmetry in these results, in that for slow growing detection thresholds, the probability above tends to $\mathbb{P}(Z_0(T_1)=0|T_1<\infty)$ as opposed to zero. This is for the degenerate reason that the type-0 population once extinct remains extinct thereafter. Indeed the event $\{Z_0(T_1)=0\}$ conditional on $T_1<\infty$ occurs with positive probability, if there is a single type-0 cell in the population just before time $T_1$, and it produces a type-1 cell via replication-free mutation at time $T_1$. But the above proposition says that as $b\to\infty$, if the detection threshold grows slower than exponentially in $b$, with high probability the type-0 population remembers its position at $T_1$ at time $\tau_{n(b)}$. Otherwise, with high probability the type-0 population forgets its position at $T_1$ at time $\tau_{n(b)}$. Hence, as the selection amount is made large and the detection threshold kept constant in Figure \ref{fig:det_time}B, the tail of the SFS can reflect the size of the wild-type population at the time of the first type-1 clone.


\subsection{Frequencies of large clones}\label{sec:large_frequencies}
To gain further insights into the behavior of the SFS at the largest frequencies, we now investigate the number of ``large families'' of driver clones using the {\em relative} driver SFS, which gives the number of distinct clones above a certain relative frequency $f\in (0,1)$, $R_f(t)=\sum_{1\leq i\leq M(t)}\mathbf{1}_{Z_1^{(i)}(t)>Z_1(t)f}=\sum_{j:j>Z_1(t)f}S_j(t)$. This version of the SFS has previously been studied by Tung and Durrett \cite{tung}.

\subsubsection{Constructing an approximation to the relative site frequency spectrum}
To investigate $R_f(t)$, we first examine a simpler quantity counting the number of ``large families'' of mutations.  Specifically, define $\widetilde{R}_x(t)=\sum_{j> xg(t)}S_j(t)$ for $x>0$ as a cumulative SFS counting clones of size $> xg(t)$ at time $t$. This type of quantity has previously been studied by Bonnet and Leman \cite{bonnet2024}.

Under a deterministic fitness advance, it should be intuitively clear that $g(t)=e^{\lambda_1 t}$ appropriately sets the scale of large mutants. 
With this scaling, $\widetilde{R}_x(t)$ is a semi-deterministic version of $R_f(t)$, in the sense $|\widetilde{R}_{fW}(t)-R_f(t)|\llas{t}{\infty}1$ since $Z_1(t)\sim We^{\lambda_1 t}$.
In the random fitness advance model, then under Assumption \ref{assumption:random}, in fact $g(t)=t^{-\frac{\lambda_1(b_{\max})}{\lambda_0}}e^{\lambda_1(b_{\max})t}$ sets the scale, in the following sense:

\begin{restatable}{proposition}{Scaling}
  Fix $x>0$. With a deterministic fitness advance and defining $g(t)=e^{\lambda_1 t}$:
  \[\lim_{t\to\infty}\mathbb{E}\left[\sum_{j> xg(t)}S_j(t)\right]=\frac{(av+\mu) q_1^{1-\lambda_0/\lambda_1}x^{-\lambda_0/\lambda_1}}{\lambda_1}\Gamma(\frac{\lambda_0}{\lambda_1},q_1 x).\]
With a random fitness advance under Assumption \ref{assumption:random} and defining $g_2(t)=t^{-\frac{\lambda_1(b_{\max})}{\lambda_0}}e^{\lambda_1(b_{\max})t}$ (here, $\lambda_1(b')=a+b'-d$ denotes the type-1 net growth rate with fitness advancement $b'$):
  \[\lim_{t\to\infty}\mathbb{E}\left[\sum_{j> xg_2(t)}S_j^B(t)\right]\]
  converges to a finite strictly positive limit.\label{prop:scaling}
\end{restatable}
\begin{proof}
  See Appendix Section \ref{app:scaling}.
\end{proof}

Interestingly, the scaled version $t^{\frac{\lambda_1(b_{\max})}{\lambda_0}}e^{-\lambda_1(b_{\max})t}Z_1^{B}(t)$ converges in distribution under a semideterministic version of our model (see Durrett et al.~\cite{durrettfoo2010}). 

\subsubsection{Almost sure convergence of the relative site frequency spectrum and small frequency limit}\label{sec:convergence_relative_SFS}
In what follows, we restrict ourselves to the model with deterministic fitness advance $b$ and set $g(t)\equiv e^{\lambda_1 t}$ in defining $\widetilde{R}_x(t)$ to study ``large families''. 
We can get a simple (but relatively opaque) convergence result for $R_f$ (and $\widetilde{R}_x$) at large times:
\begin{restatable}{proposition}{FrequencyLimit}
  Let $f\in (0,1)$ and $x>0$. Then the following a.s.~convergence result holds conditional on $\Omega_0^\infty$:
  \begin{align}
    &\lim_{t\to\infty} R_f(t)=\sum_{i\in \mathbb{N}}\mathbf{1}_{e^{-\lambda_1 T_i}Y_1^{(i)}>Wf}\label{eq:frequency_limit}\\
    &\lim_{t\to\infty} \widetilde{R}_x(t)=\sum_{i\in \mathbb{N}}\mathbf{1}_{Y_1^{(i)}> xe^{\lambda_1 T_i}}.
  \end{align}

  Furthermore, the limits are a.s.~finite.\label{prop:frequency_limit}
\end{restatable}
\begin{proof}
  See Appendix Section \ref{app:frequency_limit}
\end{proof}

Note that $e^{-\lambda_1 T_i}Y_1^{(i)}/W$ represents the asymptotic proportion of the population that the $i$th type-1 clone makes up. Additionally, it should be remarked that $R_f(t)\leq \lceil 1/f\rceil$, so that $R_f(t)$ converges boundedly and thus in $L^p$ for $p>0$ as $t\to\infty$. A trichotomy follows from the above proposition:
\begin{restatable}{corollary}{ScaleSumSFS}
    Denote $k(n):\mathbb{N}\to \mathbb{N}$ such that $k(n)\ll n$ as $n\to\infty$, $G^+(t):\mathbb{R}_{\geq 0}\to \mathbb{R}_{\geq 0}$ such that $G^+(t)\gg e^{\lambda_1 t}$ as $t\to\infty$, and $G^-(t):\mathbb{R}_{\geq 0}\to \mathbb{R}_{\geq 0}$ such that $G^-(t)\ll e^{\lambda_1 t}$ as $t\to\infty$. Then conditional on $\Omega_0^\infty$, for every $x\in (0,1)$:


  \begin{align*}
    \sum_{j=n-k(n)}^{n} S_j(\tau_n)&\overset{a.s.}{\to} 0 &\sum_{j> G^+(t)}S_j(t)&\overset{a.s.}{\to}0\\
    \sum_{j=\lceil fn\rceil }^{n} S_j(\tau_n)&\overset{a.s.}{\to} \sum_{i\in \mathbb{N}}\mathbf{1}_{e^{-\lambda_1 T_i}Y_1^{(i)}>Wf}& \sum_{j> xe^{\lambda_1 t}}S_j(t)&\overset{a.s.}{\to}\sum_{i\in \mathbb{N}}\mathbf{1}_{Y_1^{(i)}> xe^{\lambda_1 T_i}}\\
    \sum_{j=k(n)}^{n} S_j(\tau_n)&\overset{a.s.}{\to} \infty & \sum_{j> G^-(t)}S_j(t)&\overset{a.s.}{\to}\infty
  \end{align*}\label{cor:scale_sum_SFS}
\end{restatable}
\begin{proof}
  See Appendix Section \ref{app:frequency_limit}.
\end{proof}
If we look at the large detection-size limit column first, we see that there is an asymptotically vanishing contribution to the cumulative SFS from clones with size near $n$, an order 1 contribution from clones of size $fn$ for $f\in (0,1)$, and an infinite contribution from clones of small size. The middle expression gives a connection between the fixed-size SFS viewed at absolute frequencies $j = \lceil fn \rceil$ and above and the relative SFS observed at a fixed time, showing that they converge to the same limit.
We see similar results for the large-time limit column. 
Finally, compare the above corollary with Theorem \ref{thm:SFSSparsity}. The scale that sets the sparsity of mutations \textit{at a particular frequency} $j(t)$ is $j(t)\sim e^{\frac{\lambda_0\lambda_1}{\lambda_0+\lambda_1}t}$ and the scale that sets the sparsity of mutations of \textit{at least} a particular frequency $j(t)$ is $j(t)\sim e^{\lambda_1 t}$.

We next prove a convergence result for $\widetilde{R}_x(t)$ and $R_f(t)$ that gives a better insight into the limit in Proposition \ref{prop:frequency_limit}.
\begin{restatable}{theorem}{ASRelativeSFS}
   On $\Omega_0^\infty$, we obtain the following small frequency scaling a.s.:
  \[\lim_{x\to 0^+}\lim_{t\to\infty}x^{\lambda_0/\lambda_1}\widetilde{R}_x(t)=\lim_{f\to 0^+}\lim_{t\to\infty}f^{\lambda_0/\lambda_1}W^{\lambda_0/\lambda_1}R_f(t)=\frac{(av+\mu) q_1^{1-\lambda_0/\lambda_1}}{\lambda_1}\Gamma(\frac{\lambda_0}{\lambda_1})Y\]
  \label{prop:ASRelativeSFS}
\end{restatable}
\begin{proof}
  See Appendix Section \ref{app:relSFS}.
\end{proof}

\begin{figure}[H]
	\centering
	\includegraphics[width=\textwidth]{assets/RelativeSFS.pdf}
	\caption{Large time relative plots of $R_f(t)$. }\label{fig:relative_sfs}
\end{figure}

The proof of the large-time, small $x$ asymptotics for $\widetilde{R}_x(t)$ uses similar techniques as in the proof of Theorem \ref{prop:convergence_results} and Proposition \ref{prop:fraction_as}. The asymptotics for $R_f(t)$ follow from the approximation $R_f(t)\approx \widetilde{R}_{x_f}(t)$ with $x_f=fW$ and taking $x_f\to 0$. Hence, conditional on $\Omega_0^\infty$, we obtain the following power law result:
\begin{align}
  \lim_{t\to\infty}R_f(t)\overset{f\to 0^+}{\sim} f^{-\lambda_0/\lambda_1}\frac{(av+\mu) q_1^{1-\lambda_0/\lambda_1}}{\lambda_1}\Gamma(\frac{\lambda_0}{\lambda_1})YW^{-\lambda_0/\lambda_1}\label{eq:relative_large_scaling}
\end{align}
Note that if we substitute $f^{-\lambda_0/\lambda_1} \Gamma(\lambda_0/\lambda_1) $ with $\frac{d}{df} \left( f^{-\lambda_0/\lambda_1} \Gamma(\lambda_0/\lambda_1)\right) = -  f^{-\lambda_0/\lambda_1-1} \Gamma(\lambda_0/\lambda_1+1)$ and ignore the stochastic terms, we recover the $j^{-\lambda_0/\lambda_1-1}$ power law in Eq.~\ref{eq:large_j_equivalent} of Proposition \ref{prop:SFS_mean_asymptotics} (irrespective of sign). We can thus say that Eq.~\ref{eq:relative_large_scaling} is a cumulative version of the power law result we derived before.

Tung and Durett analyze the distribution of $\lim_{t\to\infty}R_f(t
)$ (Eq. \ref{eq:frequency_limit}) using results from Pitman and Yor concerning the distribution of normalized points in a Poisson point process over the positive real line \cite{tung,pitman1997}. In particular, they analyze a semideterministic approximation of our model, where the type-$0$ population grows as $Z_0^*(t)=Ye^{\lambda_{0,0} t}$, and the mutation times $\{T_i\}$ are generated as an inhomogenous Poisson point process with intensity $\Lambda(dt)=vYe^{\lambda_{0,0} t}\,dt$ conditional on $Y$\footnote{Recall that $\lambda_{0,v}=a(1-v)-d$ denotes the growth rate with no replication-independent mutations and with the replication-dependent mutation probability explicitly shown. Note that $\lambda_{0,0}$ is used as their semi-deterministic model does not incorporate mutations into the type-0 growth rate}. They find that the tail follows a $f^{-\lambda_{0,0}/\lambda_1}$ shape in expectation. They obtain:
\begin{align}
  \mathbb{E}[\lim_{t\to\infty}R_f(t)]\overset{T.D.}{=}\frac{\sin\left(\pi\frac{\lambda_{0,0}}{\lambda_1}\right)}{\pi\frac{\lambda_{0,0}}{\lambda_1}}\left(\frac{1}{f}-1\right)^{\lambda_{0,0}/\lambda_1},\label{eq:tung_durrett}
\end{align}
where $\overset{T.D.}{=}$ denotes equality in the Tung and Durrett model. In the small mutation limit, we actually recover the above up to leading order as $f\to 0^+$.
\begin{restatable}{proposition}{SmallMutationRelativeSFS}
  If mutations are only associated with births:
  \[\lim_{v\to0^+}\mathbb{E}\left[\lim_{f\to0^+}\lim_{t\to\infty}f^{\lambda_0/\lambda_1}R_{f}(t)|\Omega_0^\infty\right]=\lim_{v\to 0^+}\mathbb{E}\left[\frac{av q_1^{1-\lambda_0/\lambda_1}}{\lambda_1}\Gamma(\frac{\lambda_0}{\lambda_1})YW^{-\lambda_0/\lambda_1}|\Omega_0^\infty\right]=\frac{\sin\left(\pi\frac{\lambda_{0,0}}{\lambda_1}\right)}{\pi\frac{\lambda_{0,0}}{\lambda_1}}.\]\label{prop:small_mutation_relative_SFS}
\end{restatable}
\begin{proof}
  See Appendix Section \ref{app:small_mutation_relative_SFS}.
\end{proof}

The intuition behind this correspondence lies in a deeper connection between our two-type branching process model in the small mutation limit and the semi-deterministic model Tung and Durrett study (for example, for the case $\mu=0$, see Cheek an Antal \cite{cheek2018} Proposition 5.1; also see the Discussion for more commentary on this).

By Figure \ref{fig:relative_sfs}, we observe quite good agreement between simulations and Tung and Durett's result (Eq. \ref{eq:tung_durrett}) when mutation rates are small. Our leading order asymptotic (Eq. \ref{eq:relative_large_scaling} and Proposition \ref{prop:small_mutation_relative_SFS}) is good for small $f$ but is an overestimate of both the simulation and Tung and Durett result for large $f$.

\subsubsection{The largest clone}\label{sec:smallmut_largeclone}

We conclude our investigation of the large-frequency end of the SFS by examining the size of the largest clone.

\begin{restatable}{proposition}{LargestClone}
  We have the following large time and large detection-size asymptotics for the largest clone on $\Omega_0^\infty$:
  \begin{align*}
    &e^{-\lambda_1 t}\max_{T_i\leq t}\left(Z_1^{(i)}(t-T_i)\right)\overset{a.s.}{\to}\sup_{1\leq i<\infty}\left(e^{-\lambda_1 T_i}Y_1^{(i)}\right)\\
    &n^{-1}\max_{T_i\leq \tau_n}\left(Z_1^{(i)}(\tau_n-T_i)\right)\overset{a.s.}{\to}\sup_{1\leq i<\infty}\left(e^{-\lambda_1 T_i}Y_1^{(i)}\right)/W.\\
  \end{align*}
    In particular:
    \begin{align*}
    &\max_{T_i\leq t}\left(Z_1^{(i)}(t-T_i)\right)/Z_1(t)\overset{a.s.}{\to}\sup_{1\leq i<\infty}\left(e^{-\lambda_1 T_i}Y_1^{(i)}\right)/W\\
    &\max_{T_i\leq \tau_n}\left(Z_1^{(i)}(\tau_n-T_i)\right)/Z_1(\tau_n)\overset{a.s.}{\to}\sup_{1\leq i<\infty}\left(e^{-\lambda_1 T_i}Y_1^{(i)}\right)/W.\\
  \end{align*}\label{prop:largest_clone_convergence}
\end{restatable}
\begin{proof}
  See Appendix Section \ref{app:largest_clone_convergence}.
\end{proof}
Note that:
\[\lim_{t\to\infty}\sup\text{supp}(R_f(t))=\sup\text{supp}(\lim_{t\to\infty}R_f(t))=\sup_{1\leq i<\infty}\left(e^{-\lambda_1 T_i}Y_1^{(i)}\right)/W.\]
We can get further insights into the time limits in Proposition \ref{prop:largest_clone_convergence} when the mutation rate is small. In particular, define $\alpha=\lambda_{0,0}/\lambda_1$ and $\gamma=\left(\frac{q_1^{1-\alpha}a\Gamma(\alpha)}{\lambda_1 q_0}\right)^{1/\alpha}.$ In the theorem below $V$ denotes the largest element of a Poisson-Dirichlet distribution with parameters $(\alpha,0)$ and $Z$ denotes a random variable with a log-logistic distribution with scale $\gamma$ and shape $\alpha$.
\begin{restatable}{proposition}{LargestCloneSmall}
 If mutations are only attached to births, then conditional on $\Omega_0^\infty$:
  \begin{align*}
    \sup_{1\leq i<\infty}\left(e^{-\lambda_1 T_i}Y_1^{(i)}\right)\overset{v\to 0^+}{\Rightarrow}Z& & \sup_{1\leq i<\infty}\left(e^{-\lambda_1 T_i}Y_1^{(i)}\right)/W\overset{v\to 0^+}{\Rightarrow}V.
  \end{align*}
  \label{prop:largest_clone_small}
\end{restatable}
\begin{proof}
  See Appendix Section \ref{app:largest_clone_small}.
\end{proof}


\begin{remark}
  A complicated expression for the CDF of $V$ can be obtained explicitly from Proposition 20 of Pitman and Yor \cite{pitman1997}. All moments of $V$ are known using Proposition 17 of Pitman and Yor, with the first moment being:
  \[\mathbb{E}[V]=\int_0^\infty t^{-\alpha}e^{-t}\left(\Gamma(1-\alpha)+\Gamma(-\alpha,t)\right)^{-1}\,dt.\]
Furthermore, as long as $\frac{\sin(\pi \alpha)}{\pi\alpha}\geq \frac{1}{2}$, ($0<\alpha\lesssim 0.603$, that is $b$ is sufficiently large), the median of $V$ is:
\[\frac{1}{1+\left(\frac{\pi\alpha}{2\sin(\pi\alpha)}\right)^{1/\alpha}}\]
Finally, $V$ is small very infrequently:
\[\mathbb{P}(V\leq x)\lsimas{x}{0^+}e^{-cx^{-1}}\]
where $c$ can be explicitly determined. We give proofs of some of the statements in Appendix Section \ref{app:poisson_dirichlet}.\label{rem:poisson_dirichlet}
\end{remark}

In Durrett et al.~\cite{durrettgenetics2010}, under a semi-deterministic version of our model, the degree to which the largest clone dominates asymptotically is also determined via the Laplace transform of $1/V$.

\begin{remark}
    In the propositions where we restrict mutation accumulation to occur during births and look in the small mutation limit (Propositions \ref{prop:small_mutation_relative_SFS} and \ref{prop:largest_clone_small}), it is likely possible to include replication-independent mutations and specify the small mutation limit to be $v,\mu\to 0^+$ such that $v/\mu\to \theta$ where $\infty>\theta>0$.
\end{remark}

\section{Discussion}

In this work, we have characterized the driver SFS $S_j(t)$ and the relative driver SFS $R_f(t)$ under various limiting regimes, for a two-type branching process in which driver mutations confer a selective advantage.   
See Figure \ref{fig:paper_snapshot} for a graphical summary of our results and the various characteristics of the SFS they address.

\begin{figure}[H]
	\centering
	\includegraphics[width=\textwidth]{assets/FigureSnapshot.pdf}
	\caption{Diagram of various results from this paper and the corresponding region of $S_j(t)$ or $R_f(t)$ they describe. $\approx$ denotes asymptotic equivalence $\sim$ up to some (possibly random) constant. The large arrows depict how $(j(t),S_{j(t)})$ varies as $t$ ranges in $\{t_1,t_2\}$. The encircled regions represent the end of the support of $j\to S_j(t)$ and $f\to R_f(\infty)$.}\label{fig:paper_snapshot}
\end{figure}

\textbf{Comparing how clones of different sizes contribute to the SFS.} The connection between $S_j(t)$ and $R_f(t)$ is that they both measure the number of clones existing at a certain frequency $j$ or $\geq f$, respectively; the difference being the notion of frequency used -- $j$ is in units of number of cells whereas $f$ is in units of the fraction of the type-1 population. In other words, $S_j(t)$ considers the contribution of clones of size $\mathcal{O}(1)$ and $R_f(t)$ considers the contribution of clones of size $\mathcal{O}(\exp(\lambda_1 t))$ 

Taking the $t\to\infty$ limit of $S_j(t)$ and $R_f(t)$ (appropriately scaled) reveals that both have power-law tails $j^{-\lambda_0/\lambda_1-1}$ and $f^{-\lambda_0/\lambda_1}$ as $j\to\infty$ and $f\to0^+$ respectively. There is an intuitive heuristic for why we observe these tails. Suppose that $\gamma_j$ is the initiation time of a clone that grows to size $j$ at the detection time $t$. This clone has time $t-\gamma_j$ to grow to size $j$, so solving $\exp(\lambda_1 (t-\gamma_j))\approx j$ we obtain $\gamma_j\approx t-(1/\lambda_1)\log j$. If a clone is of size $j$, it will take time $\approx 1/(j(a+b+d))$ for the clone to change its size. Thus, the number of clones of size $j$ at time $t$ is:
\[S_j(t)\approx \int_{\gamma_j}^{\gamma_j+1/(j(a+b+d))}(av+\mu)q_1e^{\lambda_0 s}\,ds\approx Ce^{\lambda_0 t}j^{-1-\lambda_0/\lambda_1}.\]

Note then that clones that contribute to the SFS at absolute frequency $j$ occur at $\mathcal{O}(1)$ time units before $t$. For the relative SFS, we instead consider $\gamma_f$ as the initiation time of a clone that grows to size $ fZ_1(t)\approx fWe^{\lambda_1 t}$. Thus, solving as above, we obtain that $\gamma_f= -\frac{1}{\lambda_1}\log fW$ approximately. Clones that form earlier also contribute so that:
\[R_f(t)\approx \int_{0}^{\gamma_f}(av+\mu)q_1 e^{\lambda_0 s}\,ds\approx CW^{-\lambda_0/\lambda_1}f^{-\lambda_0/\lambda_1}.\]
The clones that contribute to the relative SFS at frequency $\geq f$ occur within a $\mathcal{O}(1)$ window after the start of the process at time $0$.

Finally, when counting clones of size on order $1\ll g(t)\ll \exp(t\lambda_0\lambda_1/(\lambda_0+\lambda_1))$, Theorem \ref{thm:SFSSparsity}, Point 3 and Corollary \ref{cor:sparsity_part1_2_3}, show that $\exp(-\lambda_0t)g(t)^{1+\lambda_0/\lambda_1}S_{\lceil xg(t)\rceil}(t)\to x^{-1-\lambda_0/\lambda_1}YL_{\lceil xj(\cdot)\rceil}=YCx^{-1-\lambda_0/\lambda_1}$ for $x\in (0,\infty)$, with:
\[C=\frac{\lambda_1(av+\mu)}{(a+b)^2}\left(\frac{a+b}{a+b-d}\right)^{\frac{\lambda_1+\lambda_0}{\lambda_1}}\Gamma(1+\lambda_0/\lambda_1).\]
Thus, we recover power law tails when taking $x\to0^+$ and $x\to\infty$. So, the power-law tails are in some sense universal irrespective of whether we consider frequencies of $\mathcal{O}(1)$, intermediate, or large clones.

\textbf{Estimating parameters from the SFS.} We found that in principle, the SFS can be used to discern certain evolutionary parameters of this continuous-time linear birth-death process. For example, the tails of 
$$\lim_{t\to\infty}S_j(t)/S_{\geq 1}(t)=\lim_{n\to\infty}S_j(\tau_n)/S_{\geq 1}(\tau_n)
$$ and 
$$\lim_{t\to\infty}R_f(t)=\lim_{n\to\infty}R_f(\tau_n)$$
on $\Omega_0^\infty$ follow power-laws $j^{-\lambda_0/\lambda_1-1}$ and $f^{-\lambda_0/\lambda_1}$ as $j\to\infty$ and $f\to 0^+$ respectively, which can be used to estimate the selection effect (i.e. relative increase in fitness between type-0 and type-1 cells).  

We can also construct a consistent estimator for this fitness increase $b$ using the small $j=1$ region of the SFS, assuming all other parameters are known. For example, the parameters of the wild-type population may be known, but the fitness increase $b$ associated with mutation may not be. Define $F(b)=\frac{\int_0^\infty e^{-\lambda_0 s}P_1(s)\,ds}{\int_0^{\infty}e^{-\lambda_0 s}(1-P_0(s))}$. Using the explicit form of $P_j$ (for example provided in \cite{durrettbranching2015}), $P_1$ and $P_0$ are decreasing in $b$, so in total $F(b)$ is monotonically decreasing and is continuous. Hence, $F^{-1}$ is defined and continuous\footnote{That is if $F^{-1}$ is taken as a function over domain $[0,F(0)]$. We can extend $F^{-1}$ continuously onto $[0,1]$ by defining it over $(F(0),1]$ to be $0$}. We may then consider the estimator $\widehat{b}(t)=F^{-1}(f_1(t))$ where $f_1(t):=S_1(t)/S_{\geq 1}(t)$ is the proportion of type-1 lineages at time $t$ of size $1$. By Proposition \ref{prop:fraction_as} on $\Omega_0^\infty$, $f_1(t)\to F(b)$ a.s. as $t\to\infty$ and $f_1(\tau_n)\to F(b)$ as $n\to\infty$, so by continuous mapping $\widehat{b}(t)\to b$ a.s. and $\widehat{b}(\tau_n)\to b$ a.s.

\textbf{A comment on the semi-deterministic model.} We note that Durrett and Moseley \cite{durrett2010} studied an approximation to our model where the type-0 population grows semi-deterministically; this model was later used to investigate tumor heterogeneity, allow for random fitness increases, and quantify the relative SFS with selection \cite{durrett2010,durrettfoo2010,durrett2013,tung}. The justification for the semi-deterministic model is that for small mutation rates, the typical time of forming mutations (conditional on $\Omega_0^\infty$) occurs when the type-0 population is large and thus approximately growing according to $Ye^{\lambda_0 t}$. Our results here recover some quantities (like $R_f(t)$ and the size of the largest clone) computed in prior semi-deterministic models as small-mutation limits of the fully stochastic model (Proposition \ref{prop:small_mutation_relative_SFS} and Proposition \ref{prop:largest_clone_small}).

\textbf{Future work.} The model considered in this work considers a single driver mutation, and does not allow for sequential drivers. In addition, we consider only the SFS of the driver mutation population here, not the SFS of the full tumor population.  

In upcoming work, we extend our analyses to study different components of the site frequency spectrum, including the hitchhiker SFS $S_j^h(t)$, in a more general multi-type population. Additionally, we will explore second-order asymptotics and rare-event deviations in the driver SFS.

\section{Acknowledgements}
The work of K. Leder and J. Foo were partially supported by NSF DMS award 2526511. The work of J. Foo and A. Ahmed were partially supported by CMMI 2228034.  

\printbibliography

\appendix

\section{Appendix}

\subsection{Useful tools in our analysis}

\subsubsection{$\sigma$-filtrations used in our analyses}

Here, we define some important $\sigma$-filtrations

\begin{itemize}
    \item Let $\mathcal{F}_t^{\leq i}$ for $i\in \{0,1\}$ denote the smallest $\sigma$-algebra generated either (if $i=0$) by the type-0 process in the time interval $[0,t]$ or (if $i=1$) by the type-0 and type-1 process in the time interval $[0,t]$. The $i=0$ case is most often used, so that we also write $\mathcal{F}_t:=\mathcal{F}_t^{\leq 0}$. We also let $\mathcal{F}_\infty^{\leq i}=\sigma\left(\bigcup_{t\geq 0}\mathcal{F}_t^{\leq i}\right)$. Once again, for shorthand we write $\mathcal{F}_\infty:=\mathcal{F}_\infty^{\leq 0}$.
    \item Let $\mathcal{B}_t^{\leq i}$ for $i\in \{0,1\}$ denote the smallest $\sigma$-algebra containing $\mathcal{F}_t^{\leq i}$ as well as the timings of the type-1 mutations that occur within the interval $[0,t]$. We also write $\mathcal{B}_t:=\mathcal{B}_t^{\leq0}$. We also let $\mathcal{B}_\infty^{\leq i}=\sigma\left(\bigcup_{t\geq 0}\mathcal{B}_t^{\leq i}\right)$. Once again, for shorthand we write $\mathcal{B}_\infty:=\mathcal{B}_\infty^{\leq 0}$.
    \item Let $\mathcal{I}_t$ be the smallest $\sigma$-algebra containing $\mathcal{B}_t^{\leq 1}$ as well as the identities of the clones to which each type-1 cell belongs to. Finally, define $\mathcal{I}_\infty=\sigma\left(\bigcup_{t\geq 0}\mathcal{I}_t\right)$.
\end{itemize}

Clearly, $\mathcal{F}_t^{\leq i}\subseteq \mathcal{B}_t^{\leq i}\subseteq \mathcal{I}_t$. $(Z_0(t),Z_1(t))$ is adapted to $\mathcal{F}_t^{\leq 1}$. $(Z_0(t),Z_1^B(t))$ and $S_j(t)$ are adapted to $\mathcal{I}_t$

\subsubsection{An alternate representation of the model}\label{app:alt_representation_model}

There is an alternate construction of the model we refer to in some places of the appendix, which is equal in distribution to the model introduced in Section \ref{sec:model_notation}. Start with a one-type linear birth-death process with birth rate $a(1-v)$ and death rate $d+\mu$. Equivalently, each type-0 cell has an independent exponential clock with rate $a(1-v)+d+\mu$. Once the clock rings, the cell divides with probability $a(1-v)/(a(1-v)+d+\mu)$ and dies otherwise. We let $\{D_i\}$ denote the timings of all death events in this type-0 process.

Now, we describe the initiation times of type-1 clones. These type-1 clones, conditional on their initiation, grow as branching processes with birth rate $a+b$ and death rate $d$, independently from the other type-1 clones and from the type-0 process. To every cell, attach an independent homogeneous Poisson Point Process of rate $av$ over its lifetime. Conditional on $\mathcal{F}_\infty$, the summation of these points processes over all cells yields $\{T_i'\}$, a Poisson Point Process with rate $av Z_0(s)$. These mark the initiation times of the type-1 clones produced in a replication-dependent manner. Let $\{T_i''\}$ be an independent $\mu/(d+\mu)$-thinning of $\{D_i\}$. Then, $\{T_i''\}$ marks the initiation times of the type-1 clones produced in a replication-free manner. Then, $\{T_i\}=\{T_i'\}\cup \{T_i''\}$ are the initiation times of all type-1 clones.


\subsubsection{A comparison criterion for almost sure convergence}

Here, we reformulate Proposition 4 of \cite{gunnarsson2025} (which extends analysis used in Theorem 21 of Harris \cite{harris1963}) to work with more general classes of stochastic processes.

\begin{proposition}
  Let $(\Omega,\mathcal{F},\mathbb{P})$ be a probability space and consider jointly-measurable, stochastic processes $X_t$ and $Y_t$. We further take $X_t$ to be non-negative and non-decreasing a.s. Suppose there exists $\lambda>0$ such that $e^{-\lambda t}Y_t \to \widetilde{Y}_\infty$ almost surely, where $\widetilde{Y}_\infty$ is a.s. finite and non-negative. If, in addition, the following condition is satisfied:

  \[\int_0^\infty \mathbb{E}[e^{-\lambda t}|X_t-Y_t|]dt<\infty\]
then $e^{-\lambda t}X_t\to \widetilde{Y}_\infty$ a.s.\label{prop:convergence_tool}
\end{proposition}

\begin{proof}
It suffices to show that $\limsup_{t\to\infty} e^{-\lambda t}(X_t-Y_t)>0$ and $\liminf_{t\to\infty} e^{-\lambda t}(X_t-Y_t)<0$ both occur on sets of probability 0. Note that the integrability condition implies that a.s. $\int_0^\infty e^{-\lambda t}|X_t-Y_t|<\infty$. 

First consider sample paths such that $\limsup_{t\to\infty} e^{-\lambda t}(X_t-Y_t)>0$. For each such path there exists $\delta>0$ and a sequence of times $\{t_k;k\geq 1\}$ with $t_{k+1}-t_k\geq 1$ for which $e^{-\lambda t_k}(X_{t_k}-Y_{t_k})>4\delta$. Due to the almost sure convergence of $\{e^{-\lambda t}Y_t;t\geq 0\}$ there is a $T$ where $t>T$ implies $|e^{-\lambda t}Y_t-\widetilde{Y}_{\infty}|\leq \delta$. For any $k$ sufficiently large (so that $t_k\geq T$):

  \[e^{-\lambda t_k}X_{t_k}>4\delta+e^{-\lambda t_k}Y_{t_k}\geq 3\delta +\widetilde{Y}_{\infty}.\]
Next since $\{X_t;t\geq 0\}$ is monotone we know that for any $0<\Delta<1$, if $t\in [t_k,t_k+\Delta]$:
  \[e^{-\lambda t}X_t\geq e^{-\lambda (t-t_k)}e^{-\lambda t_k}X_{t_k}\geq e^{-\lambda \Delta}(3\delta+\widetilde{Y}_\infty).\]
We can then choose $\Delta$ sufficiently small such that:
  \[e^{-\lambda \Delta}(3\delta+\widetilde{Y}_\infty)\geq (2\delta + \widetilde{Y}_\infty).\]
Then combining with the estimate $\widetilde{Y}_\infty\geq e^{-\lambda t}Y_t-\delta$, we finally obtain:
  \[e^{-\lambda t}(X_t-Y_t)\geq \delta,\]
for $t$ in $\bigcup_{k=K}^\infty[t_k,t_k+\Delta]$. Therefore on the set $\{\omega\in \Omega; \limsup_{t\to\infty}e^{-\lambda t}(X_t-Y_t)>0\}$
  \[\int_0^\infty e^{-\lambda t}|X_t-Y_t|\,dt\geq \int_{\bigcup_{k=K}^\infty[t_k,t_k+\Delta]}\delta=\infty.\]
In particular, $\limsup_{t\to\infty}e^{-\lambda t}(X_t-Y_t)\leq 0$ a.s.

We now consider sample paths such that $\liminf_{t\to\infty}e^{-\lambda t}(X_t-Y_t)<0$. For each such path there exists  $\delta>0$ and $t_k$ a sequence with $t_{k+1}-t_k\geq 1$ (and $t_1\geq 1$) such that $e^{-\lambda t_k}(X_{t_k}-Y_{t_k})<-5\delta$. Fix any $0<\Delta<1$, (and take $t\in [t_k-\Delta,t_k]$) then:

  \[-5\delta >e^{-\lambda t_k}(e^{\lambda\Delta}X_t-e^{\lambda \Delta} Y_{t_k})\geq e^{-\lambda t_k}(e^{\lambda (t_k-t)}X_t-e^{\lambda\Delta}Y_{t_k})=e^{-\lambda t}X_t-e^{\lambda\Delta}e^{-\lambda t_k}Y_{t_k}.\]
Then for $t$ sufficiently large we can again use the a.s. convergence of $\{e^{-\lambda t} Y_t;t\geq 0\}$ to find $T<\infty$ such that $t\geq T$ implies $|\widetilde{Y}_\infty-e^{-\lambda t}Y_t|<\delta$ and thus
  \[-5\delta +\delta e^{\lambda\Delta}\geq e^{-\lambda t}X_t - e^{\lambda\Delta}\widetilde{Y}_\infty.\]
Take $\Delta$ sufficiently small such that $|e^{\lambda\Delta}\widetilde{Y}_\infty-\widetilde{Y}_\infty|<\delta$ and $\delta e^{\lambda\Delta}<2\delta$. Combining these bounds yields 

  \[-2\delta\geq e^{-\lambda t}X_t-\widetilde{Y}_\infty.\]
We can then replace $\widetilde{Y}_\infty$ with $e^{-\lambda t}Y_t$ at the cost of one more $\delta$ to arrive at 
  \[-\delta\geq e^{-\lambda t}X_t-e^{-\lambda t}Y_t\]
for $t\geq T$. In particular, for almost all sample paths such that $\liminf_{t\to\infty}e^{-\lambda t}(X_t-Y_t)<0$ $\int_0^\infty e^{-\lambda t}|X_t-Y_t|\,dt=\infty$, thus the theorem follows.


\end{proof}

\subsubsection{Representation of $W$}\label{app:W_representation}

The following is a useful representation of $W=\lim_{t\to\infty} e^{-\lambda_1 t}Z_1(t)$.

\begin{lemma}
  $W=\sum_{i\in \mathbb{N}} e^{-\lambda_1 T_i}Y_1^{(i)}$, where $Y_1^{(i)}$ are i.i.d. deflated exponentials where the mass at $0$ has probability $p_1$ and the rate of the exponential part is $q_1$.
\end{lemma}

\begin{proof}
  We have:

  \begin{align}
    e^{-\lambda_1 t}Z_1(t)=\sum_{T_i\leq t} e^{-\lambda_1 t}Z_1^{(i)}(t-T_i)=\sum_{i\in \mathbb{N}}\mathbf{1}_{T_i\leq t}e^{-\lambda_1 t}Z_1^{(i)}(t-T_i),\label{eq:W_limit}
  \end{align}
where $Z_1^{(i)}$ are i.i.d. branching processes of each type $1$ clone.  Once again, for the $i$ such that $T_i=\infty$, $Z_1^{(i)}(s)\equiv 0$ for all $s$. We will take the limit as $t\to\infty$ in \eqref{eq:W_limit}, and the conclusion will follow by passing the limit inside the summation.  We have pointwise limits: $\lim_{t\to\infty}\mathbf{1}_{T_i\leq t}e^{-\lambda_1 T_i}e^{-\lambda_1(t-T_i)}Z_1^{(i)}(t-T_i)=e^{-\lambda_1 T_i}Y_1^{(i)}$. Furthermore, we can bound Eq. \ref{eq:W_limit} by:

  \begin{align}
    \sum_{i\in \mathbb{N}}e^{-\lambda_1 T_i}\sup_{0\leq s< \infty}e^{-\lambda_1s}Z_1^{(i)}(s).\label{eq:W_summable}
  \end{align}

  To apply dominated convergence, it suffices to check the right hand side of Eq. \ref{eq:W_summable} has finite expectation. Now, $e^{-\lambda_1s}Z_1^{(i)}(s)$ is a martingale with respect to the filtration generated by $Z_1^{(i)}$. Noting that $e^{-\lambda_1 s}Z_1^{(i)}(s)$ is cadlag and uniformly integrable, we obtain by Doob's maximal inequality that:
  \begin{align*}
      \mathbb{E}[\sup_{0\leq s\leq t}e^{-\lambda_1 s}Z_1^{(i)}(s)]\leq C\mathbb{E}[e^{-\lambda_1 s}Z_1^{(i)}(s)\log(e^{-\lambda_1 s}Z_1^{(i)}(s))+1]
  \end{align*}
  For all $t\geq 0$ and a constant $C$ not depending on $t$ ($x\log x$ evaluated at $0$ is taken to be $0$). Taking $t\to\infty$, by monotone convergence, the left hand side converges to $\mathbb{E}[\sup_{0\leq s<\infty}e^{-\lambda_1 s}Z_1^{(i)}(s)]$. It is further known that $\sup_{0\leq s<\infty}\mathbb{E}[(e^{-\lambda_1 s}Z_1^{(i)}(s))^2]<\infty$ so $\sup_{0\leq s<\infty}\mathbb{E}\left[\left(e^{-\lambda_1 s}Z_1^{(i)}(a)\log(e^{-\lambda_1 s}Z_1^{(i)}(s))\right)^{1+\varepsilon}\right]<\infty$ for any $0<\varepsilon<1$. Hence since the integrand of the upper bound converges to $Y\log Y+1$, by Vitali's convergence theorem we obtain:
  
  \[\mathbb{E}\left[\sup_{0\leq s< \infty}e^{-\lambda_1s}Z_1^{(i)}(s)\right]\leq C\mathbb{E}[Y\log Y+1]<\infty.\]

  Since $Z_1^{(i)}$ is independent from $T_i$, it then suffices to check that $\sum_{i\in \mathbb{N}}e^{-\lambda_1 T_i}$ has finite expectation, but this directly follows from Campbell's Theorem.

\end{proof}

\subsubsection{Approximating Hitting Times}

The following approximation lemmas will be useful in our analyses.

\begin{lemma}\label{lem:large_size_time}
  Let $X_t$ be a jointly measurable stochastic process such that there exists a $\lambda>0$ such that for $\omega\in A$, a probability $1$ event, $t\to X_t(\omega)$ is bounded over compact sets and $\lim_{t\to\infty}e^{-\lambda t}X_t(\omega)=\theta(\omega)>0$. Then defining $\tau_n=\inf\{t:X(t)\geq n\}$ and $t_n=\frac{1}{\lambda}\log(n/\theta)$, $|\tau_n-t_n|\to 0$.
\end{lemma}

\begin{proof}
  This follows the same proof as Proposition 1 of Gunnarrsson 2025 \cite{gunnarsson2025}
\end{proof}

\begin{lemma}\label{lem:point_process_times}
  Let $\Phi=\{T_k\}_{k=1}^\infty$ be a simple point process over $[0,\infty)$ with $\Phi([0,\infty))=\infty$. Denote $M_t=\Phi([0,t])$ for $t\geq 0$. Suppose there exists a $\lambda$ such that $\lim_{t\to\infty}e^{-\lambda t}M_t(\omega)=\theta(\omega)>0$. Then, defining $\widehat{T}_k$ by $k=\theta e^{\lambda \widehat{T}_k}$, $|\widehat{T}_k-T_k|\to 0$.
\end{lemma}

\begin{proof}
  Note that for all $\varepsilon>0$, there is $T(\omega)$ sufficiently large such that $t>T$ implies:

    \[(1-\varepsilon)\theta e^{\lambda t}\leq M_t\leq (1+\varepsilon)\theta e^{\lambda t}.\]
Fix $0<\varepsilon<1/2$ and define $\delta=-\frac{1}{\lambda}\log(1-\varepsilon)$ and $\delta_2=\frac{2}{\lambda}\log(1+\varepsilon)$. Choose $K$ sufficiently large such that $k>K$ implies $\widehat{T}_k+\delta>T$ and $\widehat{T}_k-\delta_2>T$. Note then that
\begin{align*}
    &M_{\widehat{T}_k+\delta}\geq(1-\varepsilon) \theta e^{\lambda \widehat{T}_k}e^{\lambda\delta}=k(1-\varepsilon)e^{\lambda \delta}= k,\\
    &M_{\widehat{T}_k-\delta_2}\leq (1+\varepsilon)\theta e^{\lambda \widehat{T}_k}e^{-\lambda \delta_2}=\frac{k}{1+\varepsilon}<k.
\end{align*}
Hence $\widehat{T}_k+\delta\geq T_k$ and $\widehat{T}_k-\delta_2\leq T_k$ and we obtain $|T_k-\widehat{T}_k|\leq \max\{\delta,\delta_2\}$. Hence, $\lim_{k\to\infty} |T_k-\widehat{T}_k|=0$
\end{proof}

\subsection{Exact Representation of Moments of Fixed-Time SFS}\label{app:SFS_mean}
Here we prove Proposition \ref{prop:SFS_mean}:
\SFSmean*
\begin{proof}
  We first establish the results \eqref{eq:mean_SFS} and \eqref{eq:mean_SFS_random_fitness}. Let $\mathcal{M}_{[s,s+dt]}$ denote the event of having precisely one event - being mutation - in the specified interval and $\mathcal{B}_j(s,t)$ the event that the clone created in $(s,s+dt]$ grew to size $j$ at time $t$; with the assumption that $\mathcal{B}_j(s,t)=\emptyset$ if no mutations are created in the interval $(s,s+dt]$. Then, the contribution of mutations arising within $(s,s+dt]$ to $E[S_j(t)]$ is:
  \begin{align*}
    \mathbb{P}(\mathcal{M}_{[s,s+dt]}\cap \mathcal{B}_j(s,t))+o(dt)&=\mathbb{P}(\mathcal{M}_{[s,s+dt]})\mathbb{P}(\mathcal{B}_j(s,t)|\mathcal{M}_{(s,s+dt]})+o(dt)\\
    &=(av+\mu)\mathbb{E}[Z_0(s)]P_j(t-s)\,dt+o(dt)
  \end{align*}
  where $o(dt)$ accounts for the possibility of multiple events occurring and that $P_j(t-s)$ ought be $P_j(t-s')$ for $s'\in (s,s+dt]$. Summing such contributions recovers Eq. \ref{eq:mean_SFS}. The formula for $\mathbb{E}[S_j^B(t)]$ analogously follows by conditioning $\mathbb{P}(\mathcal{B}_j(s,t)|\mathcal{M}_{[s,s+dt]})$ on the selection amount, distributed as $B$, of that mutant clone.

Next we consider the second moment results \eqref{eq:choose_SFS} and \eqref{eq:choose_SFS_random_fitness}.
Interpret $S_j(t)\choose 2$ as the number of distinct pairs of driver mutations (without considering order) that reach size $j$ at time $t$.
For $s<r$, the contribution to $S_j\choose 2$ from mutant pairs created in $[s,s+dt]$ and $[r,r+dt]$ is:

    \[\mathbb{P}(\mathcal{M}_{[s,s+dt]}\cap \mathcal{M}_{[r,r+dt]})P_j(t-s)P_j(t-r),\]
up to $o(dt^2)$ factor. Now, partition $\mathcal{M}_{[s,s+dt]}=\mathcal{M}_{[s,s+dt]}^b\sqcup \mathcal{M}_{[s,s+dt]}^f$, with the superscript $b,f$ denoting mutation occurring in a replication-dependent and -independent manner, respectively. Then:

    \begin{align*}
        \mathbb{P}(\mathcal{M}_{[s,s+dt]}\cap \mathcal{M}_{[r,r+dt]}|\mathcal{F}_s)&=\mathbb{P}(\mathcal{M}_{[r,r+dt]}|\mathcal{F}_s,\mathcal{M}_{[s,s+dt]}^b)\mathbb{P}(\mathcal{M}_{[s,s+dt]}^b|\mathcal{F}_s)\\
        &+\mathbb{P}(\mathcal{M}_{[r,r+dt]}|\mathcal{F}_s,\mathcal{M}_{[s,s+dt]}^f)\mathbb{P}(\mathcal{M}_{[s,s+dt]}^f|\mathcal{F}_s)\\
        &=((av+\mu)e^{(r-s)\lambda_0}Z_0(s)\,dt)avZ_0(s)\,dt+((av+\mu)e^{(r-s)\lambda_0}(Z_0(s)-1)\,dt)\mu Z_0(s)\,dt,
    \end{align*}
where equalities are up to $o(dt^2)$ factors. The last line of the previous display follows since given $\mathcal{F}_s$ and $\mathcal{M}^f_{[s,s+dt]}$, the type $0$ population at time $s+dt$ is $Z_0(s)-1$ (and given $\mathcal{F}_s$ and $\mathcal{M}^b_{[s,s+dt]}$, the type $0$ population remains $Z_0(s)$ at time $[s,s+dt]$). Taking expectation and using that $\mathbb{E}[Z_0(s)Z_0(r)]=e^{\lambda_0(r-s)}\mathbb{E}[Z_0(s)^2]$ and $\mathbb{E}[Z_0(r)]=e^{(r-s)\lambda_0}\mathbb{E}[Z_0(s)]$, the infinitesimal contribution is thus:

    \[\mathbb{P}(\mathcal{M}_{[s,s+dt]}\cap \mathcal{M}_{[r,r+dt]})P_j(t-s)P_j(t-r)=\left((av+\mu)^2\mathbb{E}[Z_0(r)Z_0(s)]-\mu(av+\mu)\mathbb{E}[Z_0(r)]\right)P_j(t-s)P_j(t-r)(dt)^2.\]
We recover Eq. \ref{eq:choose_SFS} by summing over $0\leq s<r\leq t$. The formula for the random fitness case analogously follows by conditioning the probability of the clones growing to size $j$ on the selection amount $B_1,B_2$ of those mutant clones.
\end{proof}

\subsection{Large Time Scaling of Mean of Fixed-Time SFS}\label{app:SFS_mean_asymptotics}
\SFSmeanasymptotics*
\begin{proof}
  Combine Propositions \ref{prop:exact_time_asymptotic_mean_SFS}, \ref{prop:time_rate_convergence_SFS_mean}, and  \ref{prop:mean_SFS_random_large_freq} below.
\end{proof}
We break up the statement into three Propositions and prove each one. It also depends on the form of $P_j(s)$ (for $j\geq 1$), which is (see \cite{durrettbranching2015}):
\begin{align}
    P_j(t)=\left(1-\frac{de^{\lambda_1 t}-d}{(a+b)e^{\lambda_1 t}-d}\right)\left(1-\frac{(a+b)e^{\lambda_1 t}-(a+b)}{(a+b)e^{\lambda_1 t}-d}\right)\left(\frac{(a+b)e^{\lambda_1 t}-(a+b)}{(a+b)e^{\lambda_1 t}-d}\right)^{j-1}\label{eq:P_j}
\end{align}
\begin{proposition}
  The scaled SFS has the following exact form:
  
  \[e^{-\lambda_0t}\mathbb{E}[S_j(t)]=(av+\mu)\int_0^t e^{-\lambda_0 s}P_j(s)\,ds=(av+\mu)\frac{\Gamma(j)\Gamma(\lambda_0/\lambda_1+1)}{\Gamma(j+\lambda_0/\lambda_1+1)}\sum_{n=0}^\infty \frac{(\lambda_0/\lambda_1)_n(j)_n}{(j+1+\lambda_0/\lambda_1)_nn!}p_1^n\]

  And, the following frequency asymptotic:

  \[\lim_{t\to\infty}e^{-\lambda_0t}\mathbb{E}[S_j(t)]=(av+\mu)\frac{q_1\Gamma(\lambda_0/\lambda_1+1)}{\lambda_1j^{\lambda_0/\lambda_1+1}}(1-p_1)^{-\lambda_0/\lambda_1}+\mathcal{O}\left(\frac{1}{j^{\lambda_0/\lambda_1+2}}\right)\]
  \label{prop:exact_time_asymptotic_mean_SFS}
\end{proposition}
\begin{proof}
  From Proposition \ref{prop:SFS_mean}:
  \begin{align}
    e^{-\lambda_0t}\mathbb{E}[S_j(t)]= e^{-\lambda_0t}(av+\mu)\int_0^t\mathbb{E}[Z_0(s)]P_j(t-s)ds=(av+\mu)\int_0^t e^{-\lambda_0 s}P_j(s)ds,\label{eq:exact_mean_driver_sfs}
\end{align}
we can then easily see
\begin{align}
  \lim_{t\to\infty}e^{-\lambda_0t}\mathbb{E}[S_j(t)]=(av+\mu)\int_0^\infty e^{-\lambda_0 s}P_j(s)\,ds.\label{eq:asymp_mean}
\end{align}
Using the definition of $P_j(s)$ we can write the previous integral (Eq. \ref{eq:asymp_mean}) as
$$
\int_0^\infty e^{-\lambda_0 s}P_j(s)\,ds=q_1^2\int_0^\infty \frac{e^{-s(\lambda_0+\lambda_1)}}{(1-p_1e^{-\lambda_1s})^2}\left(\frac{1-e^{-\lambda_1s}}{1-p_1e^{-\lambda_1s}}\right)^{j-1}ds.
$$
Using the changes of variable $x=e^{-\lambda_1 s}$ and $y=\frac{1-x}{1-p_1x}$, we can rewrite the previous display as
\begin{align}
&\frac{q_1^2}{\lambda_1}\int_0^1\frac{x^{\lambda_0/\lambda_1}}{(1-p_1x)^2}\left(\frac{1-x}{1-p_1x}\right)^{j-1}dx=\frac{q_1}{\lambda_1}\int_0^1\left(\frac{1-y}{1-p_1y}\right)^{\lambda_0/\lambda_1}y^{j-1}dy.\label{eq:SFS_mean_integral}
\end{align}
One way of analyzing this integral is to relate it to a hyper-geometric series. Noticing that $(1-y)^{-a}=\sum_{n=0}^\infty \frac{(a)_{n}}{n!}y^n$, where for positive number $(a)_{n}=\Gamma(a+n)/\Gamma(a)$ denotes the rising factorial, we can rewrite the integral on the right-hand-side of the previous display as:
\begin{align}
    \label{eq:ExpectedSFS_Series}
\frac{\Gamma(j)\Gamma(\lambda_0/\lambda_1+1)}{\Gamma(j+\lambda_0/\lambda_1+1)}\sum_{n=0}^\infty \frac{(\lambda_0/\lambda_1)_{n}(j)_{n}}{(j+1+\lambda_0/\lambda_1)_{n}n!}p_1^n.
\end{align}
By Watson's Lemma, noting the cofactor to the sum is the beta function ($B(j,\lambda_0/\lambda_1+1)$):

\[\frac{\Gamma(j)\Gamma(\lambda_0/\lambda_1+1)}{\Gamma(j+\lambda_0/\lambda_1+1)}=\frac{\Gamma(\lambda_0/\lambda_1+1)}{j^{\lambda_0/\lambda_1+1}}+\mathcal{O}\left(\frac{1}{j^{\lambda_0/\lambda_1+2}}\right).\]
Since the rising factorial function is increasing we know that
$$
\frac{(\lambda_0/\lambda_1)_{n}(j)_{n}}{(j+1+\lambda_0/\lambda_1)_{n}n!}p_1^n\leq \frac{(\lambda_0/\lambda_1)_n}{n!}p_1^n,
$$
and furthermore basic asymptotics of the Gamma function give us that as $j\to\infty$ the summands of \eqref{eq:ExpectedSFS_Series} converge to $\frac{(\lambda_0/\lambda_1)_n}{n!}p_1^n.$
Thus by dominated convergence we can conclude
\[\int_0^1 \left(\frac{1-y}{1-p_1y}\right)^{\lambda_0/\lambda_1}y^{j-1}\,dy\sim \frac{\Gamma(\lambda_0/\lambda_1+1)}{j^{\lambda_0/\lambda_1+1}}\sum_{n=0}^\infty \frac{(\lambda_0/\lambda_1)_n}{n!}(p_1)^n=\frac{\Gamma(\lambda_0/\lambda_1+1)}{j^{\lambda_0/\lambda_1+1}}(1-p_1)^{-\lambda_0/\lambda_1}.\]

Now, to get a rate of convergence in terms of $j$ as $j\to\infty$ on:

\[\sum_{n=0}^\infty \frac{(\lambda_0/\lambda_1)_n(j)_n}{(j+1+\lambda_0/\lambda_1)_nn!}p_1^n\]
it suffices to examine how quickly

\[\sum_{n=0}^\infty \frac{(\lambda_0/\lambda_1)_n}{n!}p_1^n\left(1-\frac{(j)_n}{(j+1+\lambda_0/\lambda_1)_n}\right)\to 0\]
as $j\to\infty$. Notice that:


\begin{align*}
  j\left(1-\frac{(j)_n}{(j+1+\lambda_0/\lambda_1)_n}\right)&=j\cdot \left(\frac{(j+1+\lambda_0/\lambda_1)_n-(j)_n}{(j+1+\lambda_0/\lambda_1)_n}\right)\\
  &\overset{j\to\infty}{\sim} j\left(\frac{(j+1+\lambda_0/\lambda_1)^n-j^n}{(j+1+\lambda_0/\lambda_1)^n}\right)\\
  &\overset{j\to\infty}{\sim} j\left((1+\frac{1+\lambda_0/\lambda_1}{j})^n-1\right)\\
  &= j\left(n\frac{1+\lambda_0/\lambda_1}{j}+\mathcal{O}(j^{-2})\right)\\
  &\overset{j\to\infty}{\sim} n(1+\lambda_0/\lambda_1)
\end{align*}
Notably then $f_j(n)=j\frac{(\lambda_0/\lambda_1)_n}{n!}p_1^n\left(1-\frac{(j)_n}{(j+1+\lambda_0/\lambda_1)_n}\right)$ has a pointwise limit as $j\to\infty$. Note that:

\[|f_j(n)|=\left|\frac{(\lambda_0/\lambda_1)_n}{n!}p_1^n\right|\cdot j\cdot \left|\left(1-\prod_{k=0}^{n-1} \frac{(j+k)}{(j+1+\lambda_0/\lambda_1+k)}\right)\right|\]
Using the Weierstrass product inequality:

\begin{align*}
  &=|\frac{(\lambda_0/\lambda_1)_n}{n!}p_1^n|\cdot j\cdot |\left(1-\prod_{k=0}^{n-1} \left(1-\frac{(1+\lambda_0/\lambda_1)}{(j+1+\lambda_0/\lambda_1+k)}\right)\right)|\\
  &\leq |\frac{(\lambda_0/\lambda_1)_n}{n!}p_1^n|\cdot j\cdot \sum_{k=0}^{n-1}\frac{(1+\lambda_0/\lambda_1)}{(j+1+\lambda_0/\lambda_1+k)}\\
  &\leq |\frac{(\lambda_0/\lambda_1)_n}{n!}p_1^n|\cdot n(1+\lambda_0/\lambda_1)
\end{align*}
which is summable against $n$ (and upper bounds $f_j$ for all $j$), hence we can use dominated convergence to interchange $\lim_{j\to\infty}$ and the sum in $\lim_{j\to\infty}\sum_{n=0}^\infty j\frac{(\lambda_0/\lambda_1)_n}{n!}p_1^n\left(1-\frac{(j)_n}{(j+1+\lambda_0/\lambda_1)_n}\right)$. We further notice that:

\[\lim_{j\to\infty}\sum_{n=0}^\infty j\frac{(\lambda_0/\lambda_1)_n}{n!}p_1^n\left(1-\frac{(j)_n}{(j+1+\lambda_0/\lambda_1)_n}\right)=\sum_{n=0}^\infty \frac{(\lambda_0/\lambda_1)_n}{n!}p_1^nn(1+\lambda_0/\lambda_1)=L\in (0,\infty)\]

This the implies that the original sum of interest

\[\sum_{n=0}^\infty \frac{(\lambda_0/\lambda_1)_n(j)_n}{(j+1+\lambda_0/\lambda_1)_n(n)!}p_1^n\]

converges at rate $\mathcal{O}(j^{-1})$ to its limit as $j\to\infty$. We thus finally obtain:

\begin{align*}
  \lim_{t\to\infty}e^{-\lambda_0 t}\mathbb{E}[S_j(t)]&=(av+\mu)\frac{q_1}{\lambda_1}\int_0^1 \left(\frac{1-y}{1-p_1y}\right)^{\lambda_0/\lambda_1}y^{j-1}\,dy\\
  &=(av+\mu)\frac{q_1\Gamma(\lambda_0/\lambda_1+1)}{\lambda_1j^{\lambda_0/\lambda_1+1}}(1-p_1)^{-\lambda_0/\lambda_1}+\mathcal{O}\left(\frac{1}{j^{\lambda_0/\lambda_1+2}}\right)
\end{align*}
\end{proof}
\begin{proposition}
  \[\left(\lim_{t\to\infty}e^{-\lambda_0 t}\mathbb{E}[S_j(t)]\right)-e^{-\lambda_0s}\mathbb{E}[S_j(s)]\overset{s\to\infty}{\sim}(av+\mu)\frac{q_1}{\lambda_0/\lambda_1+1}e^{-(\lambda_0+\lambda_1)s}\]
  \label{prop:time_rate_convergence_SFS_mean}
\end{proposition}
\begin{proof}
  Following our substitutions in the above proof, the integral limits under the substitution $x=e^{-\lambda_1 s}$ goes from $0,t$ to $e^{-\lambda_1t},1$. Then under the substitution $y=\frac{1-x}{1-p_1x}$, the limits become $0,\frac{1-e^{-\lambda_1 t}}{1-p_1 e^{-\lambda_1 t}}$. Let $r(t)=1-\frac{1-e^{-\lambda_1 t}}{1-p_1 e^{-\lambda_1 t}}$. Then:

  \[\left(\lim_{t\to\infty}e^{-\lambda_0 t}\mathbb{E}[S_j(t)]\right)-e^{-\lambda_0s}\mathbb{E}[S_j(s)]=(av+\mu)\int_{1-r(s)}^1 \left(\frac{1-y}{1-p_1y}\right)^{\lambda_0/\lambda_1}y^{j-1}\,dy.\]
For any $\varepsilon>0$ (sufficiently small so that $\frac{1}{(1-p_1)^{\lambda_0/\lambda_1}}>\varepsilon$), we have $s$ sufficiently large such that:

  \[\left(\frac{1}{(1-p_1)^{\lambda_0/\lambda_1}}-\varepsilon\right)\int_{1-r(s)}^1 (1-y)^{\lambda_0/\lambda_1}\,dy\leq \int_{1-r(s)}^1 \left(\frac{1-y}{1-p_1y}\right)^{\lambda_0/\lambda_1}y^{j-1}\,dy \leq \frac{1}{(1-p_1)^{\lambda_0/\lambda_1}}\int_{1-r(s)}^1 (1-y)^{\lambda_0/\lambda_1}\,dy.\]
Since $\int_{1-r(s)}^1 (1-y)^{\lambda_0/\lambda_1}\,dy=\frac{\lambda_1}{\lambda_0+\lambda_1}r(s)^{\lambda_0/\lambda_1+1}$ and $\varepsilon$ is arbitrary we can conclude that as $s\to\infty$

  \[\int_{1-r(s)}^1 \left(\frac{1-y}{1-p_1y}\right)^{\lambda_0/\lambda_1}y^{j-1}\,dy\sim \frac{1}{(1-p_1)^{\lambda_0/\lambda_1}}\frac{\lambda_1}{\lambda_0+\lambda_1}r(s)^{\lambda_0/\lambda_1+1}\sim (1-p_1)\frac{\lambda_1}{\lambda_1+\lambda_0}e^{-(\lambda_0+\lambda_1)s}.\]
\end{proof}
\begin{proposition}
  Under the random fitness advance model with Assumption \ref{assumption:random}:
  \[\lim_{t\to\infty}e^{-\lambda_0t}\mathbb{E}[S_j^B(t)]\overset{j\to\infty}{\sim} \frac{C}{j^{\frac{\lambda_0}{\lambda_1(b_{\max})}+1}\log j}\]
  $C$ is a function of the parameters of the model which can be explicitly determined.\label{prop:mean_SFS_random_large_freq}
\end{proposition}
\begin{proof}
  Let $\widetilde{C}(b)=(av+\mu)\frac{q_1(b)\Gamma(\lambda_0/\lambda_1(b)+1)}{\lambda_1(b)}(1-p_1(b))^{-\lambda_0/\lambda_1(b)}h(b)$. From Proposition \ref{prop:SFS_mean_asymptotics} and Fubini-Tonelli:
  \begin{align}
      \lim_{t\to\infty}e^{-\lambda_0 t}\mathbb{E}[S_j^{B}(t)]&=\mathbb{E}\left[(av+\mu)\frac{q_1(B)\Gamma(\lambda_0/\lambda_1(B)+1)}{\lambda_1(B)j^{\lambda_0/\lambda_1(B)+1}}(1-p_1(B))^{-\lambda_0/\lambda_1(B)}+\mathcal{O}\left(\frac{1}{j^{\lambda_0/\lambda_1(B)+2}}\right)\right]\nonumber\\
      &\overset{j\to\infty}{\sim} \int_0^{b_{\text{max}}}\widetilde{C}(b)e^{-\left(\lambda_0/\lambda_1(b)+1\right)\log j}\,db.\label{eq:large_time_mean_random_goal}
  \end{align}
 Therefore, we examine the rate at which \ref{eq:large_time_mean_random_goal} goes to $0$ as $j\to\infty$. Let $\phi(b)=\lambda_0/\lambda_1(b)+1$ and perform the change of variables $\ell=\phi(b)-\phi(b_{\max})$. Let $C(\ell):=\frac{\tilde{C}(b(\ell))}{|\phi'(b(\ell))|}$ and let $C:=C(0)$, where $C$ is as in the proposition statement. Then:

\[\int_0^{b_{\text{max}}}\tilde{C}(b)e^{-\left(\lambda_0/\lambda_1(b)+1\right)\log j}\,db=e^{-\phi(b_{\max})\log j}\int_0^{\ell_{\text{max}}}C(\ell)e^{-\ell\log j}\,d\ell.\]
$C(\ell)$ remains continuous at $0$ with $C(0)>0$ and is bounded over $[0,\ell_{\max}]$ by quantity $M$ from Assumption \ref{assumption:random}. Fix $\varepsilon>0$ and choose $\delta>0$ such that the $\sup_{-\delta<y<\delta}|C(0)-C(y)|<\varepsilon$ (note $\delta$ is independent of $j$).
Then:

\[(C(0)-\varepsilon)(1-e^{-\delta\log j})\leq (\log j)\int_0^\delta C(\ell)e^{-\ell \log j}\,d\ell\leq (C(0)+\varepsilon)(1-e^{-\delta\log j}).\]
Also notice that $(\log j)\int_\delta^{\ell_{\max}} C(\ell)e^{-\ell \log j}\,d\ell\to 0$ as $j\to\infty$ for $\delta>0$ since $C$ is bounded on the interval $[0,\ell_{max}].$
We thus have that
\[C(0)-\varepsilon\leq \liminf_{j\to\infty}(\log j)\int_0^{\ell_{\max}} C(\ell)e^{-\ell \log j}\,d\ell \leq \limsup_{j\to\infty}(\log j)\int_0^{\ell_{\max}} C(\ell)e^{-\ell \log j}\,d\ell\leq C(0)+\varepsilon\]
and since $\varepsilon$ is arbitrary, we conclude
\[e^{-\phi(b_{\max})\log j}\int_0^{\ell_{\text{max}}}C(\ell)e^{-\ell\log j}\,d\ell\sim \frac{C(0)}{j^{\phi(b_{\max})}\log j}.\]
\end{proof}
\subsection{SFS first-order convergence results}\label{app:SFS_convergence}

We prove Theorem \ref{prop:convergence_results} and Proposition \ref{prop:fraction_as} by proving several lemmas. We use the same overall proof structure as in Gunnarsson et al. \cite{gunnarsson2025} where the SFS was decomposed into a sum of monotone processes. However, our $L^2$ approximation scheme is simpler.
\subsubsection{Setup}
First, we introduce the notation as in \cite{gunnarsson2025}. Define $\tau_{j,-}^i(0)=T_i$ (the initiation time of the $i$-th driver clone produced) and recursively for $k\geq 1$, define
\begin{align*}
  &\tau_{j,+}^i(k)=\inf\{s>\tau_{j,-}^i(k-1):Z_1^{(i)}(s-T_i)=j\}\\
  &\tau_{j,-}^i(k)=\inf\{s>\tau_{j,+}^i(k):Z_1^{(i)}(s-T_i)\neq j\},
\end{align*}
with the convention $\inf\emptyset=\infty$. Then $\tau_{j,+}^i(k)$ is the time that the $i$th mutant clone hits size $j$ for the $k$th time, and $\tau_{j,-}^i(k)$ is the time that the clone exits out of size $j$ for the $k$th time. Then define:
\begin{align*}
  I_{j,\pm}^{(i)}(t)=\sum_{k\in \mathbb{N}}\mathbf{1}_{\tau_{j,\pm}^i(k)\leq t}
\end{align*}
which are the number of times the $i$th clone has entered into or exited size $j$ by time $t$. Finally, we define:
\begin{align}
  S_{j,\pm}^{k}(t)=\sum_{i\in\mathbb{N}} \mathbf{1}_{I^{(i)}_{j,\pm}(t)\geq k}\label{eq:def_S_j_pm}
\end{align}
which are the number of mutant clones that have reached size $j$ at least $k$ times (or exited size $j$ at least $k$ times) by time $t$. Note that $t\to S_{j,\pm}^{k}(t)$ is monotone increasing and:
\begin{align}
  S_j(t)=\sum_{k\in \mathbb{N}} S_{j,+}^{k}(t)-\sum_{k\in \mathbb{N}} S_{j,-}^{k}(t).\label{eq:S_j_decomposition}
\end{align}
We also define:
\[P^k_{j,\pm}(t)=\mathbb{P}\left(I_{j,\pm}^{(1)}(t)\geq k|Z_1^{(1)}(0)=1,Z_0(0)=0\right)\]
which is the probability that a type-1 clone has entered into (or exited from) size $j$ at least $k$ times by time $t$. Note that
\begin{align*}
\mathbf{1}_{\{Z_1(t)=j\}}
= \sum_{k\ge1}\big(\mathbf{1}_{\{I^{(1)}_{j,+}(t)\ge k\}}
-
\mathbf{1}_{\{I^{(1)}_{j,-}(t)\ge k\}}\big),
\end{align*}
since being at size $j$ at time $t$ means the clone has entered size $j$ exactly one more time than it has exited it; taking expectations gives 
\begin{align}
    P_{j}(t)=\sum_{k\in \mathbb{N}}\left(P_{j,+}^k(t)-P_{j,-}^k(t)\right).\label{eq:P_j_decomposition}
\end{align}

We also use the following approximating processes:

\begin{table}[H]
  \centering
\begin{tabular}{|c|p{12cm}|}
\hline
\textbf{Notation} & \textbf{Definition} \\
\hline
$\overline{S}_{j,\pm}^{k}(t)$ & $(av+\mu)\int_0^t Z_0(s)P_{j,\pm}^k(t-s)\,ds$\\
$\widehat{S}_{j,\pm}^{k}(t)$ & $(av+\mu)\int_0^t Ye^{\lambda_0 s}P_{j,\pm}^k(t-s)\,ds$\\
$\overline{S}_j(t)$ & $(av+\mu)\int_0^t Z_0(s)P_{j}(t-s)\,ds$\\
$\widehat{S}_j(t)$ & $(av+\mu)\int_0^t Ye^{\lambda_0 s}P_{j}(t-s)\,ds$\\
\hline
\end{tabular}
\caption{SFS notation}\label{table:SFS_notation}
\end{table}

The following uniform bound will be useful in what follows.
\begin{lemma}\label{lem:uniform_bound_kth_entrance}
  There exists $\theta \in (0,1)$, independent of $j,k,b,t$, such that $P_{j,\pm}^k(t)\lesssim \theta^k$ as $k\to\infty$.\label{lem:P_j_pm_summable}
\end{lemma}
\begin{proof}
  Let $\mathcal{Z}(n)$ denote a discrete time-embedding of the continuous-time linear birth-death process associated with a type-1 clone. Let $\mathcal{Z}(0)=j$. If the first step is a birth, then $\mathbb{P}(\text{exists some time }n\geq 2\text{ with }\mathcal{Z}(n)=j|\mathcal{Z}(1)=j+1)=\frac{d}{a+b}$, since this the extinction probability. Thus 
\[  
\mathbb{P}(\text{exists some time }n\geq 1\text{ with }\mathcal{Z}(n)=j|\mathcal{Z}(0)=j)\leq \frac{d}{a+b}\frac{a+b}{a+b+d}+1\cdot \frac{d}{a+b+d}=2\frac{d}{a+b+d}\leq \frac{2d}{a+d}<1,
\]
and we can thus take $\theta=\frac{2d}{a+d}$. Notice that $P_{j,-}^k(t)\leq P_{j,+}^k(t)$ since $\tau_{j,-}^{(1)}(k)\geq \tau_{j,+}^{(1)}(k)$ for all $k\geq 1$, which implies $I_{j,+}^{(1)}(t)\geq I_{j,-}^{(1)}(t)$ for all $t\geq 0$. The result then follows by the strong Markov property.
\end{proof}

\begin{lemma}\label{lem:CenteredMuts_Martingale}
  For $t>0$ define $R_t=M_t-(av+\mu)\int_0^t Z_0(s)\,ds$ then $\{R_t;t\geq 0\}$ is a 0-mean $\mathcal{B}_t$-martingale. 
\end{lemma}

\begin{proof}
By forming a partition $\Pi_t$ over $[0,t]$ of intervals of length $dt$, let $0<s<t$ and over one such increment $\mathbb{E}[M_s-M_{s-dt}]=(av+\mu)\mathbb{E}[Z_0(s)]dt+o(dt^2)$. Summing over all contributions $\mathbb{E}[M_t]=(av+\mu)\mathbb{E}\left[\int_0^t Z_0(s)\,ds\right].$

  Now fix $t>s\geq 0$, notice:

  \[\mathbb{E}\left[M_t-(av+\mu)\int_0^tZ_0(r)\,dr|\mathcal{B}_s \right]=\mathbb{E}\left[(M_t-M_s)-(av+\mu)\int_s^tZ_0(r)\,dr|\mathcal{B}_s \right]+M_s-(av+\mu)\int_0^sZ_0(r)\,dr\]

  However, using the strong Markov property and independence of lineages:
  \[\mathbb{E}\left[(M_t-M_s)-(av+\mu)\int_s^tZ_0(r)\,dr|\mathcal{B}_s \right]=Z_0(s)\mathbb{E}[M_{t-s}-(av+\mu)\int_0^{t-s}Z_0(r)\,ds]=0\]
\end{proof}

We next show an $L^2$ approximation result for the monotonic processes $\{S_{j,\pm}^k(t);t\geq 0\}$:

\begin{proposition}\label{prop:L2_driver}
  There exists $A<\infty$ such that $\mathbb{E}\left[(S_{j,\pm}^{k}(t)-\overline{S}_{j,\pm}^{k}(t))^2\right]\leq A\theta^k e^{\lambda_0 t}$
\end{proposition}

\begin{proof}
Let $C_t=\mathbb{E}[S_{j,\pm}^k(t)|\mathcal{B}_t]$ and then note that

\begin{align}\label{eq:2ndMomentDecomposition}
    &\mathbb{E}\left[\left(S_{j,\pm}^{k}(t)-\overline{S}_{j,\pm}^{k}(t)\right)^2\right]\nonumber\\
    &=\mathbb{E}\left[\left(S_{j,\pm}^{k}(t)-C_t+C_t-\overline{S}_{j,\pm}^{k}(t)\right)^2\right]\nonumber\\
    &=\mathbb{E}\left[\left(S_{j,\pm}^{k}(t)-C_t\right)^2\right]+2\mathbb{E}\left[\left(S_{j,\pm}^{k}(t)-C_t\right)\left(C_t-\overline{S}_{j,\pm}^{k}(t)\right)\right]+\mathbb{E}\left[\left(C_t-\overline{S}_{j,\pm}^{k}(t)\right)^2\right].
\end{align}

Note that $C_t,\overline{S}_{j,\pm}^{k}(t)$ are both $\mathcal{B}_t$-measurable. Thus using towering the middle term vanishes:

\[\mathbb{E}\left[\left(S_{j,\pm}^{k}(t)-C_t\right)\left(C_t-\overline{S}_{j,\pm}^{k}(t)\right)|\mathcal{B}_t\right]=0.\]
Now note that the first term in \eqref{eq:2ndMomentDecomposition} is the mean of the conditional variance: $\mathbb{E}\left[\mathbb{V}\left(S_{j,\pm}^{k}(t)|\mathcal{B}_t\right)\right].$
Recall our definition of $S_{j,\pm}^{k}(t)$ (Eq. \ref{eq:def_S_j_pm}) which is a sum of indicators of functions of $Z_1^{(i)}$ (which are i.i.d.) and $T_i$. Thus conditional on $\mathcal{B}_t$, the sum is over independent random indicators\footnote{Note that these indicators would not be independent if we wanted to prove analogous limit theorems for hitchiker/passenger SFS.}. Also the variance of the indicator, conditional on $\mathcal{B}_t$, is simply $P_{j,\pm}^{k}(t-T_i)(1-P_{j,\pm}^{k}(t-T_i))$ and we thus have

\begin{align}
  \mathbb{E}\left[\mathbb{V}\left(S_{j,\pm}^{k}(t)|\mathcal{B}_t\right)\right]&=\mathbb{E}\left[\sum_{T_i\leq t} P_{j,\pm}^{k}(t-T_i)(1-P_{j,\pm}^{k}(t-T_i))\right]\nonumber\\
  &=\int_0^t (av+\mu)\mathbb{E}[Z_0(s)]P_{j,\pm}^{k}(t-s)(1-P_{j,\pm}^{k}(t-s))\,ds\leq (A/2)\theta^k e^{\lambda_0 t}.\label{eq:L2_S_to_C}
\end{align}
With the second to last equality due to Campbell's theorem, and the inequality an application of Lemma \eqref{lem:uniform_bound_kth_entrance}.

For the third term in \eqref{eq:2ndMomentDecomposition}, observing that $C_t=\sum_{T_i\leq t}P_{j,\pm}^{k}(t-T_i)=\int_0^t P_{j,\pm}^k(t-s,1)\,dM_s$, and recalling the definition of $\overline{S}_{j,\pm}^{k}(t)$, we can can conclude that $C_t-\overline{S}_{j,\pm}^{k}(t)=\int_0^t P_{j,\pm}^{k}(t-s)\,dR_s$, where $R_t=M_t-(av+\mu)\int_0^t Z_0(s)\,ds$ is a $0$-mean martingale via Lemma \ref{lem:CenteredMuts_Martingale}. The quadratic variation of $R_t$, denoted $[R]_t$ is:
\begin{align*}
    [R]_t=\sum_{0\leq v\leq t}(\Delta R_v)^2=\sum_{0\leq v\leq t}(\Delta M_v)^2=\sum_{0\leq v\leq t}\Delta M_v=M_t.
\end{align*}
So by Ito's isometry and Campbell's theorem:
\begin{align}
  \mathbb{E}\left[(C_t-\overline{S}_{j,\pm}^{k}(t))^2\right]=\mathbb{E}\left[\int_0^t P_{j,\pm}^{k}(t-s)^2\,dM_s\right]=\int_0^t (av+\mu)e^{\lambda_0 s}P_{j,\pm}^{k}(t-s)^2\,ds\leq (A/2)\theta^k e^{\lambda_0t}\label{eq:L2_C_int_Z}.
\end{align}
Combining \ref{eq:L2_S_to_C} and \ref{eq:L2_C_int_Z} gives the result.
\end{proof}

We will need an estimate of the error between $Z_0$ and $Ye^{\lambda_0 s}$ in $L^2$. This can be calculated explicitely:
\begin{lemma}\label{lem:L2_distance}
    Let $Z(s)$ be a one-type continuous-time linear birth-death process with birth rate $\alpha$ and death rate $\beta$ with $\alpha>\beta$. Denote $\lambda=\alpha-\beta$. Let $Q_t=e^{-\lambda t}Z(t)$ and $Y=\lim_{t\to\infty} Q_t$, almost surely. Then:
    \[\mathbb{E}\left[\left(e^{-\lambda t}Z(t)-Y\right)^2\right]=\frac{\alpha+\beta}{\lambda}e^{-\lambda t}\]
\end{lemma}
\begin{proof}
    This can be computed using the Kolmogorov backwards equation and using various properties of branching processes. There is however an easier way. Note that the quadratic variation:
    \[[Q]_t=\sum_{T\in\Phi_{\text{event}}} e^{-2\lambda T},\]
    where $\Phi_{\text{event}}$ is the point process of event timings, having mean intensity $(\alpha+\beta)e^{\lambda t}\,dt$. $Q_t$ is a $\sigma(\{Z(s)\}_{0\leq s\leq t})$-martingale. Then, we can write $Y-e^{-\lambda t}Z(t)$ as a stochastic integral and use Ito's isometry followed by Campbell's theorem:
    \begin{align*}
        \mathbb{E}\left[\left(Y-e^{-\lambda t}Z(t)\right)^2\right]=\mathbb{E}\left[\left(\int_t^\infty\,dQ_s\right)^2\right]&=\mathbb{E}\left[\int_t^\infty d[Q]_s\right]\\
        &=(\alpha+\beta)\int_t^\infty e^{-\lambda s}\,ds\\
        &=\frac{\alpha+\beta}{\lambda}e^{-\lambda t}
    \end{align*}
\end{proof}

\begin{lemma}\label{lem:pm_approx_via_Y}
There exists a positive constant $A$ such that for all $t>0$
\[
\mathbb{E}\left[(\overline{S}_{j,\pm}^{k}(t)-\widehat{S}_{j,\pm}^{k}(t))^2\right]\leq At\theta^ke^{\lambda_0 t}.
\]
\end{lemma}

\begin{proof}
  By Cauchy-Schwarz:
  \[\left(\int_0^t (av+\mu)P_{j,\pm}^{k}(t-s)(Z_0(s)-Ye^{\lambda_0 s})\,ds\right)^2\leq (av+\mu)^2\int_0^t P_{j,\pm}^{k}(t-s)^2\,ds\int_0^t (Z_0(s)-Ye^{\lambda_0 s})^2\,ds\]
  Note the bounds:
  \[\int_0^t P_{j,\pm}^{k}(t-s)^2\,ds\leq C\theta^k t\]
  and:
  \[\mathbb{E}\left[\left(Z_0(s)-Ye^{\lambda_0 s}\right)^2\right]\lsimas{s}{\infty}e^{\lambda_0 s},\]
  from the prior lemma. The proof is complete combining the above.
\end{proof}

We can now combine our previous lemmas and propositions to establish the following a.s. convergence.

\begin{proposition}
  $\lim_{t\to\infty} e^{-\lambda_0 t}S_j(t)=(av+\mu)Y\int_0^\infty e^{-\lambda_0 s}P_j(s)\,ds$ on $\Omega_0^\infty$.\label{prop:almost_sure_conv}
\end{proposition}

\begin{proof}
  Using Proposition \ref{prop:L2_driver}, we have:

    \[\mathbb{E}\left|S_{j,\pm}^{k}(t)-\overline{S}_{j,\pm}^{k}(t)\right|\leq \sqrt{\mathbb{E}\left|S_{j,\pm}^{k}(t)-\overline{S}_{j,\pm}^{k}(t)\right|^2}\leq C_2\theta^{k/2}e^{(\lambda_0/2) t}.\]
From Lemma \ref{lem:pm_approx_via_Y} we obtain:

\[\mathbb{E}\left|\overline{S}_{j,\pm}^{k}(t)-\widehat{S}_{j,\pm}^{k}(t)\right|\leq C_3\sqrt{t}\theta^{k/2}e^{\lambda_0 t/2},\]
and using the triangle inequality:
\[\mathbb{E}\left|S_{j,\pm}^{k}(t)-\widehat{S}_{j,\pm}^{k}(t)\right|\leq C_4 \theta^{k/2}\sqrt{t}e^{(\lambda_0t)/2}.\]

Define $\widehat{S}_{j,\pm}(t):=\sum_{k\geq 1}\widehat{S}_{j,\pm}^{k}$, which is summable from Lemma \ref{lem:P_j_pm_summable}, and then via the triangle inequality:

\[\mathbb{E}\left|\sum_{k=1}^\infty S_{j,\pm}^{k}(t)-\widehat{S}_{j,\pm}(t)\right|\leq \sum_{k=1}^\infty \mathbb{E}\left|S_{j,+}^{k}(t)-\widehat{S}_{j,+}^{k}(t)\right|\lesssim\sqrt{t}e^{\lambda_0/2\cdot t}\]
Thus, $\int_0^\infty e^{-\lambda_0 t}\mathbb{E}\left|\sum_{k\geq 1}S_{j,\pm}^{k}(t)-\widehat{S}_{j,\pm}(t)\right|<\infty$. We get a.s. convergence of the sum by noting $e^{-\lambda_0 t}\widehat{S}_{j,\pm}(t)$ converges a.s. and by invoking Proposition \ref{prop:convergence_tool}. By writing $S_j(t)$ as in Eq. \ref{eq:S_j_decomposition} and using the identity Eq. \ref{eq:P_j_decomposition}, the statement is proven.
\end{proof}

We can also extend the above analysis to get a.s. convergence in the case where the selection amount is random:

\begin{corollary}
    Suppose that each clone has random fitness $B$, where $B$ has support within a compact interval $[0,b_{\text{max}}]$. Then:

    \[\lim_{t\to\infty}e^{-\lambda_0t}S_j^B(t)=(av+\mu)Y\int_0^\infty e^{-\lambda_0 s}\mathbb{E}\left[P_j(s|B)\right]\,ds\]
 almost surely.\label{cor:L2_as_random_fitness}
\end{corollary}

\begin{proof}
  Define analogs of $P_{j,\pm}^k$ and the approximations $\overline{S}_{j,\pm}^k,\widehat{S}_{j,\pm}^k$ by replacing $P_{j,\pm}^k$ with the expectation of the probability conditional on selection amount $B$:
  \begin{align*}
      P_{j,\pm}^{k,B}(t)=\mathbb{E}[\mathbb{P}(I_{j,\pm}^{(i)}(t)\geq k|Z_1^{(1)}(0)=1,Z_0(0)=0,B)],
  \end{align*}
  and replacing $P_{j,\pm}^k$ with $P_{j,\pm}^{k,B}$ in the definitions of $\overline{S}_{j,\pm}^k,\widehat{S}_{j,\pm}^k$. Remark still that:
  \[P_{j}^B(t)=\sum_{k\in \mathbb{N}}\left(P_{j,+}^{k,B}(t)-P_{j,-}^{k,B}(t)\right)\]
  by linearity of expectation. Importantly, Lemma \ref{lem:uniform_bound_kth_entrance} follows since $\theta$ is independent of the selection amount. Thus, performing the same analysis, the $L^2$ errors from Proposition \ref{prop:L2_driver} and Lemma \ref{lem:pm_approx_via_Y} are still of the form $\theta^k e^{\lambda_0 t}$, up to a constant multiple. Therefore, the analysis of Proposition \ref{prop:almost_sure_conv} will still show that $S_j^B$ converges almost surely, as desired.

\end{proof}

Then, the large detection-size behavior of $S_j$ follows by approximating $\tau_n$ by $ (1/\lambda_1)\log(n/W)$.

\begin{proposition}
    $\lim_{n\to\infty} \frac{1}{n^{\lambda_0/\lambda_1}}S_j(\tau_n)=(av+\mu)YW^{-\lambda_0/\lambda_1}\int_0^\infty e^{-\lambda_0 s}P_j(s)\,ds$ (conditional on $\Omega_0^\infty$)\label{prop:convergence_detection_size}
\end{proposition}

\begin{proof}
Since $\tau_n\to\infty$ a.s., we can now apply Proposition \ref{prop:almost_sure_conv} to conclude that condtional on $\Omega_\infty$:
  \begin{align}
    \lim_{n\to\infty} e^{-\lambda_0\tau_n}S_j(\tau_n)=(av+\mu)Y\int_0^\infty e^{-\lambda_0 s}P_j(s)\,ds\label{eq:hitting_times_large}.
  \end{align}
    Set $t_n=\frac{1}{\lambda_1}\log\left(\frac{n}{W}\right)$ and note that $|\tau_n-t_n|\to 0$ as $n\to\infty$ a.s. by Lemma \ref{lem:large_size_time}. This can then be combined with \eqref{eq:hitting_times_large} to conclude
    \begin{align}
        \lim_{n\to\infty} e^{-\lambda_0\tau_n}S_j(\tau_n)&=\lim_{n\to\infty} \left(e^{-\lambda_1(\tau_n-t_n)}e^{-\lambda_1t_n}\right)^{\lambda_0/\lambda_1}S_j(\tau_n)\nonumber\\
        &=\lim_{n\to\infty}\frac{W^{\lambda_0/\lambda_1}}{n^{\lambda_0/\lambda_1}}S_j(\tau_n).\label{eq:size_large}
    \end{align}
The desired result then follows by combining Eqs. \ref{eq:hitting_times_large} and \ref{eq:size_large}.
\end{proof}

Finally, Theorem \ref{prop:convergence_results} can easily be shown.

\ConvergenceResults*

\begin{proof}
  $L^2$ convergence follows by adapting Proposition \ref{prop:L2_driver} and Lemma \ref{lem:pm_approx_via_Y}. Specifically, Proposition \ref{prop:L2_driver} can be adapted to show $\mathbb{E}\left[\left(S_j(t)-\overline{S}_j(t)\right)^2\right]\lesssim e^{\lambda_0 t}$ by replacing $S_{j,\pm}^{k}$ with $S_j$ and $P_{j,\pm}^{k}$ with $P_j$. Similarly, adapting Lemma \ref{lem:pm_approx_via_Y}, we get $\mathbb{E}\left[(\overline{S}_j(t)-\widehat{S}_j(t))^2\right]\lesssim te^{\lambda_0 t}$. Thus $e^{-\lambda_0 t}S_j(t)$ converges in $L^2$.
  
  Almost sure convergence for $e^{-\lambda_0 t}S_j(t)$ and $n^{-\lambda_0/\lambda_1}S_j(\tau_n)$ were established in Propositions \ref{prop:almost_sure_conv} and \ref{prop:convergence_detection_size} respectively. $L^2$ and a.s. convergence the random fitness advance setting follows from Corollary \ref{cor:L2_as_random_fitness}. 
\end{proof}

Finally, we can easily prove convergence results for $S_{\geq j}(t)$, counting the number of mutations present in $\geq j$ cells at time $t$.

\begin{corollary}
  The following a.s. convergence result hold on $\Omega_0^\infty$:
  \begin{align*}
    &\lim_{t\to\infty}e^{-\lambda_0 t}S_{\geq j}(t)=(av+\mu)Y\int_0^\infty e^{-\lambda_0 s}\sum_{k=j}^\infty P_k(s)\,ds,\\
    &\lim_{n\to\infty}\frac{1}{n^{\lambda_0/\lambda_1}}S_{\geq j}(\tau_n)=(av+\mu)YW^{-\lambda_0/\lambda_1}\int_0^\infty e^{-\lambda_0 s}\sum_{k=j}^\infty P_k(s)\,ds.
  \end{align*}
  The series of equalities given by \ref{eq:normalized_SFS_limit} holds.
\end{corollary}

\begin{proof}
  Conduct the same analysis as to show convergence for $S_j$, but take $\tau_{j,+}^i(k),\tau_{j,-}^i(k)$ to be the times the $i$th mutant clone becomes \textit{at least} size $j$ $k$ times and the times the $i$th mutant clone exits from the region $\{j,j+1,\dots\}$ $k$ times.

  For the series of equalities in \ref{eq:normalized_SFS_limit}, the first three expressions are clearly equal by the almost sure limits. Therefore, it remains to check the last equality. From Proposition \ref{prop:SFS_mean_asymptotics}:
  \begin{align*}
      \lim_{t\to\infty}e^{-\lambda_0 t}\mathbb{E}[S_j(t)]&=(av+\mu)\frac{q_1}{\lambda_1}\frac{\Gamma(j)\Gamma(\lambda_0/\lambda_1+1)}{\Gamma(j+\lambda_0/\lambda_1+1)}\sum_{n=0}^\infty \frac{(\lambda_0/\lambda_1)_n(j)_n}{(j+1+\lambda_0/\lambda_1)_n n!}p_1^n\\
      &=(av+\mu)\frac{q_1}{\lambda_1}\frac{\Gamma(j)\Gamma(\lambda_0/\lambda_1+1)}{\Gamma(j+\lambda_0/\lambda_1+1)}{}_2F_1(\lambda_0/\lambda_1,j;j+\lambda_0/\lambda_1+1;p_1).
  \end{align*}
  By monotone convergence:
  \begin{align*}
      \lim_{t\to\infty}e^{-\lambda_0 t}\mathbb{E}[S_{\geq 1}(t)]=(av+\mu)\sum_{j\geq 1}\int_0^\infty P_j(s)e^{-\lambda_0 s}\,ds
  \end{align*}
  and by using \ref{eq:SFS_mean_integral}
  \begin{align*}
      &=(av+\mu)\frac{q_1}{\lambda_1}\int_0^1 (1-y)^{\lambda_0/\lambda_1-1}(1-p_1 y)^{-\lambda_0/\lambda_1}\,dy\\
      &=(av+\mu)\frac{q_1}{\lambda_0}{}_2F_1(\lambda_0/\lambda_1,1;\lambda_0/\lambda_1+1;p_1).
  \end{align*}
  Dividing both yields the last equality.
\end{proof}



\subsection{Sparsity of mutations at sizes varying in time}\label{app:SFSSparsity}
Here we prove Theorem \ref{thm:SFSSparsity} and then Corollary \ref{cor:SFSSparsity}.
\SFSSparsity*
\begin{proof}
  This is a combination of Corollary \ref{cor:sparsity_part1_2_3}, Lemmas \ref{lem:sparsity_part2} and \ref{lem:sparsity_part3}, and Propositions \ref{prop:limsup_j_k_plus} and \ref{prop:liminf_j_k_mimus}
\end{proof}
Throughout, assume $j(t)$ is non-decreasing. We prove this Proposition in various steps; firstly, we obtain a simpler expression of $\mathbb{E}[S_{j(t)}(t)]$:
\begin{lemma}\label{lem:changing_j_mean}
  Define:
  \begin{align}
      I_t(n)=\int_{ne^{-\lambda_1t}}^{n}\frac{\lambda_1u^{\frac{\lambda_0}{\lambda_1}}}{((a+b)-d(u/n))^2}\left(\frac{(a+b)n-(a+b)u}{(a+b)n-du}\right)^{n-1}\,du\label{eq:I_t_n_def}
  \end{align}
  Then for $j(t):\mathbb{R}_{\geq 0}\to\mathbb{N}$:
  \begin{align}
    \mathbb{E}[S_{j(t)}(t)]=(av+\mu)e^{\lambda_0 t}j(t)^{-\frac{\lambda_1+\lambda_0}{\lambda_1}}I_t(j(t))\label{eq:meanSFS_varying}
  \end{align}
\end{lemma}

Furthermore, $\sup_{n\in \mathbb{N},t\geq 0}I_t(n)$ is bounded; if $j(t)\ll e^{\lambda_1 t}$, then $I_t(j(t))$ converges to a constant if $j(t)$ is non-decreasing.
\begin{proof}
  By Proposition \ref{prop:SFS_mean}
    \[\mathbb{E}[S_{j(t)}(t)]=(av+\mu)e^{\lambda_0 t}\int_0^t e^{-\lambda_0 s}P_{j(t)}(s)\,ds.\]
    Explicitly write out $P_j(s)$ using \ref{eq:P_j} and use the change of variables $u=j(t)e^{-\lambda_1 s}$:
    \begin{align}
        &=\frac{(av+\mu)}{\lambda_1}e^{\lambda_0 t}\int_{j(t)e^{-\lambda_1 t}}^{j(t)}u^{\frac{\lambda_0}{\lambda_1}}j(t)^{-\frac{\lambda_0}{\lambda_1}}\frac{\lambda_1^2(j(t)/u)}{((a+b)(j(t)/u)-d)^2}\left(\frac{(a+b)(j(t)/u)-(a+b)}{(a+b)(j(t)/u)-d}\right)^{j(t)-1}\,\frac{du}{u}\nonumber\\
        &=(av+\mu)e^{\lambda_0 t}j(t)^{-\left(\frac{\lambda_1+\lambda_0}{\lambda_1}\right)}I_t(j(t))\nonumber
    \end{align}
    proving the equation \ref{eq:meanSFS_varying}. Denote $I_\infty(n)$ by the same equation as \ref{eq:I_t_n_def} but with the lower integration bound set to $0$. Remark that by dominated convergence:
    \begin{align}
        I_\infty(n)\overset{n\to\infty}{\to} \int_0^\infty \frac{\lambda_1}{(a+b)^2}u^{\frac{\lambda_0}{\lambda_1}}e^{-\frac{a+b-d}{a+b}u}\,du\label{eq:limit_I_inf}
    \end{align}
    which is strictly positive. Using that $t\to I_t(n)$ is increasing:
    \[\sup_{n\in \mathbb{N},t\geq 0}I_t(n)=\sup_{n\in \mathbb{N}}I_\infty(n)<\infty.\]

    Finally, under the condition that $j(t)\ll e^{\lambda_1 t}$ and $j(t)\to\infty$ as $t\to\infty$, we may write:
    \begin{align}
        I_t(j(t))=I_\infty(j(t))+(I_t(j(t))-I_\infty(j(t))),\label{eq:I_t_rewrite}
    \end{align}
    where the second term converges to $0$ by dominated convergence. So, $I_t(j(t))$ converges to the same limit as on the right-hand-side of \ref{eq:limit_I_inf}. If $j(t)\to N$, then examining \ref{eq:I_t_rewrite} once again, $I_t(j(t))$ converges to $I_\infty(N)$. \end{proof}

We first prove the easiest parts of the theorem in this corollary.
\begin{corollary}\label{cor:sparsity_part1_2_3}
The following hold:
\begin{itemize}
    \item $j(t)\gg\exp(t\lambda_0\lambda_1/(\lambda_0+\lambda_1))$ implies $S_{j(t)}(t)\to 0$ in $L^1$.
    \item $j(t)\sim C\exp(t\lambda_0\lambda_1/(\lambda_0+\lambda_1))$ implies:
    \[\lim_{t\to\infty}\mathbb{E}[S_{j(t)}(t)]=\frac{\lambda_1(av+\mu)}{(a+b)^2}C^{-\frac{\lambda_1+\lambda_0}{\lambda_1}}\left(\frac{a+b}{a+b-d}\right)^{\frac{\lambda_1+\lambda_0}{\lambda_1}}\Gamma(1+\lambda_0/\lambda_1).\]
    \item $j(t)\ll \exp(t\lambda_0\lambda_1/(\lambda_0+\lambda_1))$ implies
    \begin{align*}
        \lim_{t\to\infty}\mathbb{E}[e^{-\lambda_0 t}j(t)^{(\lambda_1+\lambda_0)/\lambda_1}S_{j(t)}(t)]=\begin{cases}
            \frac{\lambda_1(av+\mu)}{(a+b)^2}\left(\frac{a+b}{a+b-d}\right)^{\frac{\lambda_1+\lambda_0}{\lambda_1}}\Gamma(1+\lambda_0/\lambda_1) & \text{if }\lim_{t\to\infty}j(t)=\infty\\
            (av+\mu)I_\infty(N) & \text{if }\lim_{t\to\infty}j(t)=N,\,N<\infty
        \end{cases}
    \end{align*}
\end{itemize}
\end{corollary}
\begin{proof}
    Let $j(t)\gg \exp(t(\lambda_0\lambda_1/(\lambda_0+\lambda_1)))$.  So, $e^{\lambda_0 t}j(t)^{-(\lambda_1+\lambda_0)/\lambda_1}\ll 1$, so that using \ref{eq:meanSFS_varying} and that $I_t(n)$ is uniformly bounded over all $t\geq 0$ and $n$:
    \[\mathbb{E}[|S_{j(t)}(t)|]=\mathbb{E}[S_{j(t)}(t)]\to 0.\]
    The limits on mean in the second two cases once again follow from using the expression \ref{eq:meanSFS_varying} with $j(t)\sim C\exp(t\lambda_0\lambda_1/(\lambda_0+\lambda_1))$ and $j(t)\ll \exp(t(\lambda_0\lambda_1/(\lambda_0+\lambda_1)-\varepsilon))$.
\end{proof}

For the proof of point (2) of the theorem, it will be useful to control the term $P_j(t)$:
\begin{lemma}\label{lem:P_j_bound}
    $P_j(t)\leq 1/j$.
\end{lemma}
\begin{proof}
    Let
    \[\beta(t)=\frac{(a+b)e^{\lambda_1 t}-(a+b)}{(a+b)e^{\lambda_1 t}-d}.\]
    Then using \ref{eq:P_j}, we can bound $P_j(s)$ for any time $t$ by writing:
    \begin{align*}
        P_j(t)\leq (1-\beta(t))\beta(t)^{j-1}.
    \end{align*}
    Note that $\beta(z)\in [0,1]$, so that
    \begin{align*}
        \leq \max_{\beta\in [0,1]}(1-\beta)\beta^{j-1}\leq \frac{1}{j}.
    \end{align*}
\end{proof}

Now, we examine characteristic functions and then use the Levy continuity theorem to prove part 2 of Theorem \ref{thm:SFSSparsity}.

\begin{lemma}\label{lem:sparsity_part2}
    Suppose $j(t)\sim C\exp(t\lambda_0\lambda_1/(\lambda_0+\lambda_1))$ and $L=\lim_{t\to\infty}\mathbb{E}[S_{j(t)}(t)]$. Then:
    \[\left|\mathbb{E}\left[e^{i\theta S_{j(t)}(t)}\mathbf{1}_{\Omega_0^\infty}\right]-\mathbb{E}\left[e^{YL(e^{i\theta}-1)}\mathbf{1}_{\Omega_0^\infty}\right]\right|\overset{t\to\infty}{\to} 0\]
\end{lemma}
\begin{proof}
    We examine the difference of the characteristic functions of $S_{j(t)}(t)$ and the proposed limit in distribution:
    \begin{align}
        \left|\mathbb{E}\left[e^{i\theta S_{j(t)}(t)}\mathbf{1}_{\Omega_0^\infty}\right]-\mathbb{E}\left[e^{YL(e^{i\theta}-1)}\mathbf{1}_{\Omega_0^\infty}\right]\right|&\leq \mathbb{E}\left[\left|\mathbb{E}[e^{i\theta S_{j(t)}(t)}|\mathcal{B}_\infty] -e^{\mathbb{E}[S_{j(t)}|\mathcal{B}_\infty](e^{i\theta}-1)}\right|\right]\label{eq:boundary_term_1}\\
        &+\mathbb{E}\left[\left|e^{\mathbb{E}[S_{j(t)}|\mathcal{B}_\infty](e^{i\theta}-1)}-e^{\overline{S}_{j(t)}(t)(e^{i\theta}-1)}\right|\right]\label{eq:boundary_term_2}\\
        &+\mathbb{E}\left[\left|e^{\overline{S}_{j(t)}(t)(e^{i\theta}-1)}-e^{\widehat{S}_{j(t)}(t)(e^{i\theta}-1)}\right|\right]\label{eq:boundary_term_3}\\
        &+\mathbb{E}\left[\left|e^{\widehat{S}_{j(t)}(t)(e^{i\theta}-1)}-e^{YL(e^{i\theta}-1)}\right|\right]\label{eq:boundary_term_4}
    \end{align}
    We now show that each term above tends to $0$ as $t\to\infty$. We start with Term \ref{eq:boundary_term_1}. Remark that:
    \begin{align*}
        \mathbb{E}[e^{i\theta S_{j(t)}(t)}|\mathcal{B}_\infty]=\prod_{T_i\leq t}\left(1+P_{j(t)}(t-T_i)(e^{i\theta}-1)\right).
    \end{align*}
    Furthermore:
    \begin{align*}
        e^{\mathbb{E}[S_{j(t)}(t)|\mathcal{B}_\infty](e^{i\theta}-1)}=\prod_{T_i\leq t}e^{P_j(t-T_i)(e^{i\theta}-1)}.
    \end{align*}
    Now, notice that for complex numbers $a_k,b_k$ where $1\leq k\leq n$, if $|a_k|,|b_k|\leq 1$ then:
    \begin{align*}
        \left|\prod a_k-\prod b_k\right|\leq \sum_{k} |a_k-b_k|.
    \end{align*}
    Remark further that complex number $z$ with $\text{Re}(z)\leq 0$ and any $k\geq 0$:
    \begin{align}
        \left|e^{z}-\sum_{\ell=0}^k z^\ell/\ell!\right|\leq \frac{|z|^{k+1}}{(k+1)!}.\label{eq:exponential_bound}
    \end{align}
    Combining these facts:
    \begin{align*}
        \left|\prod_{T_i\leq t}\left(1+P_{j(t)}(t-T_i)(e^{i\theta}-1)\right)-\prod_{T_i\leq t}e^{P_j(t-T_i)(e^{i\theta}-1)}\right|&\leq \sum_{T_i\leq t}\left|1+P_{j(t)}(t-T_i)(e^{i\theta}-1)-e^{P_{j(t)}(t-T_i)(e^{i\theta}-1)}\right|\\
        &\leq 2 \sum_{T_i\leq t}P_{j(t)}(t-T_i)^2.
    \end{align*}
    We then take expectation over the times $\{T_i\}_i$ apply Lemma \ref{lem:P_j_bound} on a single factor $P_{j(t)}(t-T_i)$ and use Lemma \ref{lem:changing_j_mean}:
    \begin{align*}
        \ref{eq:boundary_term_1}\leq \frac{1}{2j(t)}\mathbb{E}[S_{j(t)}(t)]\to 0.
    \end{align*}

    Now the terms \ref{eq:boundary_term_2}, \ref{eq:boundary_term_3}, \ref{eq:boundary_term_4} will tend to $0$ if it is shown that $\mathbb{E}[S_{j(t)}|\mathcal{B}_\infty]-\overline{S}_{j(t)}(t),\overline{S}_{j(t)}(t)-\widehat{S}_{j(t)}(t),\widehat{S}_{j(t)}(t)-YL$ converge to $0$ in $L^1$, respectively. This is because for $A,B\geq 0$:
    \[\left|e^{A\left(e^{i\theta}-1\right)}-e^{B\left(e^{i\theta}-1\right)}\right|=\left|e^{(A-B)\left(e^{i\theta}-1\right)}-1\right|\leq |A-B|\left|e^{i\theta}-1\right|\leq 2|A-B|\]
    with the inequality coming from \ref{eq:exponential_bound}. This is what we will show using the ideas from Proposition \ref{prop:L2_driver} and Lemma \ref{lem:pm_approx_via_Y}. Firstly:
    \begin{align*}
      \ref{eq:boundary_term_2}&\leq 2\mathbb{E}\left[\left|\mathbb{E}[S_{j(t)}(t)|\mathcal{B}_\infty]-\overline{S}_{j(t)}(t)\right|\right]\\
      &\leq 2\sqrt{\mathbb{E}\left[\left(\int_0^t P_{j(t)}(t-s)\,dR_s\right)^2\right]}\\
      &= 2\sqrt{\mathbb{E}\left[\int_0^t P_{j(t)}(t-s)^2\,d[R]_s\right]}\\
      &=2\frac{1}{\sqrt{j(t)}}\sqrt{\mathbb{E}[S_{j(t)}(t)]}\to 0.
    \end{align*}

    For term \ref{eq:boundary_term_3}, we notice have:
    \begin{align*}
      \ref{eq:boundary_term_3}&\leq 2\mathbb{E}\left[\left|\overline{S}_{j(t)}(t)-\widehat{S}_{j(t)}(t)\right|\right]\\
      &\leq 2(av+\mu)\int_0^t P_{j(t)}(t-s)\sqrt{\mathbb{E}[(Z_0(s)-Ye^{\lambda_0 s})^2]}\,ds.\
    \end{align*}

    From Lemma \ref{lem:L2_distance}, the $L^2$ distance between $Z_0(s)$ and $Ye^{\lambda_0 s}$ is $\mathcal{O}(e^{\lambda_0 s})$. We can also use Lemma \ref{lem:P_j_bound} and that $\frac{\lambda_0\lambda_1}{\lambda_0+\lambda_1}>\frac{1}{2}\lambda_0$ to see:
    \begin{align*}
      &\leq C\int_0^t P_{j(t)}(t-s)e^{s\lambda_0/2}\,ds\\
      &\leq C\frac{te^{t\lambda_0/2}}{j(t)}\to 0.
    \end{align*}

    Remark that $\widehat{S}_{j(t)}(t)=Y\mathbb{E}[S_{j(t)}(t)]$ so that trivially:
    \begin{align*}
      \ref{eq:boundary_term_4}\leq 2\mathbb{E}[Y]|\mathbb{E}[S_{j(t)}(t)]-L|\to 0.
    \end{align*}

\end{proof}

\begin{lemma}\label{lem:sparsity_part3}
Suppose $j(t)\ll \exp(t\lambda_0\lambda_1/(\lambda_0+\lambda_1))$. Let $L=\lim_{t\to\infty}\mathbb{E}[e^{-\lambda_0 t}j(t)^{\frac{\lambda_1+\lambda_0}{\lambda_1}}S_{j(t)}(t)]$. Then:
\[\mathbb{E}[|e^{-\lambda_0 t}j(t)^{\frac{\lambda_1+\lambda_0}{\lambda_1}}S_{j(t)}(t)-YL|]\overset{t\to\infty}{\to}0.\]
\end{lemma}

\begin{proof}
  As in Section \ref{app:SFS_convergence}, we may decompose
  \begin{align}
    |S_{j(t)}(t)-e^{\lambda_0 t}j(t)^{-\frac{\lambda_1+\lambda_0}{\lambda_1}}YL|&\leq |S_{j(t)}(t)-\mathbb{E}[S_{j(t)}(t)|\mathcal{B}_\infty]|\label{eq:beta_term_1}\\
    &+|\mathbb{E}[S_{j(t)}(t)|\mathcal{B}_\infty]-\overline{S}_{j(t)}(t)|\label{eq:beta_term_2}\\
    &+|\overline{S}_{j(t)}(t)-\widehat{S}_{j(t)}(t)|\label{eq:beta_term_3}\\
    &+|\widehat{S}_{j(t)}(t)-Ye^{\lambda_0 t}j(t)^{-\frac{\lambda_1+\lambda_0}{\lambda_1}}L|.\label{eq:beta_term_4}
  \end{align}

  Then using analogs of Proposition \ref{prop:L2_driver} and Lemma \ref{lem:pm_approx_via_Y} but replacing $P_{j,\pm}^k$ with $P_{j(t)}$ and using expression \ref{eq:meanSFS_varying}:
  \begin{align*}
    \max\{\mathbb{E}[(S_{j(t)}(t)-\mathbb{E}[S_{j(t)}(t)|\mathcal{B}_\infty])^2],\mathbb{E}[(\mathbb{E}[S_{j(t)}(t)|\mathcal{B}_\infty]-\overline{S}_{j(t)}(t))^2],\mathbb{E}[(\overline{S}_{j(t)}(t)-\widehat{S}_{j(t)}(t))^2]\}\lesssim e^{\lambda_0 t}j(t)^{-\frac{\lambda_1+\lambda_0}{\lambda_1}}.
  \end{align*}
Therefore:
  \begin{align*}
    e^{-\lambda_0 t}j(t)^{\frac{\lambda_1+\lambda_0}{\lambda_1}}\mathbb{E}[|S_{j(t)}(t)-\widehat{S}_{j(t)}(t)|]\lesssim e^{-\frac{\lambda_0}{2}t}j(t)^{\frac{\lambda_1+\lambda_0}{2\lambda_1}}\to 0.
  \end{align*}
Finally, term \ref{eq:beta_term_4} is controlled in a similar manner to how \ref{eq:boundary_term_4} was controlled in the previous lemma. Namely:
\begin{align*}
    \mathbb{E}[\ref{eq:beta_term_4}]=|\mathbb{E}[S_{j(t)}(t)]-e^{\lambda_0 t}j(t)^{-\frac{\lambda_1+\lambda_0}{\lambda_1}}L|\ll e^{\lambda_0 t}j(t)^{-\frac{\lambda_1+\lambda_0}{\lambda_1}}
\end{align*}
Combining everything:
\[\mathbb{E}\left|e^{-\lambda_0 t}j(t)^{\frac{\lambda_1+\lambda_0}{\lambda_1}}S_{j(t)}(t)-YL\right|\to 0.\]
\end{proof}

To prove the limits regarding $j_k^-(t),j_k^+(t)$, it suffices to prove the following on $\Omega_0^\infty$:

\[\frac{\lambda_0\lambda_1}{\lambda_0+\lambda_1}\leq \lim_{k\to\infty}\liminf_{t\to\infty}\frac{\log(j_k^-(t))}{t},\lim_{k\to\infty}\limsup_{t\to\infty}\frac{\log(j_k^+(t))}{t}\leq \frac{\lambda_0\lambda_1}{\lambda_0+\lambda_1}.\]
This is because we trivially have $j_k^-(t)\leq j_{k+1}^+(t)+1$ for every $t\geq 0$ and $k$. Before proving these two inequalities, we require two technical lemmas.

\begin{lemma}\label{lem:V_moments_finite}
    Define:
    \[V:=\sum_{i}e^{-2\lambda_0 T_i}.\]
    For all $k$, $\mathbb{E}[V^k]<\infty$.
\end{lemma}
\begin{proof}
    We prove this result by induction on $k$. In the $k=1$ case, the result follows by Campbell's theorem. Next assume the induction hypothesis $\mathbb{E}[V^{i}]<\infty$ for $1\leq i\leq k-1$. By examining what occurs during the first event (ocurring at time $\tau$, with $E_b,E_d,E_{v},E_{\mu}$ denoting birth without mutation, death, replication-dependent mutation, and replication-free mutation), we may write:
    \[V\overset{d}{=}\mathbf{1}_{E_b}e^{-2\lambda_0 \tau}(V_1+V_2)+\mathbf{1}_{E_v}(e^{-2\lambda_0 \tau}+e^{-2\lambda_0 \tau}V_1)+\mathbf{1}_{E_\mu}e^{-2\lambda_0 \tau},\]
    where $V_1,V_2$, conditional on the first event, are independent and distributed as $V$. Taking powers and expectation, we obtain:
    \begin{align*}
        \mathbb{E}[V^k]=\frac{a(1-v)\mathbb{E}[(V_1+V_2)^k]}{a+d+\mu+2\lambda_0 k}+\frac{av\mathbb{E}[(1+V)^k]}{a+d+\mu+2\lambda_0 k}+\frac{\mu}{a+d+\mu+2\lambda_0 k}.
    \end{align*}
    Let $C_{<k}$ denote some linear combination of terms of the form $\mathbb{E}[V^{i_1}]\mathbb{E}[V^{i_2}]$ where $0\leq i_1,i_2<k$. We can rewrite the above using the leading order terms:
    \begin{align*}
        &\mathbb{E}[V^k]=\frac{(a+a(1-v))\mathbb{E}[V^k]}{a+d+\mu+2\lambda_0 k}+C_{<k}\\
        \Longrightarrow &\mathbb{E}[V^k]= \frac{a+d+\mu+2\lambda_0 k}{(2k-1)\lambda_0}C_{<k}<\infty,
    \end{align*}
    as desired.
\end{proof}

\begin{lemma}\label{lem:large_deviation_bound}
    We have the following large-deviation bound:
    \begin{align*}
        \mathbb{P}(\inf_{t\geq 0} e^{-\lambda_0 t}Z_0(t)\leq x|\Omega_0^\infty)\lsimas{x}{0^+}\sqrt{x}.
    \end{align*}
\end{lemma}
\begin{proof}
    Remark that on $(\Omega_0^\infty)^c$, $\inf_t Z_0(t)e^{-\lambda_0 t}\leq x$ occurs (for any $x>0$). Thus we may write:
    \begin{align}
        \mathbb{P}(\inf_t Z_0(t)e^{-\lambda_0t}\leq x|\Omega_0^\infty)=\frac{\mathbb{P}(\inf_t Z_0(t)e^{-\lambda_0t}\leq x)-p_0}{q_0}\label{eq:large_deviation_rewrite}.
    \end{align}
    For each $\alpha>0$, we have $e^{-\alpha Z_0(t)e^{-\lambda_0 t}}$ is a $\mathcal{F}_t$-submartingale. So using Doob's submartingale inequality:
    \[\mathbb{P}(\inf_t Z_0(t)e^{-\lambda_0t}\leq x)=\mathbb{P}(\sup_t e^{-\alpha Z_0(t)e^{-\lambda_0 t}}\geq e^{-\alpha x})\leq e^{\alpha x}\mathbb{E}[e^{-\alpha Y}]=e^{\alpha x}\left(p_0+\frac{q_0^2}{q_0+\alpha}\right)\]
    Notice we may optimize $\alpha$ in the upper bound by choosing $\alpha^*=\frac{\sqrt{q_0^4+\frac{4p_0q_0^2}{x}}}{2p_0}-\frac{q_0^2}{2p_0}$. In particular, there exists a positive constant $D$ such that:
    \[\inf_\alpha \left(e^{\alpha x}\left(p_0+\frac{q_0^2}{q_0+\alpha}\right)\right)=p_0+D\sqrt{x}+\mathcal{O}_{x\to 0^+}(x).\]
    The leading order term is constant, but by plugging the expression into \ref{eq:large_deviation_rewrite}, this term cancels which yields the result.
\end{proof}

We will use the above result in order to bound how small the SFS can be.

\begin{lemma}\label{lem:upper_bound_Sj_strong}
    We obtain the following bound on $S_{j}(t)$ on the event $B=\{\inf_{t>0} Z_0(t)e^{-\lambda_0 t}>x\}$ for $x>0$, $k$, and $\delta\in (0,1)$:
    \begin{align}
        \mathbb{P}(S_j(t)\leq k,B)\leq C_{k,\delta}e^{-x(1-\delta)\int_0^t P_j(t-s)ave^{\lambda_0 s}\,ds}=C_{k,\delta}e^{-x(1-\delta)\frac{av}{av+\mu}\mathbb{E}[S_j(t)]}.\label{eq:S_j_upper_bound}
    \end{align}
    $C_{k,\delta}>0$ and is independent of $j$.
\end{lemma}
\begin{proof}
    Let $S_j'(t)$ denote the contribution to the SFS from clones produced in a replication-dependent manner. Trivially, $\mathbb{P}(S_j(t)\leq k,B)\leq \mathbb{P}(S_j'(t)\leq k,B)$. It suffices to bound this latter term. Remark that the initiation times of replication-dependent mutants are seeded as a Poisson point process $\Phi'$ with intensity $avZ_0(s)$ conditional on $\mathcal{F}_\infty$.

    Let $\widetilde{\Phi}_t$ be a $p_t(s)=P_j(t-s)$-thinning of $\Phi'$. Thus $\widetilde{\Phi}_t$ remains a Poisson point process with intensity $avp_t(s)Z_0(s)$ conditional on $\mathcal{F}_\infty$ and furthermore $\mathbb{P}(S_{j}'(t)\leq k,B)=\mathbb{P}(|\widetilde{\Phi}_t\cap [0,t]|\leq k,B)$. Then, since $B$ is $\mathcal{F}_\infty$-measurable:
    \begin{align*}
        \mathbb{P}(S_{j}'(t)\leq k,B)&=\mathbb{E}\left[\mathbf{1}_{B}\mathbb{E}[\mathbf{1}_{|\widetilde{\Phi}_t\cap [0,t]|\leq k}|\mathcal{F}_\infty]\right]\\
        &=\sum_{n=0}^k\mathbb{E}\left[\mathbf{1}_B\frac{1}{n!}\left(\int_0^t P_{j}(t-s)avZ_0(s)\,ds\right)^ne^{-\int_0^t P_{j}(t-s)avZ_0(s)\,ds}\right].
    \end{align*}
    At the expense of a $(1-\delta)$ multiple of the exponent term, we may remove the power $n$ term:
    \begin{align*}
        \leq C_{k,\delta}\mathbb{E}\left[\mathbf{1}_Be^{-(1-\delta)\int_0^t P_{j}(t-s)avZ_0(s)\,ds}\right]
    \end{align*}
    We can then trivially write $Z_0(s)\geq e^{\lambda_0 s}\inf_t Z_0(t)e^{-\lambda_0 t}$ and the desired result follows by using the definition of the event $B$.
\end{proof}

Finally, we prove the main results.

\begin{proposition}\label{prop:intermediate_j_k_plus}
    There exists constants $C_1>0$ such that if $k>\frac{C_1}{\varepsilon}$, then:
    \begin{align*}
        \limsup_{t\to\infty}\frac{\log(j_k^+(t))}{t}\leq \frac{\lambda_0\lambda_1}{\lambda_0+\lambda_1}+\varepsilon
    \end{align*}
    almost surely.
\end{proposition}
\begin{proof}
    Fix any $\varepsilon>0$. Define $j(t)\sim e^{\left(\frac{\lambda_0\lambda_1}{\lambda_0+\lambda_1}+\varepsilon\right)t}$ such that for each $n$, $j(t)=j_n$ is constant for $t\in [n,n+1)$. Suppose we show that:
  \begin{align}
    \sum_{n=1}^\infty \mathbb{P}\left(\exists j\geq j_n,\sup_{t\in [n,n+1)}S_{j}(t)\geq k+1\right)<\infty\label{eq:to_show_upper}
  \end{align}
  for sufficiently large $k$ as in the proposition statement. Then by Borel-Cantelli, there is a $T=T(\omega)$ (depending on the sample point) such that $t\geq T$ implies:
  \[\frac{\log(j_k^+(t))}{t}\leq \frac{\lambda_0\lambda_1}{\lambda_0+\lambda_1}+\varepsilon\Longrightarrow \limsup_{t\to\infty}\frac{\log(j_k^+(t))}{t}\leq \frac{\lambda_0\lambda_1}{\lambda_0+\lambda_1}+\varepsilon,\]
  almost surely, as desired. It remains to show \ref{eq:to_show_upper}. Perform a union bound:
  \begin{align}
      \mathbb{P}\left(\exists j\geq j_n,\,\sup_{t\in [n,n+1)}S_{j}(t)\geq k+1\right)\leq \sum_{j\geq j_n}\mathbb{P}\left(\sup_{t\in [n,n+1)}S_{j}(t)\geq k+1\right).\label{eq:union_bound_jn}
  \end{align}
  Firstly, we can break up each of the events in \ref{eq:union_bound_jn} by whether $S_{j_n}(n)$ is larger than $k+1$ or not:
  \begin{align}
    \sum_{j\geq j_n}\mathbb{P}\left(\sup_{t\in [n,n+1)}S_{j}(t)\geq k+1\right)\leq \sum_{j\geq j_n}\mathbb{P}\left(S_{j}(n)\geq k+1\right)+\sum_{j\geq j_n}\mathbb{P}\left(S_{j}(n)\leq k,\sup_{t\in [n,n+1)}S_{j}(t)\geq k+1\right).\label{eq:sup_S_j_n_trivial_bound}
  \end{align}
  Define event $E_n^j:=\{S_{j}(n)\leq k,\sup_{t\in [n,n+1)}S_{j}(t)\geq k+1\}$. Letting $m$ denote the Lebesgue measure on $\mathbb{R}$, we define the local time $L_{B,j}^{\geq k}:=m\{t\in B:S_j(t)\geq k\}$ for Borel set $B$. To control the second term, note that on event $E_n^j$, $\zeta_n:=\inf\{t\geq n:S_{j}(t)\geq k+1\}$ is an almost surely finite stopping-time (with respect to the filtration by $\mathcal{I}_t$) and stops within $[n,n+1)$. Remark that since $(Z_0,Z_1)$ is a linear birth-death process, $S_j$ changes by increments of $1$ (almost surely), so that by right-continuity $S_{j}(\zeta_n)=k+1$. Starting from time $\zeta_n$, a necessary condition for $S_j(t)\leq k$ for $t\geq \zeta_n$ is that \textit{some} event occur to one of the original $k+1$ clones of size $j_n$ present at time $\zeta_n$ by time $t$. Thus:
  \[L_{[n,n+2],j}^{\geq k+1}\geq \mathbf{1}_{E_n^j}\min\{\eta,1\},\]
  where -- by the Strong Markov property -- conditional on $\mathcal{I}_{\zeta_n}$, $\eta$ is the minimum of $k+1$ independent exponentials each of rate $(a+b+d)j$. Using that $E_n^j$ is $\mathcal{I}_{\zeta_n}$-measurable and taking expectations:
  \begin{align}
    \mathbb{P}(E_n^j)\leq \frac{j(k+1)}{1-e^{-j(k+1)(a+b+d)}}\mathbb{E}[L_{[n,n+2]}^{\geq k+1}].\label{eq:sup_bound}
  \end{align}

  Remark that:
  \begin{align}
    L_{[n,n+2]}^{\geq k+1}=\int_n^{n+2}\mathbf{1}_{S_{j}(s)\geq k+1}\,ds\label{eq:local_time_decomposition}
  \end{align}
  Therefore, to bound \ref{eq:sup_bound}, we bound $\mathbb{P}(S_{j}(s)\geq k+1)$. Recall that $S_{j}(t)$, conditional on $\mathcal{B}_\infty$, is a sum of independent indicators, so that upon taking a factorial bound:
    \begin{align}
        \mathbb{P}(S_{j}(s)\geq k+1)\leq \mathbb{E}\left[\frac{\mathbb{E}[S_{j}(s)|\mathcal{B}_\infty]^{k+1}}{(k+1)!}\right].\label{eq:factorial_bound}
    \end{align}
    Now, remark that:
    \begin{align*}
        \mathbb{E}[S_j(t)|\mathcal{B}_\infty]=\int_0^t P_j(t-s)\,dM_s=\int_0^t M_{s-}P_j'(t-s)\,ds,
    \end{align*}
    with the last step by integration by parts. Then:
    \begin{align}
        &\leq  \sup_{z\in [0,\infty)} M_ze^{-\lambda_0 z}\left(P_j(t)+\lambda_0\int_0^t e^{\lambda_0 s}P_j(t-s)\,ds\right)\nonumber\\
        &\leq  \sup_{z\in [0,\infty)} M_ze^{-\lambda_0 z}\left(P_j(t)+\lambda_0 e^{\lambda_0 t}\int_0^\infty e^{-\lambda_0 u}P_j(u)\,du\right)\nonumber\\
        &\leq C\sup_{z\in [0,\infty)} \left(M_ze^{-\lambda_0 z}\right) e^{\lambda_0 t}j^{-1-\frac{\lambda_0}{\lambda_1}},\label{eq:S_j_expectation_bound}
    \end{align}
    with the last step coming from the asymptotic for $\mathbb{E}[e^{-\lambda_0 t}S_j(t)]$ (Proposition \ref{prop:SFS_mean}). Remark that $\sup_{z\in [0,\infty)}\left(M_z e^{-\lambda_0 z}\right)$ has all finite moments. To see this, recall from Lemma \ref{lem:CenteredMuts_Martingale} the decomposition $M_z=R_z+(av+\mu)\int_0^z Z_0(s)\,ds$ where $R_z$ is a $\mathcal{B}_z$-martingale. It suffices to show then that $\sup_{z\in [0,\infty)}e^{-\lambda_0 z}|R_z|$ and $\sup_{z\in [0,\infty)}e^{-\lambda_0 z}\int_0^z Z_0(s)\,ds$ have finite moments. For the latter term:

    \begin{align*}
        \sup_{z\in [0,\infty)}e^{-\lambda_0 z}\int_0^z Z_0(s)\,ds\leq \sup_{s\in [0,\infty)}(e^{-\lambda_0 s}Z_0(s))\sup_{z\in [0,\infty)}\frac{1}{\lambda_0}(1-e^{-\lambda_0 z})=\frac{1}{\lambda_0}\sup_{z\in [0,\infty)}e^{-\lambda_0 s}Z_0(s). 
    \end{align*}
The $p$th moment of the upper bound is controlled by Doob's maximal inequality:
    \begin{align*}
        \mathbb{E}\left[\sup_{z\in [0,\infty)}\left(e^{-\lambda_0 s}Z_0(s)\right)^p\right]\leq C_p\sup_{t\in [0,\infty)}\mathbb{E}[\left(e^{-\lambda_0 t}Z_0(t)\right)^p]<\infty,
    \end{align*}
    where the final inequality follows from e.g. Lemma 5 of \cite{foo2014escape}.
To control the term $e^{-\lambda_0 z}|R_z|$, define the $\mathcal{B}_t$-martingale $Q_t=\int_0^t e^{-\lambda_0 s}dR_s$. Using integration by parts:
    \begin{align*}
        |R_t|=\left|\int_0^t e^{\lambda_0 s}\,dQ_s\right|\leq |Q_t|e^{\lambda_0 t}+\lambda_0 \int_0^t \left|Q_{s-}\right|e^{\lambda_0 s}\,ds.
    \end{align*}
    Multiplying by $e^{-\lambda_0 t}$ and taking a supremum over $z\in [0,\infty)$:
    \begin{align*}
        \sup_{z\in [0,\infty)} e^{-\lambda_0 z}|R_z|\leq 3\sup_{z\in [0,\infty)}|Q_z|.
    \end{align*}
    We can control the upper bound using the BDG-inequalities, noting that $[Q]_t=\int_0^t e^{-2\lambda_0 s}d[R]_s=\int_0^t e^{-2\lambda_0 s}\,dM_s$:
    \begin{align*}
        \mathbb{E}[\sup_{z\in [0,\infty)}|Q_z|^p]\leq C_p \mathbb{E}\left[\left(\sum_{T_i}e^{-2\lambda_0 T_i}\right)^{p/2}\right].
    \end{align*}
    The last term is finite by using Lemma \ref{lem:V_moments_finite} and applying Jensen's inequality.
    
    Combining \ref{eq:factorial_bound} and \ref{eq:S_j_expectation_bound} then gives for $s\in [n,n+2]$:
    \begin{align}
      \mathbb{P}(S_{j}(s)\geq k+1)\leq C_2 e^{(k+1)\lambda_0t}j^{-(1+\frac{\lambda_0}{\lambda_1})(k+1)}.\label{eq:S_j_n_tight_bound}
    \end{align}
    Summing over all $j\geq j_n$ gives
    \begin{align}
        \sum_{j\geq j_n}\mathbb{P}(S_j(s)\geq k+1)\leq C_3e^{(k+1)\lambda_0 t}j_n^{-(1+\frac{\lambda_0}{\lambda_1})(k+1)+1}\leq C_4e^{-(c_1^\varepsilon k-c_2)n},\label{eq:S_j_n_sum_tight_bound}
    \end{align}
    where $c_1^\varepsilon,c_2>0$, $c_1^\varepsilon$ depends on $\varepsilon$ (scales linearly with $\varepsilon$), and $s\in [n,n+2]$. Therefore, using \ref{eq:local_time_decomposition} along with the union bound \ref{eq:union_bound_jn} implies that for $k$ such that $c_1^\varepsilon k-c_2>0$ (or equivalently for some $C_1>0$, $k>C_1/\varepsilon$), the second term on the right-hand-side of \ref{eq:sup_S_j_n_trivial_bound} is summable (over $n$). \ref{eq:S_j_n_sum_tight_bound} implies that the first term on the right-hand-side of \ref{eq:sup_S_j_n_trivial_bound} is summable (over $n$) as well. Thus, \ref{eq:to_show_upper} holds.
\end{proof}

\begin{proposition}\label{prop:limsup_j_k_plus}
  \begin{align}
      \lim_{k\to\infty}\limsup_{t\to\infty}\frac{\log(j_k^+(t))}{t}\leq \frac{\lambda_0\lambda_1}{\lambda_0+\lambda_1}\label{eq:limsup_j_k_plus}
  \end{align}
  on $\Omega_0^\infty$
\end{proposition}
\begin{proof}
    Remark that $k\mapsto j_k^+(t)$ is non-increasing in $k$. Therefore, the limit on the left-hand-side of \ref{eq:limsup_j_k_plus} exists. For any $\varepsilon>0$, there exists $k'$ large enough such that:
    \[\limsup_{t\to\infty}\frac{\log(j_{k'}^+(t))}{t}\leq \frac{\lambda_0\lambda_1}{\lambda_0+\lambda_1}+\varepsilon\]
    by Proposition \ref{prop:intermediate_j_k_plus}. Thus:
    \begin{align*}
        \lim_{k\to\infty}\limsup_{t\to\infty}\frac{\log(j_k^+(t))}{t}\leq \limsup_{t\to\infty}\frac{\log(j_{k'}^+(t))}{t}\leq \frac{\lambda_0\lambda_1}{\lambda_0+\lambda_1}+\varepsilon.
    \end{align*}
    As $\varepsilon$ is arbitrary, we conclude \ref{eq:limsup_j_k_plus} holds.
\end{proof}

\begin{proposition}\label{prop:liminf_j_k_mimus}
  \[\frac{\lambda_0\lambda_1}{\lambda_0+\lambda_1}\leq \lim_{k\to\infty}\liminf_{t\to\infty}\frac{\log(j_k^-(t))}{t}\]
  on $\Omega_0^\infty$
\end{proposition}
\begin{proof}
Fix any $\varepsilon>0$. Take $\lambda_1>\beta>0$. Define $j(t)\sim \exp\left((\frac{\lambda_0\lambda_1}{\lambda_0+\lambda_1}-\varepsilon)t\right)$ such that every $n$, $j(t)=j_n$ is constant over $t\in [n,n+1)$. Let $E_n=\{\sup_{t\in [n,n+1]}Z_1(t)\leq e^{(\lambda_1+\beta)n}\}$. Notice that since $e^{-\lambda_1 t}Z_1(t)$ converges almost surely to a finite random variable, $\mathbb{P}(E_n\text{ eventually})=1$. Then, to prove the claim it suffices to show:
  \begin{align}
    \sum_{n=1}^\infty \mathbb{P}\left(\exists j\leq j_n,\inf_{t\in [n,n+1)}S_{j}(t)\leq k,\Omega_0^\infty,E_n\right)<\infty\label{eq:to_show_lower}
  \end{align}

  Letting $x_n=n^{-4}$, the probability of the event $B_n^c=\{\inf_{t}Z_0e^{-\lambda_0 t}\leq x_n\}$ conditional on $\Omega_0^\infty$ is summable by Lemma \ref{lem:large_deviation_bound} and in particular $P(B_n^c\text{ i.o.})=0$. So, to prove \ref{eq:to_show_lower} it suffices to show:
\begin{align*}
    \sum_{n=1}^\infty \mathbb{P}\left(\exists j\leq j_n,\inf_{t\in [n,n+1)}S_{j}(t)\leq k,\Omega_0^\infty,E_n,B_n\right)<\infty.
  \end{align*}

  Now, partition $[n,n+1)$ into $M_n\ordas{n}{\infty}\exp(2\lambda_1 n)$ equally-spaced half-open intervals. Denote the intervals $\{[t_i,t_{i+1})\}_{i=1}^{M_n}$. By applications of unions bounds, it suffices to show:
  \begin{align*}
      \sum_{n=1}^\infty\sum_{j\leq j_n}\sum_{i=1}^{M_n} \mathbb{P}\left(\inf_{t\in [t_i,t_{i+1})}S_{j}(t)\leq k,E_n,B_n\right)<\infty.
  \end{align*}

  By \ref{eq:meanSFS_varying}, $\inf_{s\in [n,n+1),j\leq j_n}\mathbb{E}[S_{j}(s)]\gsimas{t}{\infty}\exp\left(n\left(\frac{\lambda_0+\lambda_1}{\lambda_1}\varepsilon\right)\right)$. Thus, using the first and third expression of Eq. \ref{eq:S_j_upper_bound} of Lemma \ref{lem:upper_bound_Sj_strong} with say $\delta=0.5$, we obtain $\max_{j\leq j_n,1\leq i\leq M_n}\mathbb{P}(B_n\cap S_j(t_i)\leq 2k)\lsimas{n}{\infty}\exp\left(-0.5n^{-4}e^{\varepsilon n}\right)$. This decays doubly-exponentially so that summing $\mathbb{P}(B_n\cap S_j(t_i)\leq 2k)$ over $(n,j,i)$ is finite. Therefore, it suffices to prove:
  \begin{align*}
      \sum_{n=1}^\infty\sum_{j\leq j_n}\sum_{i=1}^{M_n} \mathbb{P}\left(S_j(t_i)\geq 2k,\inf_{t\in [t_i,t_{i+1})}S_{j}(t)\leq k,Z_1(t_i)\leq e^{(\lambda_1+\beta)n}\right)<\infty.
  \end{align*}

  For the event to occur, it must be that at least $k$ events occur to the type-1 members present at the beginning of the interval (at time $t_i$) by the end of the interval (by time $t_{i+1}$). Hence, the above probability is dominated by the tail of a binomial $B$ with $\lceil e^{(\lambda_1+\beta)n}\rceil$ trials with success probability $p=1-\exp(-(a+b+d)/M_n)$. Using $M_n\gsimas{n}{\infty}\exp(2\lambda_1 n)$ and by a factorial tail bound:
  \[\mathbb{P}\left(S_j(t_i)\geq 2k,\inf_{t\in [t_i,t_{i+1})}S_{j}(t)\leq k,Z_1(t_i)\leq e^{(\lambda_1+\beta)n}\right)\leq \mathbb{P}(B\geq k)\leq \frac{(\lceil e^{(\lambda_1+\beta)n}\rceil p)^k}{k!}\lsimas{n}{\infty}e^{(\beta-\lambda_1)kn}\]
Further using that $M_n\lsimas{n}{\infty}\exp(2\lambda_1 n)$ and that $j_n\lsimas{n}{\infty}\exp(\frac{\lambda_0\lambda_1}{\lambda_0+\lambda_1}n)$:
  \[\sum_{j\leq j_n}\sum_{i=1}^{M_n} \mathbb{P}\left(S_j(t_i)\geq 2k,\inf_{t\in [t_i,t_{i+1})}S_{j}(t)\leq k,Z_1(t_i)\leq e^{(\lambda_1+\beta)n}\right)\lsimas{n}{\infty}e^{\frac{\lambda_0\lambda_1}{\lambda_0+\lambda_1}n+2\lambda_1n+(\beta-\lambda_1)kn}.\]
  Taking $k$ large enough ensures that \ref{eq:to_show_lower} holds.
  
\end{proof}

The second part of the corollary requires a novel proof. First, we consider a lemma regarding a bound on the laplace transform of the type-1 clone.
\begin{lemma}\label{lem:MGF_bound}
  Let $G(\theta,t)=\mathbb{E}[\exp(-\theta Z_1^{(1)}(t))]$ denote the laplace transform of a single type-one clone. For any fixed $t_0$, there exists a $K>0$ such that:
  \begin{align*}
    e^{\beta \exp(t(\lambda_0\lambda_1)/(\lambda_0+\lambda_1))}G(\beta,t)\leq e^{-Kt^2}
  \end{align*}
  for all $t\in [0,t_0]$ and some $\beta>0$ (depending on $t$)
\end{lemma}
\begin{proof}
  Let $y_{\beta}=1-e^{-\beta}$, $C_t^{(1)}=\exp(\lambda_1t)$, and $C_t^{(2)}=\frac{a+b}{\lambda_1} \cdot (C_t^{(1)}-1)$, and $C_t^{(3)}=\exp(t\lambda_0\lambda_1/(\lambda_0+\lambda_1))$. It is known that (see Eq. 4 of \cite{durrettbranching2015}):
  \begin{align}
    e^{\beta C_t^{(3)}}G(\beta,t)=(1-y_{\beta})^{-C_t^{(3)}}\left(1-\frac{y_{\beta}C_t^{(1)}}{1+y_{\beta}C_t^{(2)}}\right).\label{eq:laplace_rewrite}
  \end{align}
  We upper bound \ref{eq:laplace_rewrite}. First, notice that:
  \begin{align*}
    &-\log(1-y_\beta)=y_\beta+\sum_{i=2}^\infty\frac{y_\beta^i}{i}\leq y_\beta+\frac{y_\beta^2}{2}\frac{1}{1-y_\beta},\text{ and}\\
    &1-\frac{y_\beta C_t^{(1)}}{1+y_\beta C_t^{(2)}}\leq e^{-\frac{y_\beta C_t^{(1)}}{1+y_\beta C_t^{(2)}}}\leq e^{-y_\beta C_t^{(1)}+y_\beta^2 C_t^{(1)}C_t^{(2)}}.
  \end{align*}
  Therefore:
  \begin{align*}
    \ref{eq:laplace_rewrite}\leq e^{y_\beta(C_t^{(3)}-C_t^{(1)})+y_\beta^2\left(\frac{C_t^{(3)}}{2(1-y_\beta)}+C_t^{(1)}C_t^{(2)}\right)}.
  \end{align*}
  Note that $\frac{C_t^{(1)}-C_t^{(3)}}{C_t^{(1)}C_t^{(2)}+C_t^{(3)}}$ is uniformly bounded over all $t\geq 0$, so in particular we can choose a constant $C$ such that
  $C\frac{C_t^{(1)}-C_t^{(3)}}{C_t^{(1)}C_t^{(2)}+C_t^{(3)}}\in (0,1)$. Then, we can choose $\beta^*$ (depending on $t$) so that $y^\ast_t:=y_{\beta^\ast}=C\frac{C_t^{(1)}-C_t^{(3)}}{C_t^{(1)}C_t^{(2)}+C_t^{(3)}}$. Using this $\beta^*$ we get
  \begin{align*}
    e^{\beta^\ast C_t^{(3)}}G(\beta^\ast,t)\leq e^{y_t^\ast(C_t^{(3)}-C_t^{(1)})+(y_t^\ast)^2\left(C_t^{(3)}+C_t^{(1)}C_t^{(2)}\right)}=e^{\frac{(C_t^{(3)}-C_t^{(1)})^2}{C_t^{(1)}C_t^{(2)}+C_t^{(3)}}\left(-C+C^2(C_t^{(3)}-C_t^{(1)})\right)}.
  \end{align*}
  The result follows since $C_t^{(1)}-C_t^{(3)}\overset{t\to 0^+}{\sim}\left(\lambda_1-\frac{\lambda_0\lambda_1}{\lambda_0+\lambda_1}\right)t$, and $C_t^{(1)}C_t^{(2)}+C_t^{(3)}\to 1$ as $t\to 0^+$  and then choosing $K$ sufficiently large.

\end{proof}

\begin{lemma}[Simple Poisson Tail Bounds]\label{lem:poisson_tails}
  Let $(P_n)_n$ be a sequence of Poisson random variables with rates $(\lambda_n)_n$. Then for any $k$:
  \begin{enumerate}
    \item $\mathbb{P}(P_1\geq k)\leq \frac{\lambda_1^k}{k!}$
    \item If $\lambda_n\ll 1$ as $n\to\infty$, then $\mathbb{P}(P_n\geq k)\overset{n\to\infty}{\sim} \frac{\lambda_n^k}{k!}$
  \end{enumerate}
\end{lemma}
\begin{proof}
  For item 1:
  \[\mathbb{P}(P_1\geq k)= \sum_{z=k}^\infty \frac{\lambda_1^z}{z!}e^{-\lambda_1 }\leq \frac{\lambda_1^k}{k!}\sum_{z=k}^\infty \frac{\lambda_1^{z-k}}{(z-k)!}e^{-\lambda_1}=\frac{\lambda_1^k}{k!}\]

  Now, we examine item 2. This is as easy as item 1:
  \begin{align*}
    \mathbb{P}(P_n\geq k)=\frac{\lambda_n^k}{k!}\sum_{z=k}^\infty\frac{k!}{z!}\lambda_n^{z-k}e^{-\lambda_n}=\frac{\lambda_n^k}{k!}(1+o_{n\to\infty}(1)),
  \end{align*}
  with the $1+o_{n\to\infty}(1)$ term following since the sum above goes to $1$ as $n\to\infty$ by dominated convergence.
\end{proof}

\begin{proposition}\label{prop:limsup_theta}
   Let $J_\varepsilon(t)\lsimas{t}{\infty}\exp\left(\left(\frac{\lambda_0\lambda_1}{\lambda_0+\lambda_1}+\varepsilon\right)t\right)$ be monotone, then:
    \begin{align*}
        \limsup_{t\to\infty} S_{J_\varepsilon(t)}(t)\gsimas{\varepsilon}{0^+}\varepsilon^{-1}
    \end{align*}
    on $\Omega_0^\infty$.
\end{proposition}
\begin{proof}
  Define $k_\varepsilon:=\lfloor \frac{\lambda_0\lambda_1^2}{4\varepsilon(\lambda_0+\lambda_1)^2}\rfloor $. It would then suffice to show:
  \[\limsup_{t\to\infty}S_{J(t)}(t)\geq k_\varepsilon\]
  almost surely on $\Omega_0^\infty$. To do so, we consider $S_{J(t)}'(t)$, which is the SFS contributed only by mutants formed in a replication-dependent manner. Since $S_{J(t)}'(t)\leq S_{J(t)}(t)$, it then suffices to show:
  \[\limsup_{t\to\infty}S_{J(t)}'(t)\geq k_\varepsilon\]
  almost surely on $\Omega_0^\infty$. We further define $S_j^{',B}(t)$, for $B\subseteq \mathbb{R}$ a borel set, to be the SFS contributed only by mutants formed in a replication-dependent manner and at times contained in $B$. So, $S_{j}^{',B}(t):=\sum_{T_i'\in B,T_i'\leq t}\mathbf{1}_{Z_1^{',(i)}(t-T_i')=j}$. Recall that the mutation timings of replication-dependent mutants form a Poisson process conditional on $\mathcal{F}_\infty$. Then, conditional on $\mathcal{F}_\infty$, $\{\{T_i'\}\cap [n,n+1)\}_{n\in\mathbb{N}}$ form independent Poisson point processes on the intervals $[1,2),[2,3),\dots$. Thus, $\{\{S_{J(t)}^{',[n,n+1)}(t)\}_{t\geq 0}\}_{n\in\mathbb{N}}$ form independent processes conditional on $\mathcal{F}_\infty$. By the second Borel-Cantelli to prove the proposition it will then suffice to show the last equality in:
  \begin{align}
    \sum_{n=1}^\infty \mathbb{P}(\sup_{t}S_{J(t)}^{',[n,n+1)}(t)\geq k_\varepsilon|\mathcal{F}_\infty)\geq \sum_{n=1}^\infty \mathbb{P}(\sup_{t\in [\alpha n,\alpha(n+1))}S_{J(t)}^{',[n,n+1)}(t)\geq k_\varepsilon|\mathcal{F}_\infty)= \infty\label{eq:S_J_alpha_large}
  \end{align}
  on $\Omega_0^\infty$ for some $\alpha>1$. To choose an appropriate $\alpha$ we consider the following heuristic. If a clone is produced at time $\approx n$, then it will reach the boundary $J(t)$ at time $t^*$ determined by roughly $e^{\lambda_1 (t^*-n)}\approx e^{\left(\frac{\lambda_0\lambda_1}{\lambda_0+\lambda_1}+\varepsilon\right)t^*}$, so $t^*\approx \lambda_1n/\left(\lambda_1-\lambda_0\lambda_1/(\lambda_0+\lambda_1)-\varepsilon\right)$. Using this, set $\alpha:= \lambda_1/\left(\lambda_1-\lambda_0\lambda_1/(\lambda_0+\lambda_1)-\varepsilon\right)$ (and take $\varepsilon>0$ be small enough so that $\alpha>0$) in \ref{eq:S_J_alpha_large}. Define $L_{A,B}^k:=\int_{A} J(t)\mathbf{1}_{S_{J(t)}^{',B}(t)\geq k}\,dt$ to be a ``weighted'' local time. Then, we have the following equality of events (up to null sets):
  \begin{align*}
    \left\{\sup_{t\in [\alpha n,\alpha(n+1))}S_{J(t)}^{',[n,n+1)}(t)\geq k_\varepsilon\right\}&=\left\{\exists t\in [\alpha n,\alpha(n+1)),\,S_{J(t)}^{',[n,n+1)}(t)\geq k_\varepsilon\right\}\\
    &=\left\{L_{[\alpha n,\alpha (n+1)),[n,n+1)}^{k_\varepsilon}>0\right\}.
  \end{align*}

  By the Paley-Zygmund bound it suffices to show the last equality in the series:
  \begin{align}
    \sum_{n=1}^\infty \frac{\mathbb{E}\left[L_{[\alpha n,\alpha (n+1)),[n,n+1)}^{k_\varepsilon}|\mathcal{F}_\infty\right]^2}{\mathbb{E}\left[\left(L_{[\alpha n,\alpha (n+1)),[n,n+1)}^{k_\varepsilon}\right)^2|\mathcal{F}_\infty\right]}=\sum_{n=1}^\infty \frac{1}{\frac{\mathbb{V}\left[L^{k_\varepsilon}_{[\alpha n,\alpha(n+1)),[n,n+1)}|\mathcal{F}_\infty\right]}{\mathbb{E}\left[L^{k_\varepsilon}_{[\alpha n,\alpha(n+1)),[n,n+1)}|\mathcal{F}_\infty\right]^2}+1}=\infty\label{eq:tight_limsup_to_show}
  \end{align}
  almost surely on $\Omega_0^\infty$. To do so, we show that the conditional expectation grows sufficiently quickly and the conditional variance grows sufficiently slowly (but still both grow to $\infty$). Defining $\mu_n(t):=av\int_n^{n+1}Z_0(s)P_{J(t)}(t-s)\,ds$, $\mu_n(t)$ is the Poisson rate for $S_{J(t)}^{',[n,n+1)}$ conditional on $\mathcal{F}_\infty$. Hence by Fubini-Tonelli:
  \begin{align}
    \mathbb{E}\left[L_{[\alpha n,\alpha (n+1)),[n,n+1)}^k|\mathcal{F}_\infty\right]&=\int_{\alpha n}^{\alpha(n+1)}J(t)\mathbb{P}\left(S_{J(t)}^{',[n,n+1)}(t)\geq k|\mathcal{F}_\infty\right)\,dt\nonumber\\
    &=\int_{\alpha n}^{\alpha(n+1)}J(t)\sum_{r=k}^\infty\frac{(\mu_n(t))^r}{r!}e^{-\mu_n(t)}\,dt.\label{eq:L_limsup_larger_1}
\end{align}
Fix $t_n\in [\alpha n,\alpha(n+1))$ and $s_n\in [n,n+1)$. Choosing $n$ large enough, we can make $t_n-s_n$ arbitrarily large as $\alpha>1$. For $\varepsilon$ small, since $J_\varepsilon(t)\ll \exp(\lambda_1 t)$, by using \ref{eq:P_j}:
\begin{align*}
    P_{J(t_n)}(t_n-s_n)\overset{n\to\infty}{\sim} (a+b-d)e^{-\lambda_1 (t_n-s_n)},
\end{align*}
uniformly for any such choice of sequences $(t_n)_n,(s_n)_n$ where $(t_n,s_n)\in [\alpha,\alpha(n+1))\times [n,n+1)$. Fix $T$ large enough so that for all $s\geq T$, $\frac{1}{2}Ye^{\lambda_0 s}\leq Z_0(s)\leq 2Ye^{\lambda_0 s}$. Take $n\geq T$ and note that:
\begin{align*}
    \mu_n(t_n)\leq CY\int_n^{n+1}e^{\lambda_0 s}e^{-\lambda_1(t_n-s)}\,ds\leq CYe^{(\lambda_0+\lambda_1) n}e^{-\lambda_1\alpha n}\overset{n\to\infty}{\to}0
\end{align*}
as $n\to\infty$, since $\lambda_0+\lambda_1-\lambda_1\alpha<0$. Thus. for sufficiently large $n$:
\[e^{-\mu_n(t_n)}\geq e^{-1}.\]
Reciprocally:
\begin{align*}
    \mu_n(t_n)\geq CYe^{(\lambda_0+\lambda_1)n}e^{-\lambda_1\alpha n}
\end{align*}
for a different constant $C$. Thus, for sufficiently large $n$:
\begin{align*}
    \ref{eq:L_limsup_larger_1}&\geq CYe^{(\lambda_0+\lambda_1)kn}e^{-\lambda_1\alpha kn}\int_{\alpha n}^{\alpha(n+1)}e^{\left(\frac{\lambda_0\lambda_1}{\lambda_0+\lambda_1}+\varepsilon\right)t}\,dt\\
    &\geq CY e^{\left(\frac{\lambda_0\lambda_1}{\lambda_0+\lambda_1}+\varepsilon-k\lambda_1\right)\alpha n+(\lambda_0+\lambda_1)kn}.
\end{align*}
  Observe that setting $k=k_\varepsilon$:
  \begin{align}
    \left(\frac{\lambda_0\lambda_1}{\lambda_0+\lambda_1}+\varepsilon-k_\varepsilon\lambda_1\right)\alpha +(\lambda_0+\lambda_1)k_\varepsilon=\frac{3\lambda_0}{4}+\mathcal{O}_{\varepsilon\to 0^+}(\varepsilon)\label{eq:exp_minimum_growth}
  \end{align}
  Now we obtain upper bounds on the variance in the denominator of \ref{eq:tight_limsup_to_show}:
  \begin{align}
    \mathbb{V}\left[L_{[\alpha n,\alpha (n+1)),[n,n+1)}^k|\mathcal{F}_\infty\right]&=2\int_{\alpha n}^{\alpha(n+1)}\int_{s}^{\alpha(n+1)}J(t)J(s)\text{CoV}[\mathbf{1}_{S_{J(t)}^{',[n,n+1)}(t)\geq k},\mathbf{1}_{S_{J(s)}^{',[n,n+1)}(s)\geq k}|\mathcal{F}_\infty]\,dt\,ds\label{eq:variance_term_limsup}
  \end{align}
  We can write the covariance term as:
  \begin{align}
    &\text{CoV}[\mathbf{1}_{S_{J(t)}^{',[n,n+1)}(t)\geq k},\mathbf{1}_{S_{J(s)}^{',[n,n+1)}(s)\geq k}|\mathcal{F}_\infty]\nonumber\\
    &=\mathbb{P}\left(S_{J(t)}^{',[n,n+1)}(t)\geq k,S_{J(s)}^{',[n,n+1)}(s)\geq k|\mathcal{F}_\infty\right)-\mathbb{P}\left(S_{J(t)}^{',[n,n+1)}(t)\geq k|\mathcal{F}_\infty\right)\mathbb{P}\left(S_{J(s)}^{',[n,n+1)}(s)\geq k|\mathcal{F}_\infty\right).\label{eq:cov_decomposition}
  \end{align}
  Now, we stratify the clones which contribute towards $S_{J(t)}^{',[n,n+1)}(t)$ and $S_{J(s)}^{',[n,n+1)}(s)$. Let $\mathcal{C}_1^{s,t},\mathcal{C}_2^{s,t},\mathcal{C}_3^{s,t}\subseteq \mathbb{N}$ be random sets defined as follows:
  \begin{align*}
    &\mathcal{C}_1^{s,t}:=\{i\in \mathbb{N}:Z_1^{',(i)}(s-T_i')=J(s),Z_1^{',(i)}(t-T_i')\neq J(t)\}\\
    &\mathcal{C}_2^{s,t}:=\{i\in \mathbb{N}:Z_1^{',(i)}(s-T_i')\neq J(s),Z_1^{',(i)}(t-T_i')= J(t)\}\\
    &\mathcal{C}_3^{s,t}:=\{i\in \mathbb{N}:Z_1^{',(i)}(s-T_i')=J(s),Z_1^{',(i)}(t-T_i')= J(t)\}.
  \end{align*}
  In words, $\mathcal{C}_1^{s,t},\mathcal{C}_2^{s,t},\mathcal{C}_3^{s,t}$ respectively are the identities of the clones which contribute towards $S_{J(s)}^{',[n,n+1)}(s)$ but not $S_{J(t)}^{',[n,n+1)}(t)$, do not contribute towards $S_{J(s)}^{',[n,n+1)}(s)$ but do contribute towards $S_{J(t)}^{',[n,n+1)}(t)$, and contribute towards both $S_{J(s)}^{',[n,n+1)}(s)$ and $S_{J(t)}^{',[n,n+1)}(t)$, respectively. Let $X_i^{s,t}:=|\mathcal{C}_i^{s,t}|$ for $i\in \{1,2,3\}$ and thus $S_{J(s)}^{',[n,n+1)}(s)=X_1^{s,t}+X_3^{s,t}$ and $S_{J(t)}^{',[n,n+1)}(t)=X_2^{s,t}+X_3^{s,t}$. Here, $X_1^{s,t},X_2^{s,t},X_3^{s,t}$ are distributed as independent Poisson conditional on $\mathcal{F}_\infty$, since their distributions can be written in terms of thinnings of $\{T_i'\}_i\cap [n,n+1)$. To specify the rates of $X_i^{s,t}$, we introduce a notation for the transition function of a single type-1 clone $P_{n_1\to n_2}(t):=\mathbb{P}(Z_1^{(1)}(t)=n_2|Z_1^{(1)}(0)=n_1)$. Note that $P_{n_2}(t)=P_{1\to n_2}(t)$. $X_1^{s,t}$ has rate $\mu^{(1)}(s,t)=\int_n^{n+1}avZ_0(r)P_{J(s)}(s-r)\,dr(1-P_{J(s)\to J(t)}(t-s))$, $X_2^{s,t}$ has rate $\mu^{(2)}(s,t)=\sum_{z\neq J(s)}\int_n^{n+1}avZ_0(r)P_{z}(s-r)\,drP_{z\to J(t)}(t-s)$, and $X_3^{s,t}$ has rate $\mu^{(3)}(s,t)=\int_n^{n+1}avZ_0(r)P_{J(s)}(s-r)\,dr P_{J(s)\to J(t)}(t-s)$. Then by conditioning on the values $X_3^{s,t}$ can take:
  \begin{align*}
    \ref{eq:cov_decomposition}&\leq \mathbb{P}(X_1^{s,t}\geq k|\mathcal{F}_\infty)\mathbb{P}(X_2^{s,t}\geq k|\mathcal{F}_\infty)+\sum_{x=1}^{k-1}\mathbb{P}(X_1^{s,t}\geq k-x|\mathcal{F}_\infty)\mathbb{P}(X_2^{s,t}\geq k-x|\mathcal{F}_\infty)\mathbb{P}(X_3^{s,t}=x|\mathcal{F}_\infty)\\
    &+\mathbb{P}(X_3^{s,t}\geq k|\mathcal{F}_\infty)\\
    &-\mathbb{P}\left(S_{J(t)}^{',[n,n+1)}(t)\geq k|\mathcal{F}_\infty\right)\mathbb{P}\left(S_{J(s)}^{',[n,n+1)}(s)\geq k|\mathcal{F}_\infty\right).
  \end{align*}

  We can upper bound $\mathbb{P}(X_1^{s,t}\geq k|\mathcal{F}_\infty)\mathbb{P}(X_2^{s,t}\geq k|\mathcal{F}_\infty)$ by $\mathbb{P}\left(S_{J(t)}^{',[n,n+1)}(t)\geq k|\mathcal{F}_\infty\right)\mathbb{P}\left(S_{J(s)}^{',[n,n+1)}(s)\geq k|\mathcal{F}_\infty\right)$ in which case the first and last term cancel:
  \begin{align}
    &\leq \sum_{x=1}^{k-1}\mathbb{P}(X_1^{s,t}\geq k-x|\mathcal{F}_\infty)\mathbb{P}(X_2^{s,t}\geq k-x|\mathcal{F}_\infty)\mathbb{P}(X_3^{s,t}=x|\mathcal{F}_\infty)\label{eq:sum_limsup_bound}\\
    &+\mathbb{P}(X_3^{s,t}\geq k|\mathcal{F}_\infty).\label{eq:term_limsup_bound_2}
  \end{align}
  We control \ref{eq:sum_limsup_bound} term-by-term. For each summand, fixing $x\in \{1,\dots,k-1\}$ and using $\mu^{(1)}(s,t)\leq \mu_n(s)$, $\mu^{(2)}(s,t)\leq \mu_n(t)$, $\mu^{(3)}(s,t)\leq \min(\mu_n(s),\mu_n(t))$ and Lemma \ref{lem:poisson_tails}:
  \begin{align*}
    \mathbb{P}(X_1\geq k-x|\mathcal{F}_\infty)\mathbb{P}(X_2\geq k-x|\mathcal{F}_\infty)\mathbb{P}(X_3=x|\mathcal{F}_\infty)&\leq C_{x,k}\mu_n(s)^{k-x}\mu_n(t)^{k-x}(\mu^{(3)}(s,t))^{x}\\
    &\leq C_{x,k}\mu_n(s)^{k-x}\mu_n(t)^{k-x}\mu_n(t)^{\frac{x-1}{2}}\mu_n(s)^{\frac{x-1}{2}}\mu^{(3)}(s,t)\\
    &\leq C_{x,k}(\mu_n(t)^{2k-x-1}+\mu_n(s)^{2k-x-1})\mu^{(3)}(s,t)
  \end{align*}
Recall that $\mu_n$ goes to $0$ as $n\to\infty$, so eventually $\mu_n\leq 1$ for all large $n$ uniformly for $t,s\in [\alpha n,\alpha(n+1))$, in which case using $2k-x-1\geq k$:
  \begin{align}
    \leq C_{x,k}(\mu_n(t)^{k}+\mu_n(s)^{k})\mu^{(3)}(s,t)&=C_{x,k}\mu_n(t)^k\mu_n(s)P_{J(s)\to J(t)}(t-s)\label{eq:limsup_summand_1}\\
    &+C_{x,k}\mu_n(s)^{k+1}P_{J(s)\to J(t)}(t-s)\label{eq:limsup_summand_2},
  \end{align}
  and then using $\mu^{(3)}(s,t)=\mu_n(s)P_{J(s)\to J(t)}(t-s)$ for the equality.

  
  Define $\varepsilon_n:=\exp(-(3\lambda_0/8)n)$. The idea going forward is to control the right-hand-side of \ref{eq:variance_term_limsup} by separating the inner $t$-integration range into a small region between $s$ and $s+\varepsilon_n$ and between $s+\varepsilon_n$ and $\alpha(n+1)$. The integral in the first region is controlled by the small measure of $[s,s+\varepsilon_n]$ and the second region is controlled by the term $P_{J(s)\to J(t)}(t-s)$. The contribution of \ref{eq:limsup_summand_1} to the integral \ref{eq:variance_term_limsup} is
  \begin{align}
    &C_{x,k}\int_{\alpha n}^{\alpha(n+1)}\int_{s}^{\alpha(n+1)}J(t)J(s)\mu_n(t)^k\mu_n(s)P_{J(s)\to J(t)}(t-s)\,dt\,ds\nonumber\\
    &=C_{x,k}\int_{\alpha n}^{\alpha(n+1)}J(s)\mu_n(s)\int_{s}^{\alpha(n+1)}J(t)\mu_n(t)^kP_{J(s)\to J(t)}(t-s)\,dt\,ds\nonumber\\
    &=C_{x,k}\int_{\alpha n}^{\alpha(n+1)}J(s)\mu_n(s)\int_{s}^{s+\varepsilon_n}J(t)\mu_n(t)^kP_{J(s)\to J(t)}(t-s)\,dt\,ds\label{eq:contribution_integral_term_1}\\
    &+C_{x,k}\int_{\alpha n}^{\alpha(n+1)}J(s)\mu_n(s)\int_{s+\varepsilon_n}^{\alpha(n+1)}J(t)\mu_n(t)^kP_{J(s)\to J(t)}(t-s)\,dt\,ds.\label{eq:contribution_integral_term_2}
  \end{align}
  The first term (\ref{eq:contribution_integral_term_1}) can be controlled since the range of integration is small:
  \begin{align}
    \ref{eq:contribution_integral_term_1}
    &\leq C_{x,k}\int_{\alpha n}^{\alpha(n+1)}\mu_n(s)^{k+1}J(s)^2\varepsilon_n\,ds\nonumber\\
    &\leq C_{x,k}\mathbb{E}\left[L_{[\alpha n,\alpha (n+1)),[n,n+1)}^{k}|\mathcal{F}_\infty\right] J(\alpha(n+1))\varepsilon_n.\label{eq:contribution_integral_term_1_cont}
  \end{align}
  Now, note $\frac{\lambda_0\lambda_1}{\lambda_0+\lambda_1}\alpha=\lambda_0+\mathcal{O}_{\varepsilon\to 0^+}(\varepsilon)$ so that for sufficiently small $\varepsilon>0$, $\frac{\lambda_0\lambda_1}{\lambda_0+\lambda_1}\alpha-\frac{3\lambda_0}{8}<\ref{eq:exp_minimum_growth}$. Therefore for such small $\varepsilon>0$:
  \begin{align*}
    \ref{eq:contribution_integral_term_1_cont}\llas{n}{\infty}\mathbb{E}\left[L_{[\alpha n,\alpha (n+1)),[n,n+1)}^{k}|\mathcal{F}_\infty\right]^2.
  \end{align*}

  For the second term (\ref{eq:contribution_integral_term_2}), we use the Chernoff bound on $P_{J(s)\to J(t)}(t-s)$ for $s\in [\alpha n,\alpha(n+1))$ and $t\in [s+\varepsilon_n, \alpha(n+1))$ devised in Lemma \ref{lem:MGF_bound}:
  \begin{align*}
    P_{J(s)\to J(t)}(t-s)&\leq \mathbb{P}(Z_1^{(1)}(t-s)\leq J(t)|Z_1^{(1)}(0)=J(s))\leq\left(e^{\beta J(t)/J(s)}G(\beta,t-s)\right)^{J(s)}\\
    &\leq C e^{-KJ(s)(t-s)^2}\leq Ce^{-KJ(\alpha n)\varepsilon_n^2}\lsimas{n}{\infty} e^{-Ke^{\frac{\lambda_0}{8}n}}
  \end{align*}
  for some $K>0$ and all small $\varepsilon>0$. Therefore:
  \[
  \int_{\alpha n}^{\alpha(n+1)}J(s)\mu_n(s)\int_{s+\varepsilon_n}^{\alpha(n+1)}J(t)\mu_n(t)^kP_{J(s)\to J(t)}(t-s)\,dt\,ds\llas{n}{\infty} \mathbb{E}\left[L_{[\alpha n,\alpha (n+1)),[n,n+1)}^{k}|\mathcal{F}_\infty\right]^2.
  \]
almost surely on $\Omega_0^\infty$. Term \ref{eq:limsup_summand_2} is handled exactly analogously.

  For term \ref{eq:term_limsup_bound_2}:
  \begin{align*}
    \mathbb{P}(X_{3}^{s,t}\geq k|\mathcal{F}_\infty)\leq C_k(\mu^{(3)}(s,t))^k=C_k\mu_n(s)^kP_{J(s)\to J(t)}(t-s)^k.
  \end{align*}
  Thus the contribution of this term to the integral \ref{eq:variance_term_limsup} is precisely the same as for \ref{eq:sum_limsup_bound}.

  Therefore in total,
  \begin{align*}
    \mathbb{V}\left[L_{[\alpha n,\alpha (n+1)),[n,n+1)}^k|\mathcal{F}_\infty\right]&\llas{n}{\infty}\mathbb{E}[L_{[\alpha n,\alpha (n+1)),[n,n+1)}^k|\mathcal{F}_\infty]^2,
  \end{align*}
  so \ref{eq:tight_limsup_to_show} holds.
\end{proof}

Now, we end by proving the corollary.
\SFSSparsityCor*
\begin{proof}
    \textbf{1.} Recall that in Theorem \ref{thm:SFSSparsity}, we proved the fixed time versions of the above series of equalities. Note that $\tau_n$ is a discrete subsequence of time so that on $\Omega_0^\infty$:
    \begin{align*}
        \frac{\lambda_0\lambda_1}{\lambda_0+\lambda_1}\leq \lim_{k\to\infty}\liminf_{n\to\infty}\frac{\log j_k^-(\tau_n)}{\tau_n},\, \lim_{k\to\infty}\limsup_{n\to\infty}\frac{\log j_k^+(\tau_n)}{\tau_n}\leq  \frac{\lambda_0\lambda_1}{\lambda_0+\lambda_1}.
    \end{align*}
    Once again, since $j_k^-(\tau_n)\leq j_k^+(\tau_n)+1$ and $\tau_n\to\infty$, we immediately conclude
    \begin{align*}
        &\lim_{k\to\infty}\liminf_{n\to\infty}\frac{\log j_k^-(\tau_n)}{\tau_n}=\lim_{k\to\infty}\limsup_{n\to\infty}\frac{\log j_k^-(\tau_n)}{\tau_n}\\
        =&\lim_{k\to\infty}\liminf_{n\to\infty}\frac{\log j_k^+(\tau_n)}{\tau_n}=\lim_{k\to\infty}\limsup_{n\to\infty}\frac{\log j_k^+(\tau_n)}{\tau_n}=\frac{\lambda_0\lambda_1}{\lambda_0+\lambda_1}.
    \end{align*}
Combining this with Lemma \ref{lem:large_size_time}, which implies that $\tau_n\sim \frac{1}{\lambda_1}\log n$, proves the result.

    \textbf{2.} Remark that $\liminf_{t\to\infty}S_{j(t)}(t)\leq \liminf_{n\to\infty}S_{j(\tau_n)}(\tau_n)$, so to show the first line, it suffices to prove $\liminf_{t\to\infty}S_{j(t)}(t)=\infty$. Note that $\liminf_{t\to\infty}\frac{\log j_k^-(t)}{t}$ is decreasing in $k$ so that Theorem \ref{thm:SFSSparsity} implies that for all $k\geq 1$, $\liminf_{t\to\infty}\frac{\log j_k^-(t)}{t}\geq \frac{\lambda_0\lambda_1}{\lambda_0+\lambda_1}$. Since $\limsup_{t\to\infty}\frac{\log j(t)}{t}\leq \frac{\lambda_0\lambda_1}{\lambda_0+\lambda_1}-\varepsilon$, there is some time $T$ such that for all $t\geq T$, $j(t)<j_k^-(t)$. Then by definition $S_{j(t)}(t)>k$ for $t\geq T$, i.e., $\liminf_{t\to\infty}S_{j(t)}(t)\geq k$ and since $k$ is arbitrary, the first line is proven.

    For the second line, let $k_\varepsilon=\lceil c_2/c_1^{\varepsilon/2}+2\rceil$. Then as per Proposition \ref{prop:intermediate_j_k_plus}:
    \[\limsup_{t\to\infty}\frac{\log(j_k^+(t))}{t}\leq \frac{\lambda_0\lambda_1}{\lambda_0+\lambda_1}+\frac{\varepsilon}{2}.\]
    Thus, there is a $T(\omega)$ such that $t\geq T(\omega)$ implies $J(t)>j_k^+(t)$. For such $t$ by definition of $j_k^+(t)$, $S_{J(t)}(t)\leq k_\varepsilon\lesssim\varepsilon^{-1}$ as $\varepsilon\to 0^+$. And, once again the fixed detection-size limit follows from the detection-time limit.

    The last line follows from (1) of Theorem \ref{thm:SFSSparsity}. In particular, it implies that there is a subsequence of times $t_n\to \infty$ such that $S_{j(t_n)}(t_n)\to 0$.

    The last part of the corollary is just Proposition \ref{prop:limsup_theta}.
\end{proof}

\subsection{Tail of SFS at fixed detection size}
\subsubsection{Mass near $j=n$}\label{app:tau_n_SFS}
Here we prove Lemma \ref{lem:tau_n_SFS}:
\TauNSFS*
\begin{proof}

On $\Omega_1^\infty$, $S_n(\tau_n)$ has (along with being well-defined) support $\{0,1\}$ for $n\geq 3$. Let $A_1$ be the event that there is only 1 established type-1 clone. With probability 1 on $\Omega_1^\infty$, $\omega\in A_1$ iff eventually all cells belong to a single type-1 clone iff $S_n(\tau_n)=1$ for all large enough $n$. This proves the a.s. convergence. Note that the success probability is:
    \[\mathbb{P}(A_1|\Omega_1^\infty)=\mathbb{P}(A_1,\left(\Omega_0^\infty\right)^c|\Omega_1^\infty)=\frac{\mathbb{P}(A_1,(\Omega_0^\infty)^c)}{\pi_1}=\frac{\mathbb{P}(A_1|(\Omega_0^\infty)^c)(1-\pi_0)}{\pi_1}.\]

    

  We can compute $\mathbb{P}(A_1|(\Omega_0^{\infty})^c)$. Remark that on extinction of a supercritical one-type continuous-time linear birth-death process (with birth rate $a(1-v)$ and death rate $d+\mu$), the process becomes a subcritical continuous-time linear birth-death process with birth rate $\alpha=d+\mu$ and death rate $\beta=a(1-v)$. Thus, by using the description of our two-type model presented in Appendix Section \ref{app:alt_representation_model}, conditional on $\Omega_0^\infty$, the type-0 process evolves with birth parameter $\alpha$ and death parameter $\beta$. Each death there is an independent $\frac{q_1\mu}{d+\mu}$ probability of forming an established type-1 lineage produced in a replication-free manner (which is a further $q_1$-thinning of $\{T_i''\}$). And, further conditional on $\mathcal{F}_\infty$ and on $\Omega_0^\infty$, replication-dependent established type-1 lineages are formed with Poisson rate $q_1avZ_0(s)\,ds$ (which is a further $q_1$-thinning of $\{T_i'\}$). Thus, letting $D$ count the total number of death events that occur to the type-0 process:
  \begin{align*}
    &\mathbb{P}(A_1|(\Omega_0^{\infty})^c)\\
    &=\mathbb{E}\left[\mathbb{P}(A_1|\mathcal{F}_\infty,\left(\Omega_0^\infty\right)^c)|\left(\Omega_0^\infty\right)^c\right]\\
    &=\mathbb{E}\left[{D\choose 1}\frac{q_1\mu}{d+\mu}\left(\frac{d+p_1\mu}{d+\mu}\right)^{D-1}e^{-avq_1\int_0^\infty Z_0(s)\,ds}+\left(\frac{d+p_1\mu}{d+\mu}\right)^D\left(avq_1\int_0^\infty Z_0(s)\,ds\right)e^{-avq_1\int_0^\infty Z_0(s)\,ds}|(\Omega_0^\infty)^c\right].
  \end{align*}
  Let $Z(s)$ be a subcritical continuous-time linear birth-death process with rates $\alpha,\beta$, and again $D$ count the total number of death events. Now define $G(z,\theta)=\mathbb{E}\left[z^De^{-\theta \int_0^\infty Z(s)\,ds}\right]$. Hence, we can rewrite the above via:
  \begin{align}
      =\frac{q_1\mu}{d+\mu}\partial_z G((d+p_1\mu)/(d+\mu),avq_1)-avq_1\partial_\theta G((d+p_1\mu)/(d+\mu),avq_1)\label{eq:P_expr_in_G}.
  \end{align}

  By conditioning on the first event and using independence of lineages, and letting $\tau$ denote the time to the first event (which is exponential with rate $\alpha+\beta$):
  \[G(z,\theta)=\frac{\beta}{\alpha+\beta}z\mathbb{E}[e^{-\theta\tau}]+\frac{\alpha}{\alpha+\beta}\mathbb{E}[e^{-\theta\tau}]\left(G(z,\theta)\right)^2=\frac{\beta}{\alpha+\beta}z\frac{\alpha+\beta}{\alpha+\beta+\theta}+\frac{\alpha}{\alpha+\beta}\frac{\alpha+\beta}{\alpha+\beta+\theta}\left(G(z,\theta)\right)^2\]
  We then solve a quadratic equation to get $G$:
  \begin{align}
      G(z,\theta)=\frac{\alpha+\beta+\theta-\sqrt{(\alpha+\beta+\theta)^2-4\alpha\beta z}}{2\alpha}.\label{eq:G_expr}
  \end{align}
Note that the extinction probability of the type-0 process is that of a one-type process with birth rate $a(1-v)$ and death rate $d+\mu$ which is:
\begin{align}
    1-\pi_0=(d+\mu)/(a(1-v)).\label{eq:pi_0_expr}
\end{align}

We may then calculate $\pi_1$ by conditioning on the first event:
\[1-\pi_1=\frac{a(1-v)}{a+d+\mu}(1-\pi_1)^2+\frac{d}{a+d+\mu}+\frac{av}{a+d+\mu}p_1(1-\pi_1)+\frac{\mu}{a+d+\mu}p_1\]
Let the associated polynomial in $x=1-\pi_1$ be:
\[\mathcal{Q}(x)=a(1-v)x^2+d+avp_1x+\mu p_1-x(a+d+\mu).\]
Notice that $\mathcal{Q}''>0$ and $\mathcal{Q}(0)>0$ and $\mathcal{Q}(1)<0$, hence we are looking for the smaller root of $\mathcal{Q}$:
\begin{align}
    1-\pi_1=\frac{a+d+\mu-av p_1-\sqrt{(a+d+\mu-avp_1)^2-4a(1-v)(d+\mu p_1)}}{2a(1-v)}.\label{eq:pi_1_expr}
\end{align}

Thus, combining expressions \ref{eq:P_expr_in_G} to \ref{eq:pi_1_expr} gives an explicit expression for $\mathbb{P}(A_1|(\Omega_0^\infty)^c)$ and thus of $\mathbb{P}(A_1|\Omega_1^\infty)$.

\end{proof}


\subsubsection{Mass at tail in large selection limit}\label{app:freeze}

\Freeze*
\begin{proof}
    Decomposing the size at $Z_0(T_1)$:
    \begin{align*}
        &\mathbb{P}(\tau_{n(b)}<\infty,Z_0(\tau_{n(b)})=Z_0(T_1)|T_1\leq \tau_{n(b)}^{(0)},T_1<\infty)\\
        =&\sum_{r=0}^{n(b)-1}\mathbb{P}(\tau_{n(b)}<\infty,Z_0(\tau_{n(b)})=r|T_1\leq \tau_{n(b)}^{(0)},T_1<\infty,Z_0(T_1)=r)\mathbb{P}(Z_0(T_1)=r|T_1\leq \tau_{n(b)}^{(0)},T_1<\infty).
    \end{align*}
    Then, using the Strong-Markov property to start the process at time $T_1$:
    \begin{align}
        &=\sum_{r=0}^{n(b)-1}\mathbb{P}(\tau_{n(b)}<\infty,Z_0(\tau_{n(b)})=r|(Z_0(0)=r,Z_1^{(1)}(0)=1))\mathbb{P}(Z_0(T_1)=r|T_1\leq \tau_{n(b)}^{(0)},T_1<\infty).\label{eq:decomp_equation}
    \end{align}
    We will take $b\to\infty$. To understand how \ref{eq:decomp_equation} behaves, first notice that the summands are upper bounded by
    \begin{align*}
        &\mathbf{1}_{r\leq n(b)-1}\mathbb{P}(\tau_{n(b)}<\infty,Z_0(\tau_{n(b)})=r|(Z_0(0)=r,Z_1^{(1)}(0)=1))\mathbb{P}(Z_0(T_1)=r|T_1\leq \tau_{n(b)}^{(0)},T_1<\infty)\\
        \leq &\mathbf{1}_{r\leq n(b)-1}\frac{1}{\mathbb{P}(T_1\leq \tau_{n(b)}^{(0)}|T_1<\infty)}\mathbb{P}(Z_0(T_1)=r|T_1<\infty)\leq \frac{1}{\mathbb{P}(T_1\leq \tau_2^{(0)}|T_1<\infty)}\mathbb{P}(Z_0(T_1)=r|T_1<\infty)
    \end{align*}
    which is summable. Since $n(b)$ is monotone-increasing, $\mathbb{P}(Z_0(T_1)=r|T_1\leq \tau_{n(b)}^{(0)},T_1<\infty)\to \mathbb{P}(Z_0(T_1)=r|T_1\leq \tau_{n(\infty)}^{(0)},T_1<\infty)$. We can thus show convergence of \ref{eq:decomp_equation} by showing pointwise convergence of the first factor of each summand and invoking dominated convergence.
    
    First consider $\log n(b)\ll b$ as $b\to\infty$. We can write:
    \[\mathbb{P}(\tau_{n(b)}<\infty,Z_0(\tau_{n(b)})=r|(Z_0(0)=r,Z_1^{(1)}(0)=1))= \mathbb{P}(Z_0(\tau_{n(b)})=r|(Z_0(0)=r,Z_1^{(1)}(0)=1))\]
    if we interpret $Z_0(\infty)\neq r$. Let $E_0\sim \text{Exp}(r(a(1-v)+d+\mu))$ denote the time to the first type-0 event and $\tau_n^{(1),1}$ the time for the starting type-1 clone to reach size $n$. Remark that $E_0$ and $\tau_n^{(1),1}$ are independent. Also, $Z_1^{(1)}(t)\geq n$ implies that $\tau_n^{(1),1}\leq t$. Finally, define an intermediate time $t_b=(\log (b\cdot n(b)))/\lambda_1$. So, we obtain the further lower bound:
    \begin{align*}
        &\mathbb{P}(Z_0(\tau_{n(b)})=r|(Z_0(0)=r,Z_1^{(1)}(0)=1))\\
        &\geq \mathbb{P}(E_0>\tau_{n(b)-r}^{(1),1}|(Z_0(0)=r,Z_1^{(1)}(0)=1))\\
        &\geq \mathbb{P}(E_0>t_b|(Z_0(0)=r,Z_1^{(1)}(0)=1))\mathbb{P}(\tau_{n(b)-r}^{(1),1}<t_b|(Z_0(0)=r,Z_1^{(1)}(0)=1))\\
        &\geq \mathbb{P}(E_0>t_b|(Z_0(0)=r,Z_1^{(1)}(0)=1))\mathbb{P}(Z_1^{(1)}(t_b)\geq n(b)-r|(Z_0(0)=r,Z_1^{(1)}(0)=1)).
    \end{align*}
    Both terms in the above lower bound go to $1$. The first probability does trivially since $t_b\to 0$ as $b\to\infty$:
    \[\mathbb{P}(E_0>t_b|(Z_0(0)=r,Z_1^{(1)}(0)=1))\to 1.\]
    The second term is $1$ whenever $n(b)-r\leq 0$ so assume $n(b)-r\geq 1$. Using Bernoulli's inequality $(1-x)^n\geq 1-nx$, the CDF of $Z_1^{(1)}$, and that $\exp(\lambda t_b)=bn(b)$, we obtain:
    \[\mathbb{P}(Z_1^{(1)}(t_b)\geq n(b)-r|(Z_0(0)=r,Z_1^{(1)}(0)=1))\geq \frac{\lambda_1 bn(b)}{(a+b)bn(b)-d}\cdot \left(1-\frac{\lambda_1(n(b)-r)}{(a+b)bn(b)-d}\right)\to 1\]
    as $b\to\infty$. 
    
    Now, suppose $\log n(b)\gg b$ as $b\to\infty$ and rewrite \ref{eq:decomp_equation} as
    \begin{align*}
        &\sum_{r=0}^{n(b)-1}\mathbb{P}(\tau_{n(b)}<\infty,Z_0(\tau_{n(b)})=r|(Z_0(0)=r,Z_1^{(1)}(0)=1))\mathbb{P}(Z_0(T_1)=r|T_1\leq \tau_{n(b)}^{(0)},T_1<\infty)\\
        &=\mathbb{P}(\tau_{n(b)}<\infty|(Z_0(0)=0,Z_1^{(1)}(0)=1))\cdot \mathbb{P}(Z_0(T_1)=0|T_1\leq \tau_{n(b)}^{(0)},T_1<\infty)\\
        &+\sum_{r=1}^{n(b)-1}\mathbb{P}(\tau_{n(b)}<\infty,Z_0(\tau_{n(b)})=r|(Z_0(0)=r,Z_1^{(1)}(0)=1))\mathbb{P}(Z_0(T_1)=r|T_1\leq \tau_{n(b)}^{(0)},T_1<\infty).
    \end{align*}
    We will now show that the first term in on the RHS of the previous display converges. Note that $\mathbb{P}(\Omega_1^{\infty}|(Z_0(0)=r,Z_1^{(1)}(0)=1))\to 1$ as $b\to\infty$ and $\inf_n\mathbb{P}(\tau_{n}<\infty|\Omega_1^\infty)=\mathbb{P}(\cap \{\tau_n<\infty\}|\Omega_1^\infty)=1$. Thus:
    \[
    \mathbb{P}(\tau_{n(b)}<\infty|Z_1^{(1)}(0)=1)\geq \mathbb{P}(\tau_{n(b)}<\infty|Z_1^{(1)}(0)=1,\Omega_1^\infty)P(\Omega_1^\infty|Z_1^{(1)}(0)=1)\to 1
    \]
    as $b\to\infty$. Let $L_r=\sup\{t:Z_0(t)=r\}$ denote the last time the type-0 population returns to size $r$.  If $r>0$, since $Z_0$ converges almost surely to $0$ or $\infty$, $L_r$ is an almost surely finite random variable. Now, supposing $r\geq 1$:
    \begin{align*}
        &\mathbb{P}(\tau_{n(b)}<\infty,Z_0(\tau_{n(b)})=r|(Z_0(0)=r,Z_1^{(1)}(0)=1))\\
        &=\mathbb{P}(Z_0(\tau_{n(b)})=r|(Z_0(0)=r,Z_1^{(1)}(0)=1))\\
        &\leq \mathbb{P}(\tau_{n(b)-r}^{1}\leq L_r|(Z_0(0)=r,Z_1^{(1)}(0)=1))
    \end{align*}
    with the last inequality since $Z_0(\tau_{n(b)})=r$ implies that $Z_1(\tau_{n(b)})=n(b)-r$, so that $\tau^1_{n(b)-r}\leq \tau_{n(b)}\leq L_r$. Fix $\varepsilon>0$. There is $M_\varepsilon>0$ large enough, and independent of $b$, so that $\mathbb{P}(L_r\geq M_\varepsilon|(Z_0(0)=r,Z_1^{(1)}(0)=1))\leq \varepsilon$. By partitioning the probability in the above upper bound on the sets $\{L_r\geq M_\varepsilon\},\{L_r<M_\varepsilon\}$ and by writing the event $\tau^1_{n(b)-r}\leq L_r$ as $\sup_{s\in [0,L_r]}Z_1(s)\geq n(b)-r$:
    \[\leq \varepsilon+\mathbb{P}(\sup_{s\in [0,M_\varepsilon]}Z_1(s)\geq n(b)-r|(Z_0(0)=r,Z_1^{(1)}(0)=1)).\]
    Let $B_1(s)$ denote the total number of type-1 birth-events which have occurred up to time $s$. Then, $B_1(s)$ is increasing and $B_1(s)\geq Z_1(s)$, so that:
    \[\leq \varepsilon+\mathbb{P}(B_1(M_\varepsilon)\geq n(b)-r|(Z_0(0)=r,Z_1^{(1)}(0)=1))\leq \varepsilon+\frac{\mathbb{E}[B_1(M_\varepsilon)|(Z_0(0),Z_1^{(1)}(0)=1)]}{n(b)-r}.\]
    So, we must show that the expectation is $o(n(b))$ to conclude with the proof. Let $B_1^i(s-T_i)$ denote the number of birth events associated with clone $i$. Remark that the mean intensity of $B_1^1(s)$ is $(a+b)\mathbb{E}[Z_1^1(s)]\,ds$ so that $\mathbb{E}[B_1^1(s)]=\frac{a+b}{\lambda_1}\left(e^{\lambda_1 s}-1\right)$. Thus, writing $B_1(s)=\sum_{T_i\leq s}B_1^i(s-T_i)$, we obtain by conditioning on the mutation times $\{T_i\}_i$ and using Campbell's theorem:
    \[\mathbb{E}[B_1(M_\varepsilon)]=\frac{a+b}{\lambda_1}\int_0^{M_\varepsilon}\left(e^{\lambda_1 (M_\varepsilon-t)}-1\right)(av+\mu)e^{\lambda_0 t}\,dt\ll n(b).\]
\end{proof}
\subsection{Scaling of large families}\label{app:scaling}
Here we prove Proposition \ref{prop:scaling}.
\Scaling*
\begin{proof}
    \textbf{Part 1: }First we show:
    \[\lim_{t\to\infty}\mathbb{E}\left[\widetilde{R}_x(t)\right]=\lim_{t\to\infty}\mathbb{E}\left[\sum_{j\geq xg(t)}S_j(t)\right]=\frac{(av+\mu) q_1^{1-\lambda_0/\lambda_1}x^{-\lambda_0/\lambda_1}}{\lambda_1}\Gamma(\frac{\lambda_0}{\lambda_1},q_1 x)\]
    By Fubini-Tonnelli applied twice:
  \begin{align*}
    \lim_{t\to\infty}\mathbb{E}[\widetilde{R}_x(t)]&=(av+\mu)\lim_{t\to\infty}\sum_{j\geq xe^{\lambda_1 t}}\int_0^t e^{\lambda_0 s}\mathbb{P}(Z_1^{(1)}(t-s)=j)\,ds\\
    &=(av+\mu)\lim_{t\to\infty}\int_0^\infty \mathbf{1}_{s\leq t} e^{\lambda_0 s}\mathbb{P}(Z_1^{(1)}(t-s)\geq xe^{\lambda_1 t})\,ds
  \end{align*}

Notice pointwise that $e^{\lambda_0 s} \mathbb{P}(Z_1^{(1)}(t-s)\geq x e^{\lambda_1 t})=e^{\lambda_0 s} \mathbb{P}(e^{-\lambda_1 (t-s)}Z_1^{(1)}(t-s)\geq x e^{\lambda_1 s})\to e^{\lambda_0 s}\mathbb{P}(Y_1^{(1)}\geq x e^{\lambda_1 s})$. Furthermore by Markov inequality:

\[\mathbb{P}(Z_1^{(1)}(t-s)\geq xe^{\lambda_1 t})\leq \frac{e^{\lambda_1(t-s)}}{x e^{\lambda_1 t}}=\frac{1}{x}e^{-\lambda_1 s}.\]
Next note that $\int_0^\infty \frac{1}{x}e^{\lambda_0 s}e^{-\lambda_1 s}<\infty$, so by dominated convergence, we may pass the limit inside:
\[
=(av+\mu)\int_0^\infty e^{\lambda_0 s}(0\cdot p_1+\exp(-q_1 xe^{\lambda_1 s})\cdot q_1)\,ds.
\]
We then use the substitution $u=e^{\lambda_1 s}$:
\[
=\frac{(av+\mu)q_1}{\lambda_1}\int_1^\infty u^{\lambda_0/\lambda_1-1}e^{-q_1 x u}\,du=\frac{(av+\mu) q_1^{1-\lambda_0/\lambda_1}x^{-\lambda_0/\lambda_1}}{\lambda_1}\Gamma(\frac{\lambda_0}{\lambda_1},q_1 x).
\]
This proves the limit in the deterministic fitness advance case.

\textbf{Part 2: }Now, suppose instead that we consider the case with the random fitness advance.

Recall that $\lim_{t\to\infty}e^{-\lambda_0t}\mathbb{E}[S_j^B(t)]=(av+\mu)\int_0^\infty e^{-\lambda_0 s}\mathbb{E}[P_j(s|B)]\,ds$. Note then:

  \begin{align*}
    e^{\lambda_0 t}\left|\mathbb{E}\left[\sum_{j\geq xg_2(t)}\left(e^{-\lambda_0 t}S_j^B(t)\right)\right]-\sum_{j\geq xg_2(t)}\lim_{t\to\infty}\mathbb{E}\left[ e^{-\lambda_0 t}S_j^B(t)\right]\right|&=e^{\lambda_0 t}\sum_{j\geq xg_2(t)}\int_t^\infty e^{-\lambda_0 s}\mathbb{E}[P_j(s|B)]\,ds\\
    &=e^{\lambda_0 t}\int_t^\infty e^{-\lambda_0 s}\mathbb{E}[\mathbb{P}(Z_1(s)\geq xg_2(t)|B)]\,ds.
  \end{align*}
Let $s=t+u$, then:
\begin{align*}
    &=\int_0^\infty e^{-\lambda_0 u}\mathbb{E}[\mathbb{P}(Z_1(t+u)\geq g_2(t)|B)]\,du.
  \end{align*}
Notice that with probability $1$ that $B<b_{\max}$, so by dominated convergence the above goes to $0$. Hence to show the desired limit converges it suffices to show that:
  \[
  \lim_{t\to\infty} e^{\lambda_0 t}\sum_{j\geq xg_2(t)}\lim_{s\to\infty}e^{-\lambda_0 s}\mathbb{E}[S_j^B(s)]
  \]
  converges. Note that from Proposition \ref{prop:SFS_mean_asymptotics}, for any $\varepsilon>0$ there is $J>0$ such that $j\geq J$ implies:
  \[
  (1-\varepsilon)\frac{C}{j^{\frac{\lambda_0}{\lambda_1(b_{\max})}+1}\log j}\leq \lim_{s\to\infty}e^{-\lambda_0s}\mathbb{E}[S_j^B(s)]\leq (1+\varepsilon)\frac{C}{j^{\frac{\lambda_0}{\lambda_1(b_{\max})}+1}\log j}.
  \]
  Taking $T$ large enough such that $t\geq T$ implies $xg_2(t)\geq J$, in which case:
  \[\left|e^{\lambda_0 t}\sum_{j\geq xg_2(t)}\lim_{s\to\infty}e^{-\lambda_0 s}\mathbb{E}[S_j^B(s)]-e^{\lambda_0 t}\sum_{j\geq xg_2(t)} \frac{C}{j^{\frac{\lambda_0}{\lambda_1(b_{\max})}+1}\log j}\right|\leq \varepsilon e^{\lambda_0 t}\sum_{j\geq g_2(t)}\frac{C}{j^{\frac{\lambda_0}{\lambda_1(b_{\max})}+1}\log j}.\]
  
  We thus see that it suffices to examine convergence of:
  \[
  \lim_{t\to\infty} e^{\lambda_0 t}\sum_{j\geq x g_2(t)}\frac{C}{j^{\frac{\lambda_0}{\lambda_1(b_{\max})}+1}\log j}.
  \]
  Notice that $j\to \frac{1}{j^{\frac{\lambda_0}{\lambda_1(b_{\max})}+1}\log j}$ is monotonically decreasing in $j$, so by integral comparisons it suffices to examine convergence of:
  \begin{align*}
    \lim_{t\to\infty} e^{\lambda_0 t}\int_{xg_2(t)}^\infty\frac{C}{z^{\frac{\lambda_0}{\lambda_1(b_{\max})}+1}\log z}\,dz.
  \end{align*}
We can rewrite the integral in the previous display in terms of the exponential integral function as
\[
\int_{xg_2(t)}^\infty\frac{C}{z^{\frac{\lambda_0}{\lambda_1(b_{\max})}+1}\log z}\,dz=E_1\left(\frac{\lambda_0}{\lambda_1(b_{\max})}\log (xg_2(t))\right)
\]
where $E_1(t)$ is the exponential integral. Finally utilize the asymptotic, $E_1(t)\sim e^{-t}/t$ as $t\to\infty$ to conclude the result.
\end{proof}
\subsection{Relative SFS time limit}\label{app:frequency_limit}
Here we prove Proposition \ref{prop:frequency_limit} and Corollary \ref{cor:scale_sum_SFS}. But first we prove an important lemma:
\begin{lemma}
     Let $Y^{(i)}(t)=e^{-\lambda_1(t-T_i)}Z_1^{(i)}(t-T_i)$ and $Y_*^{(i)}=\sup_{T_i\leq t<\infty} Y^{(i)}(t)$. Then, $\lim_{i\to\infty} e^{-\lambda_1 T_i}Y_*^{(i)}=0$ a.s. If $i\geq \lim_{t\to\infty}M_t$, set for notation $Y^{(i)}(t)\equiv 0$\label{lem:Y_small}.
\end{lemma}

\begin{proof}
    It suffices to show that for all $\varepsilon>0$ there is an $I$ for which $i\geq I$ implies $e^{-\lambda_1 T_i}Y_*^{(i)}\leq \varepsilon$ (a.s.).

    Define $A_{i,\varepsilon}=\{e^{-\lambda_1 T_i}Y_*^{(i)}>\varepsilon\}$. By Borel-Cantelli, it suffices to show that for each $\varepsilon>0$, $\sum_{i=1}^\infty \mathbb{P}(A_{i,\varepsilon})<\infty$. By Doob's Martingale inequality:
\[
\mathbb{P}(A_{i,\varepsilon})=\mathbb{E}[\mathbb{P}(Y_*^{(i)}>\varepsilon e^{\lambda_1 T_i}|T_i)]\leq \frac{1}{\varepsilon}\mathbb{E}[e^{-\lambda_1 T_i}].
\]
It thus suffices to show that $\sum_{i=1}^\infty \mathbb{E}\left[e^{-\lambda_1 T_i}\right]<\infty$, which is immediate from Campbell's theorem.
\end{proof}

\FrequencyLimit*
\begin{proof}
  Conditional on $\Omega_0^\infty$, there a.s exists a $\varepsilon>0$ such that $Wf>\varepsilon>0$. But from Lemma \ref{lem:Y_small}, for $i>I$ sufficiently large, $e^{-\lambda_1 T_i}Y_1^{(i)}<\varepsilon $, thus the proposed limit is a.s. finite.
  
  Recall that:
  \[
  R_f(t)=\sum_{1\leq i\leq M_t} \mathbf{1}_{Z_1^{(i)}(t-T_i)>Z_1(t)f}=\sum_{1\leq i\leq M_t} \mathbf{1}_{e^{-\lambda_1 T_i}e^{-\lambda_1(t-T_i)}Z_1^{(i)}(t-T_i)>e^{-\lambda_1 t}Z_1(t)f}.
  \]

  Let $\delta=\varepsilon/2$ and note there exists a time $T$ such that for all $t>T$, $|e^{-\lambda_1 t}Z_1(t)f-Wf|\leq \delta$. Without loss of generality, consider $T$ even larger such that $i>M_T$ implies that $|e^{-\lambda_1 T_i}Y_*^{(i)}|\leq \delta$, so:
  \begin{align*}
    \lim_{t\to\infty} R_f(t)&=\lim_{t\to\infty} \sum_{1\leq i\leq M_T} \mathbf{1}_{Z_1^{(i)}(t-T_i)>Z_1(t)f}+\sum_{M_T< i\leq M_t} \mathbf{1}_{Z_1^{(i)}(t-T_i)>Z_1(t)f}\\
    &=\sum_{1\leq i\leq M_T}\mathbf{1}_{e^{-\lambda_1 T_i}Y_1^{(i)}>Wf}+0\\
    &=\sum_{i\in \mathbb{N}}\mathbf{1}_{e^{-\lambda_1 T_i}Y_1^{(i)}>Wf}.
  \end{align*}
In the second equality, we are justified in taking the limit within the indicator in the first sum as long as $e^{-\lambda_1 T_i}Y_1^{(i)}/W\neq f$ (which is a probability 1 event), from continuous mapping. An analogous proof holds for $\widetilde{R}_x$.
\end{proof}

\ScaleSumSFS*
\begin{proof}
  The same proof showing the convergence of $R_f(t)$ in Proposition \ref{prop:frequency_limit} also shows that $\sum_{j>f(Z_1(t)+Z_0(t))}S_j(t)$ converges to the same limit as $R_f$. Thus, the middle row limits hold.

  On $\Omega_0^\infty$, $W>0$. Furthermore, infinitely many established type-1 clones will be formed, say the first one is $Z_1^{(n_1)},Z_1^{(n_2)}$. Therefore:
  \[
  \frac{Z_1^{(n_1)}(t-T_{n_1})}{Z_1(t)}=e^{-\lambda_1 T_{n_1}}\frac{e^{-\lambda_1(t-T_{n_1})}Z_1^{(n_r)}(t-T_{n_1})}{e^{-\lambda_1 t}Z_1(t)}\overset{t\to\infty}{\to} L_1>0.
  \]
So, the largest clone makes up a fraction $<1-\varepsilon$ of the population for some $\varepsilon>0$ and all sufficiently large times $t>T$ (for example, with $\varepsilon=L_1/2$ say). Take $n>N$ such that $\tau_n>T$ and such that $k(n)/n<\varepsilon/2$ and thus $\sum_{j=n-k(n)}^n S_j(\tau_n)=0$. For the right upper limit, eventually $G^+(t)>Z_1(t)$ for all $t>T$.

  The lower left and right limits follow by similar arguments.
\end{proof}

\subsection{Relative SFS Small Frequency Behavior}\label{app:relSFS}

Here we prove Theorem \ref{prop:ASRelativeSFS}.

\ASRelativeSFS*

It is proved by breaking it up into a lemma and one corollary. Let $Z_x=\sum_{i\in \mathbb{N}}\mathbf{1}_{Y_1^{(i)}> xe^{\lambda_1 T_i}}$, we then show that $\lim_{x\to 0^+}x^{\lambda_0/\lambda_1}Z_x$ converges. Using the change of variables $x=e^{-r}$ and defining $Z(r)=Z_{e^{-r}}$, it suffices to show $\lim_{r\to \infty}e^{-\frac{\lambda_0}{\lambda_1}r}Z(r)$ converges.

\begin{lemma}
  On $\Omega_0^\infty$:
  \[\lim_{x\to0^+}\lim_{t\to\infty}x^{\lambda_0/\lambda_1}\widetilde{R}_x(t)=\frac{(av+\mu)q_1^{1-\lambda_0/\lambda_1}}{\lambda_1}\Gamma\left(\frac{\lambda_0}{\lambda_1}\right)Y\]
  \label{lem:small_x_behavior_approx}
\end{lemma}

\begin{proof}
  As per our discussion above, we show $e^{-\frac{\lambda_0}{\lambda_1}r}Z(r)$ converges to the above limit. Notice that $Z(r)$ is non-decreasing in $r$, so we aim to use the convergence tool directly. We use the same type of notation to define approximations of $Z(r)$ as in Table \ref{table:SFS_notation}. Following the same logic as Proposition \ref{prop:L2_driver} and using the intermediate approximation $e^{-\frac{\lambda_0}{\lambda_1}r}\mathbb{E}[Z(r)|\mathcal{B}_\infty]$, we get:

\[\mathbb{E}\left[\left(e^{-\frac{\lambda_0}{\lambda_1}r}Z(r)-e^{-\frac{\lambda_0}{\lambda_1}r}\overline{Z}(r)\right)^2\right]\lsimas{r}{\infty} e^{-\frac{\lambda_0}{\lambda_1}r},\]
where $\overline{Z}(r)=\int_0^\infty \mathbb{P}(Y_1^{(1)}\geq e^{-r}e^{\lambda_1 s})(av+\mu)Z_0(s)\,ds$.

Defining $\widehat{Z}(r)=\int_0^\infty \mathbb{P}(Y_1^{(1)}\geq e^{-r}e^{\lambda_1 s})(av+\mu)Ye^{\lambda_0 s}\,ds$, we have from above:
\[\widehat{Z}(r)=\frac{(av+\mu) q_1^{1-\lambda_0/\lambda_1}e^{\lambda_0/\lambda_1 r}}{\lambda_1}\Gamma(\frac{\lambda_0}{\lambda_1},q_1 e^{-\lambda_0/\lambda_1 r})Y.\]
Hence $e^{-\frac{\lambda_0}{\lambda_1}r}\widehat{Z}(r)\to \frac{(av+\mu) q_1^{1-\lambda_0/\lambda_1}}{\lambda_1}\Gamma(\frac{\lambda_0}{\lambda_1})Y$ as $r\to\infty$. We further have that $\mathbb{E}\left[\left(\widehat{Z}(r)-\overline{Z}(r)\right)^2\right]\lesssim \exp\left(\frac{\lambda_0}{\lambda_1}r\right)$ as $r\to\infty$.
Hence $\int_0^\infty e^{-\frac{\lambda_0}{\lambda_1}r}\mathbb{E}[|Z(r)-\widehat{Z}(r)|]<\infty$, so the limit follows from Proposition \ref{prop:convergence_tool}.
\end{proof}

\begin{corollary}
  On $\Omega_0^\infty$:
\[
\lim_{f\to0^+}\lim_{t\to\infty}f^{\lambda_0/\lambda_1}R_f(t)=\frac{(av+\mu)q_1^{1-\lambda_0/\lambda_1}}{\lambda_1}\Gamma\left(\frac{\lambda_0}{\lambda_1}\right)(W^{-\lambda_0/\lambda_1}Y).
\]
  
\end{corollary}

\begin{proof}
  Notice that $\lim_{t\to\infty}R_f(t)=\sum_{i\in \mathbb{N}}\mathbf{1}_{Y_1^{(i)}\geq fW e^{\lambda_1 T_i}}=Z_{fW}$ by Proposition \ref{prop:frequency_limit}. By the same argument as in Proposition \ref{prop:frequency_limit}, $\lim_{t\to\infty}\widetilde{R}_{fW}(t)=Z_{fW}$ a.s. Notice letting $x=fW$ and taking $x\to 0^+$, from Lemma \ref{lem:small_x_behavior_approx}:

  \[\lim_{f\to 0^+}(fW)^{\lambda_0/\lambda_1}Z_{fW}=\frac{(av+\mu)q_1^{1-\lambda_0/\lambda_1}}{\lambda_1}\Gamma\left(\frac{\lambda_0}{\lambda_1}\right)Y.\]
  Multiplying by $W^{-\lambda_0/\lambda_1}$ concludes the result.
\end{proof}
\subsection{Small mutation behavior of $(Y,W)$}\label{app:small_mutation_relative_SFS}
Here we prove Proposition \ref{prop:small_mutation_relative_SFS}. We first show that $(Y,v^{-\lambda_1/\lambda_0}W)$ converges in distribution as $v\to 0$.
\begin{proposition}
  As $v\to 0^+$, $(Y_v,v^{-\lambda_1/\lambda_0}W_v)|\Omega_{0,v}^\infty$ converges in distribution to a random vector with Laplace Transform:
  \begin{align}
    M_0(x,y)=\frac{q_{0,0}}{q_{0,0}+x+y^{\lambda_{0,0}/\lambda_1}\frac{aq_1^{1-\lambda_{0,0}/\lambda_1}}{\lambda_1}\frac{\pi}{\sin(\pi\lambda_{0,0}/\lambda_1)}}\label{eq:laplace_transform_limit}
  \end{align}
  The $v$-subscript in $Y_v,W_v,\Omega_{0,v}^\infty,q_{0,v},\lambda_{0,v}$ denotes the dependence on the mutation rate $v$.\label{prop:conv_dist}
\end{proposition}
\begin{proof}
  We also explicitly note wite $Z_{0,v}(t)=Z_0(t)$ to denote the type-0's dependence on $v$. When mutations are only attached to births, note that the mutation times $\{T_i\}=\Phi$ form a Cox process with respect to $\{Z_0(s);s\geq 0\}$. In particular, conditional on $\mathcal{F}_\infty$, $\widetilde{\Phi}=\sum_{1\leq i\leq \Phi(\mathbb{R}_{\geq 0})}\delta_{(T_i,Y_1^{(i)})}$ is a Poisson Point Process on $\mathbb{R}_{\geq 0}\times \mathbb{R}_{\geq 0}$. Let $f_y(t,\ell)=ye^{-\lambda_1 t}\ell$ and denote the distribution of $Y_1^{(i)}$ by $\mathbb{Y}$. Then:

\begin{align*}
  M(x,y)=\mathbb{E}\left[e^{-xY_v-yW_v}\right]&=\mathbb{E}[e^{-xY_v}\mathbb{E}[e^{-\int f_y(t,\ell)\widetilde{\Phi}(dt\times d\ell)}|\mathcal{F}_{\infty}]]\\
  &=\mathbb{E}[e^{-xY}e^{-\int_0^\infty \int_0^\infty (1-e^{-ye^{-\lambda_1 s}\ell})\,\mathbb{Y}(d\ell)avZ_{0,v}(s)\,ds}]\\
  &=\mathbb{E}[e^{-xY}e^{-\int_0^\infty q_1av (1-\frac{q_1}{q_1+ye^{-\lambda_1 s}})Z_{0,v}(s)\,ds}].
\end{align*}

First note that the exponent term in the expectation of $M(0,v^{-\lambda_1/\lambda_{0,v}}y)$ can be simplified as follows:

  \[
  \int_0^\infty q_1 av\left(1-\frac{q_1}{q_1+yv^{-\lambda_1/\lambda_{0,v}}e^{-\lambda_1 s}}\right)Z_{0,v}(s)\,ds=\int_0^\infty q_1 av\left(\frac{yv^{-\lambda_1/\lambda_{0,v}}}{q_1e^{\lambda_1 s}+yv^{-\lambda_1/\lambda_{0,v}}}\right)Z_{0,v}(s)\,ds,
  \]
letting $u=e^{\lambda_1 s}$ we then get
\begin{align*}
  &=\frac{q_1a}{\lambda_1}\int_{1}^\infty v\frac{yu^{-1}}{q_1 uv^{\lambda_1/\lambda_{0,v}}+y}Z_{0,v}\left(\frac{1}{\lambda_1}\log u\right)\,du
\end{align*}
we then make the substitution $r=uv^{\lambda_1/\lambda_{0,v}}$ to see
\begin{align}
  &=\frac{q_1 a}{\lambda_1}\int_{0}^\infty \mathbf{1}_{r\geq v^{\lambda_1/\lambda_{0,v}}}\frac{yr^{-1}}{q_1 r+y}vZ_{0,v}\left(\frac{1}{\lambda_1}\log\left(rv^{-\lambda_1/\lambda_{0,v}}\right)\right)\,dr\nonumber\\
  &=\frac{q_1 a}{\lambda_1}\int_{0}^\infty \mathbf{1}_{r\geq v^{\lambda_1/\lambda_{0,v}}}\frac{yr^{\frac{\lambda_{0,v}}{\lambda_1}-1}}{q_1 r+y}e^{-\lambda_0 t_v(r)}Z_{0,v}\left(t_v(r)\right)\,dr\nonumber
\end{align}
where $t_v(r)=\frac{1}{\lambda_1}\log\left(rv^{-\lambda_1/\lambda_{0,v}}\right)$. Now, we can approximate $e^{-\lambda_{0,v}t_v(r)}Z_{0,v}(t_v(r))$ via $Y_v$ in the above integral as follows:
\begin{align*}
    &\mathbb{E}\left[\left|\frac{q_1 a}{\lambda_1}\int_{0}^\infty \mathbf{1}_{r\geq v^{\lambda_1/\lambda_{0,v}}}\frac{yr^{\frac{\lambda_{0,v}}{\lambda_1}-1}}{q_1 r+y}\left(e^{-\lambda_0 t_v(r)}Z_{0,v}\left(t_v(r)\right)-Y_v\right)\,dr\right|\right]\\
    &\leq \frac{q_1 a}{\lambda_1}\int_{0}^\infty \mathbf{1}_{r\geq v^{\lambda_1/\lambda_{0,v}}}\frac{yr^{\frac{\lambda_{0,v}}{\lambda_1}-1}}{q_1 r+y}\mathbb{E}\left[\left|e^{-\lambda_0 t_v(r)}Z_{0,v}\left(t_v(r)\right)-Y_v\right|\right]\,dr\\
    &\leq \frac{q_1 a}{\lambda_1}\int_{0}^\infty \mathbf{1}_{r\geq v^{\lambda_1/\lambda_{0,v}}}\frac{yr^{\frac{\lambda_{0,v}}{\lambda_1}-1}}{q_1 r+y}\sqrt{\mathbb{E}\left[\left|e^{-\lambda_0 t_v(r)}Z_{0,v}\left(t_v(r)\right)-Y_v\right|^2\right]}\,dr\\
    &= \frac{q_1 a}{\lambda_1}\sqrt{\frac{a(1-v)+d+\mu}{\lambda_0}}\int_{0}^\infty \mathbf{1}_{r\geq v^{\lambda_1/\lambda_{0,v}}}\frac{yr^{\frac{\lambda_{0,v}}{\lambda_1}-1}}{q_1 r+y}e^{-\frac{\lambda_0}{2}t_v(r)}\,dr\\
    &\leq  v^{1/2}\frac{q_1 a}{\lambda_1}\sqrt{\frac{a(1-v)+d+\mu}{\lambda_0}}\int_{0}^\infty\frac{yr^{\frac{\lambda_{0,v}}{2\lambda_1}-1}}{q_1 r+y}\,dr\overset{v\to 0^+}{\to} 0.
\end{align*}
The equality follows from Lemma \ref{lem:L2_distance}. Remark that $Y_v$ is $0$ with probability $p_{0,v}$ and exponential with rate $q_{0,v}$ with probability $q_{0,v}$, so it is easy to deduce that $Y_v\Rightarrow Y_0$ in distribution as $v\to 0^+$. Thus:
\begin{align}
    -xY_v+\frac{q_1 aY_v}{\lambda_1}\int_{0}^\infty \mathbf{1}_{r\geq v^{\lambda_1/\lambda_{0,v}}}\frac{yr^{\frac{\lambda_{0,v}}{\lambda_1}-1}}{q_1 r+y}\,dr\Rightarrow -xY_0+\frac{q_1 a Y_0}{\lambda_1}\int_{0}^\infty \frac{yr^{\frac{\lambda_{0,0}}{\lambda_1}-1}}{q_1 r+y}\,dr.\label{eq:limit_laplace}
\end{align}
Thus by an application of Slutsky's theorem and remarking that for $z\geq 0$ and $X$ a non-negative random variable $e^{-zX}\leq 1$ is bounded, we obtain:
\begin{align}
  M(x,v^{-\lambda_1/\lambda_{0,v}}y)&\to \mathbb{E}\left[\exp\left(-Y_0x-Y_0\frac{q_1 a}{\lambda_1}\int_0^\infty\frac{yr^{\frac{\lambda_0}{\lambda_1}-1}}{q_1r+y}\,dr\right)\right]=\mathbb{E}\left[\exp\left(-Y_0x-Y_0\frac{q_1 a}{\lambda_1}\frac{y^{\lambda_0/\lambda_1}}{q_1^{\lambda_0/\lambda_1}}\frac{\pi}{\sin\left(\pi\lambda_0/\lambda_1\right)}\right)\right]\nonumber\\
  &=\frac{q_{0,0}}{q_0+x+y^{\lambda_{0,0}/\lambda_1}\frac{ aq_1 ^{1-\lambda_{0,0}/\lambda_1}}{\lambda_1}\frac{\pi}{\sin(\pi\lambda_{0,0}/\lambda_1)}}:=M_0(x,y).\nonumber
\end{align}
Also note that on $(\Omega_1^{\infty})^c$, $(Y,W)=(0,0)$ a.s. Hence:
\begin{align}
  \mathbb{E}[e^{-xY_v-yW_v}|\Omega_1^\infty]=\frac{1}{\mathbb{P}(\Omega_1^\infty)}\left(M(x,y)-\mathbb{P}(\Omega_1^{\infty,c})\right).\label{eq:trivial_cond}
\end{align}

Now, by Eq. \ref{eq:trivial_cond}, $(Y_v,v^{-\lambda_1/\lambda_{0,v}}W_v)|\Omega_{1,v}^\infty$ converges in distribution. But then since $\Omega_{0,v}^\infty\subseteq \Omega_{1,v}^\infty$ (up to a probability 0 set) and $\mathbb{P}(\Omega_{1,v}^\infty\Delta \Omega_{0,v}^\infty)\to 0$ as $v\to 0^+$ (where $\Delta$ denotes the symmetric difference) we obtain $(Y_v,v^{-\lambda_1/\lambda_{0,v}}W_v)|\Omega_{0,v}^\infty$ converges with the Laplace Transform specified in the proposition statement.

\end{proof}
\begin{remark}
  The marginal laplace transforms (setting $x=0$ or $y=0$ in Eq. \ref{eq:laplace_transform_limit}) agrees with the Mittag-Leffler distribution derived in a semi-deterministic version of our model in \cite{nicholson2023}.
\end{remark}
Then by the continuous mapping theorem $vY_v W_v^{-\lambda_{0,v}/\lambda_1}|\Omega_{0,v}^\infty$ converges in distribution as $v\to 0^+$. In order to establish Proposition \ref{prop:small_mutation_relative_SFS} we will need to show uniform integrability of  $\{vY_v W_v^{-\lambda_{0,v}/\lambda_1}|\Omega_{0,v}^\infty;0<v<\delta\}$ for $\delta\in (0,1).$

\begin{lemma}
  There is a $\varepsilon>0,\delta<1$ such that $\sup_{\delta>v>0}\mathbb{E}\left[\left(vY_vW^{-\lambda_{0,v}/\lambda_1}\right)^{1+\varepsilon}|\Omega_{0,v}^\infty\right]<\infty$\label{lem:uniform_integrability}
\end{lemma}
\begin{proof}
  On $\Omega_{0,v}^\infty$, there is a first time $T_1^*$ of forming an established type-1 clone. Let the associated size of the scaled, established type-1 clone be $Y_1^{(1),*}$. Then on $\Omega_{0,v}^\infty$:
  \[\left(vY_vW^{-\lambda_{0,v}/\lambda_1}\right)^{1+\varepsilon}\leq \left(vY_ve^{\lambda_{0,v} T_1^*}(Y_1^{(1),*})^{-\lambda_{0,v}/\lambda_1}\right)^{1+\varepsilon}.\]
  Note that an exponentially distributed random variable has finite negative moment of power $-\alpha$ for $0<\alpha<1$
  \[
  \mathbb{E}\left[\left(vY_ve^{\lambda_{0,v} T_1^*}(Y_1^{(1),*})^{-\lambda_{0,v}/\lambda_1}\right)^{1+\varepsilon}|\Omega_{0,v}^\infty\right]=\mathbb{E}\left[\left(Y_1^{(1),*}\right)^{-\frac{\lambda_{0,v}}{\lambda_1}(1+\varepsilon)}\right]\mathbb{E}\left[(vY_v)^{1+\varepsilon}\mathbb{E}\left[e^{\lambda_{0,v}(1+\varepsilon)T_1^*}|\mathcal{F}_\infty,\Omega_{0,v}^\infty\right]|\Omega_{0,v}^\infty\right].
  \]
  Note that $T_1^*$ is the first point of a Cox process with rate $avq_1Z_0(s)$, thus denoting $K>\sup_{\delta>v>0}\mathbb{E}\left[\left(Y_1^{(1),*}\right)^{-\frac{\lambda_{0,v}}{\lambda_1}(1+\varepsilon)}\right]$:
\begin{align*}
  &\leq K\mathbb{E}\left[(vY_v)^{1+\varepsilon}\int_0^\infty e^{\lambda_{0,v}(1+\varepsilon)t}avq_1 Z_0(t)e^{-\int_0^t avq_1 Z_0(s)\,ds}\,dt|\Omega_{0,v}^\infty\right].\\
\end{align*}
Set $u(t)=\int_0^t avq_1 Z_0(s)\,ds$ and on $\Omega_{0,v}^\infty$, $u$ is strictly increasing with range $[0,\infty)$. We can then define $t(u)=u^{-1}(t)$, to rewrite the integral in the previous display as
\begin{align*}
  &=K\mathbb{E}\left[(vY_v)^{1+\varepsilon}\int_0^\infty e^{\lambda_{0,v}(1+\varepsilon)t(u)}e^{-u}\,du|\Omega_{0,v}^\infty\right].
\end{align*}
Let $M(t)=e^{-\lambda_0 t}Z_0(t)$ and notice that on $\Omega_{0,v}^\infty$:
\[
\frac{d}{du}e^{\lambda_{0,v}t(u)}=e^{\lambda_{0,v}t(u)}\frac{\lambda_{0,v}}{avq_1}\frac{1}{Z_0(t(u))}=\frac{\lambda_{0,v}}{avq_1 M(t(u))}.
\]
Hence:
\[\mathbb{E}\left[(vY_v)^{1+\varepsilon}\int_0^\infty e^{\lambda_{0,v}(1+\varepsilon)t(u)}e^{-u}\,du|\Omega_{0,v}^\infty\right]=\mathbb{E}\left[\int_0^\infty \left(vY_v\left(\frac{\lambda_{0,v}}{avq_1}\int_0^u \frac{1}{M(t(z))}\,dz+1\right)\right)^{1+\varepsilon}e^{-u}\,du|\Omega_{0,v}^\infty\right].\]
Then by Minkowsky's inequality:
\[\leq \int_0^\infty \left(\mathbb{E}\left[\left(\frac{\lambda_{0,v}}{aq_1}\int_0^u \frac{Y_v}{M(t(z))}\,dz\right)^{1+\varepsilon}|\Omega_{0,v}^\infty\right]^{1/(1+\varepsilon)}+\mathbb{E}\left[(vY_v)^{1+\varepsilon}|\Omega_{0,v}^\infty\right]^{1/(1+\varepsilon)}\right)^{1+\varepsilon}e^{-u}\,du.\]
Now, it suffices to examine:
\[\mathbb{E}\left[\mathbf{1}_{\Omega_{0,v}^\infty}\left(\frac{\lambda_{0,v}}{aq_1}\int_0^u \frac{Y_v}{M(t(z))}\,dz\right)^{1+\varepsilon}\right]^{1/(1+\varepsilon)}.\]
By Minkowsky's integral inequality:
\[\leq \frac{\lambda_{0,v}}{aq_1}\int_0^u \mathbb{E}\left[\mathbf{1}_{\Omega_{0,v}^\infty}\left(\frac{Y_v}{M(t(z))}\right)^{1+\varepsilon}\right]^{1/(1+\varepsilon)}\,du.\]
Observe that $t(z)$ is a $\mathcal{F}_t$-stopping time (since $\{t(z)\leq r\}=\{u(r)\geq z\}\in\mathcal{F}_r$). On $\Omega_{0,v}^\infty$, $t(z)<\infty$. Therefore, conditioning on $\mathcal{F}_{t(z)}$:

\begin{align*}
    \mathbb{E}\left[\mathbf{1}_{\Omega_{0,v}^\infty}\left(\frac{Y_v}{M(t(z))}\right)^{1+\varepsilon}\right]&\leq \mathbb{E}\left[\sum_{k=1}^\infty \mathbb{E}\left[\left(\frac{Y_v}{M(t(z))}\right)^{1+\varepsilon}|\mathcal{F}_{t(z)}\right]\mathbf{1}_{Z_0(t(z))=k,t(z)<\infty}\right]\\
    &=\mathbb{E}\left[\sum_{k=1}^\infty \mathbb{E}\left[\left(\frac{Y_v}{e^{-\lambda_{0,v}t(z)}k}\right)^{1+\varepsilon}|\mathcal{F}_{t(z)}\right]\mathbf{1}_{Z_0(t(z))=k,t(z)<\infty}\right]
\end{align*}

On $t(z)<\infty$, we may decompose $Y_v=e^{-\lambda_{0,v}t(z)}\sum_{j=1}^k Y_v^{(j)}$ for i.i.d. $Y_v^{(j)}$ (independent of $\mathcal{F}_{t(z)}$). Therefore:
\begin{align*}
    &=\mathbb{E}\left[\sum_{k=1}^\infty \mathbb{E}\left[\left(\sum_{j=1}^k Y_v^{(j)}\frac{1}{k}\right)^{1+\varepsilon}|\mathcal{F}_{t(z)}\right]\mathbf{1}_{Z_0(t(z))=k,t(z)<\infty}\right]\leq \mathbb{E}[(Y_v^{(1)})^{1+\varepsilon}].
\end{align*}

Thus, since positive moments of $Y_v$ are finite and uniformly bounded for $v$ in any compact set, we thus have uniform integrability.
\end{proof}

\SmallMutationRelativeSFS*
\begin{proof}
The first equality is just Theorem \ref{prop:ASRelativeSFS}. For the second equality, first remark that:
\[
\mathbb{E}[\lim_{v\to 0^+}Y_vvW_v^{-\lambda_{0,v}/\lambda_1}|\Omega_{0,v}^\infty]=-\frac{1}{\Gamma(\lambda_{0,0}/\lambda_1)\mathbb{P}(\Omega_{0,0}^\infty)}\int_0^\infty \partial_xM_0(0,\theta)\theta^{\lambda_{0,0}/\lambda_1-1}\,d\theta,
\]
  where $M_0$ is as in Eq. \ref{eq:laplace_transform_limit} and by using formulas for inverse moments of random variables. Here $\lim_v$ is taken in the distribution sense. But from the uniform integrability in Lemma \ref{lem:uniform_integrability}, we can exchange limit and expectation in the left-hand-side of the above equation. Then we can write:
  \begin{align*}
    \lim_{v\to 0^+}\mathbb{E}[Y_vvW_v^{-\lambda_{0,v}/\lambda_1}|\Omega_{0,v}^\infty]&=\frac{1}{\Gamma(\lambda_1/\lambda_{0,0})\mathbb{P}(\Omega_{0,0}^\infty)}\int_0^\infty \frac{1}{\left(1+y^{\lambda_{0,0}/\lambda_1}\frac{ aq_1 ^{1-\lambda_{0,0}/\lambda_1}}{\lambda_1 q_0}\frac{\pi}{\sin(\pi\lambda_{0,0}/\lambda_1)}\right)^2}y^{\lambda_{0,0}/\lambda_1-1}\,dy\\
    &=\frac{1}{\Gamma(\lambda_1/\lambda_{0,0})\mathbb{P}(\Omega_{0,0}^\infty)}\frac{\lambda_1}{\lambda_{0,0}\frac{ aq_1 ^{1-\lambda_{0,0}/\lambda_1}}{\lambda_1 q_0}\frac{\pi}{\sin(\pi\lambda_{0,0}/\lambda_1)}}.
  \end{align*}
The result then follows after simple algebraic manipulations.
\end{proof}

\subsection{Largest clone convergence in the large time limit}\label{app:largest_clone_convergence}
Here we prove Proposition \ref{prop:largest_clone_convergence}.
\LargestClone*
\begin{proof}
For any $K$ and sufficiently large $t$, we have that:
\[
\max_{T_i\leq t}\left( e^{-\lambda_1 T_i}Y_1^{(i)}(t)\right)\geq \max_{1\leq i\leq K}\left( e^{-\lambda_1 T_i}Y_1^{(i)}(t)\right).
\]
Taking the $\liminf$ on both sides and noting that $\lim_{t\to\infty}Y^{(i)}(t)=Y_1^{(i)}$, we thus have:
\[
\liminf_{t\to\infty}\max_{1\leq i\leq M(t)}\left( e^{-\lambda_1 T_i}Y_1^{(i)}(t)\right)\geq \max_{1\leq i\leq K}\left( e^{-\lambda_1 T_i}Y_1^{(i)}\right).
\]
This holds for any $K$, thus we may take the limit as $K\to\infty$:
\[
\liminf_{t\to\infty}\max_{1\leq i\leq M(t)}\left( e^{-\lambda_1 T_i}Y_1^{(i)}(t)\right)\geq \sup_{1\leq i<\infty}\left( e^{-\lambda_1 T_i}Y_1^{(i)}\right).
\]
Now, we use that $\lim_{i\to\infty}e^{-\lambda_1 T_i}Y_*^{(i)}=0$ (Lemma \ref{lem:Y_small}). For any $\varepsilon>0$, there is sufficiently large $I$ such that for all $i\geq I$, $e^{-\lambda_1 T_i}Y_*^{(i)}\leq \varepsilon$. Hence:
\begin{align*}
        \limsup_{t\to\infty}\max_{1\leq i\leq N(t)}\left( e^{-\lambda_1 T_i}Y_1^{(i)}(t)\right)&\leq \limsup_{t\to\infty}\max\{\max_{1\leq i\leq I}\left(e^{-\lambda_1 T_i}Y_1^{(i)}(t)\right),\varepsilon\}\\
        &=\max\{\max_{1\leq i\leq I}\left(e^{-\lambda_1 T_i}Y_1^{(i)}\right),\varepsilon\}\\
        &\leq \max\{\max_{1\leq i<\infty}\left(e^{-\lambda_1 T_i}Y_1^{(i)}\right),\varepsilon\}.
\end{align*}
Now take $\varepsilon\to0^+$, and we have our result. The large-detection size limit directly follows from taking the discrete time sequence $(\tau_n)_n$ and using Lemma \ref{lem:large_size_time} to observe:
    \[
    e^{-\lambda_1 \tau_n}=e^{-\lambda_1(\tau_n-t_n)}e^{-\lambda_1 t_n}\to \frac{W}{n}.
    \]
\end{proof}

\subsection{Largest clone convergence in small mutation limit}\label{app:largest_clone_small}
Here we prove Proposition \ref{prop:largest_clone_small}. Recall that  $\alpha=\lambda_{0,0}/\lambda_1$, $\gamma=\left(\frac{q_1^{1-\alpha}a\Gamma(\alpha)}{\lambda_1 q_0}\right)^{1/\alpha}$,  $V$ denotes the largest element of a Poisson-Dirichlet distribution with parameters $(\alpha,0)$ and $Z$ denotes a random variable with a log-logistic distribution with scale $\gamma$ and shape $\alpha$.
\LargestCloneSmall*
\begin{proof}
  Let $T_i'$ denote the time of forming the $i$th established mutant lineage. Let $Z_1^{(i)'}$ denote the corresponding established type-1 clone, and let $Y_1^{(i)'}=\lim_{t\to\infty}e^{-\lambda_1 t}Z_1^{(i)'}(t)$. Then $\sum \delta_{(T_i',Y_1^{(i)'})}$ conditional on $\mathcal{F}_\infty$ is a Poisson process with intensity $q_1 avZ_0(s)(q_1e^{-q_1 x}\,dx)\,ds$ by the marking theorem. Let $J_v^{(i)}=v^{-\lambda_1/\lambda_{0}}e^{-\lambda_1 T_i'}Y_1^{(i)'}$ and define $\mathcal{N}_v=\sum \delta_{J_v^{(i)}}$. We take this to be a point process over $((0,\infty],r)$ with the vague topology. This metric space is complete and separable using $r(x,y)=\left|\frac{1}{x}-\frac{1}{y}\right|$. In what follows we will first show that $\mathcal{N}_v$ converges weakly to a point process $\mathcal{N}_0$ as $v\to 0$, and then use the continuous mapping theorem to establish the convergence of the maximum and normalized maximum clone sizes. 
  
First note that $v^{-\lambda_1/\lambda_{0}}W=\sum J_v^{(i)}$. Let $f\in C_{0,b}((0,\infty])$ -- namely, $f:((0,\infty],r)\to \mathbb{R}_{\geq 0}$ is a continuous function which is bounded and has bounded support (that is, is $0$ close to $0$). We define further $h:(0,\infty]\to \mathbb{R}$ via $h(j)=\theta j+f(j)$. Call $L_v(\theta,f):[0,\infty)\times C_{0,b}((0,\infty])\to \mathbb{R}_{\geq 0}$ the joint laplace function of $(v^{-\lambda_1/\lambda_0}W,\mathcal{N}_v)|\Omega_{0,v}^\infty$:

\begin{align*}
  L_v(\theta,f)&=\mathbb{E}\left[\exp\left(-\theta v^{-\lambda_1/\lambda_{0}}W-\sum f(J_v^{(i)})\right)\middle| \Omega_{0,v}^\infty\right]=\mathbb{E}\left[\exp\left(-\sum h(J_v^{(i)})\right)\middle| \Omega_{0,v}^\infty\right]\\
  &=\mathbb{E}\left[\exp\left(-avq_1\int_0^\infty \int_0^\infty (1-e^{-h(v^{-\lambda_1/\lambda_0}e^{-\lambda_1 s}x)})q_1 e^{-q_1 x}\,dxZ_0(s)\,ds\right)\middle| \Omega_{0,v}^\infty\right]\\
  &=\mathbb{E}\left[\exp\left(-avq_1\int_0^\infty \int_0^\infty (1-e^{-h(v^{-\lambda_1/\lambda_0}e^{-\lambda_1 s}x)})Z_0(s)\,ds q_1 e^{-q_1 x}\,dx\right)\middle| \Omega_{0,v}^\infty\right].
\end{align*}
Let $r=v^{-\lambda_1/\lambda_0}e^{-\lambda_1 s}x$, hence $s(r)=\lambda_1^{-1}\log r^{-1}v^{-\lambda_1/\lambda_0}x$

\begin{align*}
  &=\mathbb{E}\left[\exp\left(-\lambda_1^{-1}avq_1\int_0^\infty \int_0^{v^{-\lambda_1/\lambda_0}x} (1-e^{-h(r)})e^{\lambda_0 s(r)}e^{-\lambda_0 s(r)}Z_0(s(r))\,\frac{dr}{r} q_1 e^{-q_1 x}\,dx\right)\middle| \Omega_{0,v}^\infty\right]\\
  &=\mathbb{E}\left[\exp\left(-\lambda_1^{-1}aq_1\int_0^\infty \int_0^{v^{-\lambda_1/\lambda_0}x} (1-e^{-h(r)})r^{-\lambda_0/\lambda_1}x^{\lambda_0/\lambda_1}e^{-\lambda_0 s(r)}Z_0(s(r))\,\frac{dr}{r} q_1 e^{-q_1 x}\,dx\right)\middle| \Omega_{0,v}^\infty\right]\\
  &\overset{v\to 0^+}{\to} \mathbb{E}\left[\exp\left(-\lambda_1^{-1}aq_1\int_0^\infty \int_0^{\infty} (1-e^{-h(r)})r^{-\lambda_0/\lambda_1}x^{\lambda_0/\lambda_1}e^{-\lambda_0 s(r)}Y\,\frac{dr}{r} q_1 e^{-q_1 x}\,dx\right)\middle| \Omega_{0,0}^\infty\right]\\
  &=\mathbb{E}\left[\exp\left(-C\cdot Y\int_0^\infty(1-e^{-h(r)})r^{-1-\lambda_0/\lambda_1}\,dr\right)\middle|\Omega_{0,0}^\infty\right],
\end{align*}
where the convergence is justified by similar arguments as in \ref{prop:conv_dist}. Notice that we have integrability near $0$ since $1-e^{-h(r)}\overset{r\to 0^+}{\sim} \theta r$ implying $(1-e^{-h(r)})r^{-\lambda_0/\lambda_1}\overset{r\to 0^+}{\sim}\theta r^{1-\lambda_0/\lambda_1}$. We have integrability near $\infty$ as $r^{-1-\lambda_0/\lambda_1}$ is integrable there. Hence $(v^{-\lambda_1/\lambda_0}W,\mathcal{N}_v)|\Omega_{0,v}^\infty$ converges in distribution. We thus conclude $\mathcal{N}_v$ converges to a point process $\mathcal{N}_0$. From the Laplace functional we conclude that $\mathcal{N}_0$ is doubly-stochastic, where conditional on $Y$, it is a Poisson point process with intensity of a $\lambda_{0,0}/\lambda_1$-subordinator. 

Let $N((0,\infty])$ denote the space of boundedly-finite counting measures with the vague topology. We define the function $G:\mathbb{R}_{>0}\times N((0,\infty])\to \mathbb{R}$ via:
\[
G(s,\mu)=\frac{\sup\text{supp}(\mu)}{s}.
\]
Let $\mu\in N((0,\infty])$ have a largest finite element of the support ($x_{\max}$) which is not a limit point. Let $\mu_n\to\mu$ vaguely. Note that for some $\varepsilon'>0$ with $x_{\max}>\varepsilon'$, $[x_{\max}-\varepsilon',x_{\max}]$ contains only one element in $\text{supp}(\mu)$. Then for any $0<\varepsilon<\varepsilon'$, and for large enough $N$ with $n\geq N$ implying $\mu_n([x_{\max}-\varepsilon,x_{\max}])=1$, we can conclude that $x_{\max}^n\geq x_{\max}-\varepsilon$. Here, $x_{\max}^n=\sup\text{supp }\mu_n$. Also $[x_{\max}+\varepsilon,\infty]$ (which is bounded under $r$) contains no element in $\text{supp}(\mu)$, thus $x_{\max}^n \leq x_{\max}+\varepsilon$. Hence $\mu$ is a continuity point of $\sup\text{supp}(\mu)$. For such $\mu$ and $s>0$, we have $(s,\mu)$ is a continuity point of $G$. By continuous mapping on $\Omega_0^\infty$:
\[
\frac{\sup\left(e^{-\lambda_1 T_i}Y_1^{(i)}\right)}{W}=\frac{\sup\left(v^{-\lambda_1/\lambda_0}e^{-\lambda_1 T_i}Y_1^{(i)}\right)}{v^{-\lambda_1/\lambda_0}W}\overset{v\to 0^+}{\Rightarrow} \sup\left(\frac{\text{supp}(\mathcal{N}_0)}{\sum_{x\in \mathcal{N}_0}x}\right)=V.
\]
By Pitman and Yor \cite{pitman1997}, the above is the largest point of a $\text{PD}(\alpha,0)$ distribution. The CDF of $V$ follows from Proposition 20 of Pitman and Yor.

To establish the limiting distribution of the unnormalized maximum we can again use the fact that the limiting point process is a point of continuity of $\mu\to \sup\text{supp }\mu$. Then by the continuous mapping theorem we get
\begin{align*}
\lim_{v\to 0}\mathbb{P}(\sup_i(e^{-\lambda_1 T_i}Y_1^{(i)})\leq x|\Omega_0^\infty)&=\mathbb{P}(\mathcal{N}_0((x,\infty))=0|\Omega_0^\infty)=\mathbb{E}\left[\exp\left(-CY\int_{x}^\infty y^{-1-\alpha}dy\right)|\Omega_0^\infty\right]\\
&=
\mathbb{E}\left[\exp\left(-\frac{C}{\alpha}Yx^{-\alpha}\right)|\Omega_0^\infty\right].
\end{align*}
On the event $\Omega_0^\infty$, $Y\sim \text{Exp}(q_0)$, and the result then follows.

\end{proof}

\subsection{Poisson-Dirichlet Distribution Factoids}\label{app:poisson_dirichlet}
Here we prove some facts about the largest point in the Poisson-Dirichlet distribution, specifically those stated in Remark \ref{rem:poisson_dirichlet}. Here we let $\alpha=\lambda_{0,0}/\lambda_1$ and mutations are only attached to births. $V$ denotes the largest point 
\begin{lemma}
    As long as $\frac{\sin(\pi \alpha)}{\pi\alpha}\geq \frac{1}{2}$, $(0<\alpha\lesssim 0.603$, that is $b$ is sufficiently large), the median of $V$ is:
\[\frac{1}{1+\left(\frac{\pi\alpha}{2\sin(\pi\alpha)}\right)^{1/\alpha}}.\]
\end{lemma}
\begin{proof}
    Let $p_V$ denote the probability density function of $V$. If $1>u>1/2$, then by Proposition 20 of Pitman and Yor:
    \[p_V(u)=\frac{1}{\Gamma(\alpha)\Gamma(1-\alpha)}\frac{(1-u)^{\alpha-1}}{u^{\alpha+1}}.\]
    Remark that $\int_{1/2}^1 p_V(u)\,du\geq \frac{1}{2}$ iff $\frac{\sin\pi\alpha}{\pi\alpha}\geq \frac{1}{2}$ and under this condition the median ($m$) satisfies $m\geq 1/2$ and can be solved via $\int_m^1 p_V(u)\,du=1/2$.
\end{proof}
\begin{lemma}
    $V$ is small very infrequently:
\[\mathbb{P}(V\leq x)\lsimas{x}{0^+}e^{-cx^{-1}}\]
Where $c$ can be explicitly determined.
\end{lemma}
\begin{proof}
    Remark that the characteristic function of $1/V$ is known (Corollary 12 of Pitman and Yor) to be:
    \[G(t)=\mathbb{E}\left[e^{it/V}\right]=\frac{e^{it}}{1+\alpha\int_0^1 \left(1-e^{itx}\right)x^{-\alpha-1}\,dx}\]
    The numerator and denominator are entire (the denominator is entire since for any disk $D\subseteq \mathbb{C}$ with its closure compact $\int_{\varepsilon}^1 \left(1-e^{itx}\right)x^{-\alpha-1}\,dx$ converges uniformly over $D$ as $\varepsilon\to 0^+$). Letting $t=-iy$, remark that $\int_0^1 (1-e^{yx})x^{-\alpha-1}\,dx$ is strictly decreasing, goes to $-\infty$ as $y\to\infty$, and is $1$ when $y=0$. In particular, there is a unique $y^*>0$ where the denominator of $G(-iy)$ is $0$. By Theorem 2.4b of \cite{vernon_laplace_1946}:
    \[-y^*=-\limsup_{x\to\infty}\frac{\mathbb{P}(1/V\geq x)}{x}\]
    The lemma thus follows with $c$ being $-y^*+\varepsilon$ for any $\varepsilon>0$.
\end{proof}

\subsection{Supplementary Figures}
\begin{figure}[H]
	\centering
	\includegraphics[width=\textwidth]{assets/sfs_heatmap.pdf}
	\caption{Mean site frequency spectrum heatmap at varying sizes $j$ and times $t$}\label{fig:SFS_heatmap}
\end{figure}

\end{document}
